% Generated by tools/build.py -- do not edit; edit data/operating-systems/*.txt
% No preamble here: this file is \input by main.tex.
\chapter{Unit 5: System Software and Operating System}

\section{System Software (language levels, assemblers, macros, loaders and linkers, compilers and interpreters, debuggers)}

\begin{question}{H}
In a two-pass assembler, the symbol table is built during
\optsshort{Pass 2}{the loading phase}{Pass 1}{both passes equally}
\end{question}

\begin{question}{H}
The main reason for using a two-pass assembler is to handle
\optsshort{macro expansion}{dynamic linking}{relocation}{forward references}
\end{question}

\begin{question}{H}
Match List I with List II.\\
List I (Assembler table): A. OPTAB  B. SYMTAB  C. LITTAB  D. LOCCTR\\
List II (Contents/purpose): I. Literal constants and their addresses  II. Mnemonic operation codes and their machine equivalents  III. Address of the next statement being assembled  IV. Labels and their assigned addresses
\optsshort{A-II, B-IV, C-III, D-I}{A-II, B-I, C-IV, D-III}{A-IV, B-II, C-I, D-III}{A-II, B-IV, C-I, D-III}
\end{question}

\begin{question}{H}
Which of the following tasks is performed in Pass 2 of a two-pass assembler?
\optslong{Processing START to set the initial location counter}{Translating opcodes and symbolic operands and writing the object program}{Saving label values in the symbol table}{Assigning addresses to all statements}
\end{question}

\begin{question}{H}
Which of the following are assembler directives (pseudo-operations) that do not generate machine instructions?\\
A. START  B. EQU  C. LDA  D. ORG  E. LTORG\\
Choose the correct answer from the options given below:
\optsshort{A, B and C only}{C and E only}{B, C and D only}{A, B, D and E only}
\end{question}

\begin{question}{H}
In a SIC assembly program the location counter is 1003 (hexadecimal) when the statement BUFFER RESB 4096 is processed. The address assigned to the next statement is
\optsshort{5099}{2006}{1003}{2003}
\end{question}

\begin{question}{H}
In SIC assembly, the statement TABLE RESW 100 reserves
\optsshort{200 bytes}{400 bytes}{300 bytes}{100 bytes}
\end{question}

\begin{question}{M}
In SIC assembly, how many bytes are generated by BYTE C'EOF' and BYTE X'F1' respectively?
\optsshort{6 and 2}{1 and 1}{3 and 1}{3 and 2}
\end{question}

\begin{question}{M}
Match List I with List II.\\
List I (SIC/XE object program record): A. H  B. T  C. E  D. M\\
List II (Contents): I. Address of the first executable instruction  II. Program name, starting address and length  III. Address and length of a field to be changed at load time  IV. Object code and its starting address
\optsshort{A-II, B-IV, C-I, D-III}{A-II, B-IV, C-III, D-I}{A-IV, B-II, C-I, D-III}{A-II, B-I, C-IV, D-III}
\end{question}

\begin{question}{M}
A one-pass assembler handles forward references by
\optslong{asking the loader to resolve every symbol}{disallowing labels}{leaving the address blank and filling it in (backpatching) when the symbol is defined}{using a second pass over the source}
\end{question}

\begin{question}{M}
The EQU directive is used to
\optslong{mark the end of the program}{define a symbol and assign it a value}{place the literal pool}{reserve memory for a variable}
\end{question}

\begin{question}{M}
The LTORG directive instructs the assembler to
\optslong{reset the location counter}{start a new control section}{place the literals collected so far into a literal pool at the current location}{end the macro definition}
\end{question}

\begin{question}{H}
Given below are two statements: one is labelled as Assertion A and the other is labelled as Reason R.\\
Assertion A: A program that uses macros generally executes faster than an equivalent program that uses subroutines.\\
Reason R: A macro call is expanded in-line at assembly time, so there is no call and return overhead at run time.
\optslong{Both A and R are true and R is the correct explanation of A}{Both A and R are true but R is NOT the correct explanation of A}{A is true but R is false}{A is false but R is true}
\end{question}

\begin{question}{H}
Match List I with List II.\\
List I (Macro processor data structure): A. MNT  B. MDT  C. ALA\\
List II (Purpose): I. Stores the body of macro definitions  II. Maps formal parameters to actual arguments  III. Stores macro names with a pointer to their definitions
\optsshort{A-II, B-I, C-III}{A-III, B-II, C-I}{A-III, B-I, C-II}{A-I, B-III, C-II}
\end{question}

\begin{question}{M}
Given below are two statements.\\
Statement I: A macro is expanded every time it is called, which increases the size of the program.\\
Statement II: A subroutine call is expanded in-line wherever it appears in the program.
\optslong{Both Statement I and Statement II are true}{Both Statement I and Statement II are false}{Statement I is true but Statement II is false}{Statement I is false but Statement II is true}
\end{question}

\begin{question}{M}
Conditional macro expansion in IBM 360 assembler is supported by the pseudo-operations
\optsshort{START and END}{AIF and AGO}{USING and DROP}{EQU and ORG}
\end{question}

\begin{question}{H}
The four basic functions of a loader are
\optslong{allocation, linking, relocation and loading}{scanning, parsing, linking and loading}{compilation, assembly, linking and loading}{allocation, compaction, relocation and paging}
\end{question}

\begin{question}{H}
In an absolute loader, the functions of allocation, linking, relocation and loading are performed respectively by
\optslong{the loader, the loader, the loader and the loader}{the assembler, the programmer, the loader and the loader}{the programmer, the assembler, the loader and the loader}{the programmer, the programmer, the assembler and the loader}
\end{question}

\begin{question}{H}
Match List I with List II.\\
List I (Loader): A. Compile-and-go loader  B. Absolute loader  C. Direct linking loader  D. Bootstrap loader\\
List II (Characteristic): I. A relocatable loader that uses ESD, TXT and RLD cards  II. Loads the operating system when the computer is switched on  III. The assembler places object code directly in memory and transfers control  IV. Loads code at the addresses fixed in the object program
\optsshort{A-III, B-I, C-IV, D-II}{A-II, B-IV, C-I, D-III}{A-III, B-IV, C-I, D-II}{A-IV, B-III, C-I, D-II}
\end{question}

\begin{question}{M}
A disadvantage of the compile-and-go (assemble-and-go) loading scheme is that
\optslong{the assembler occupies memory during execution and the program must be retranslated every time it is run}{it supports only one segment}{the program cannot be relocated}{it requires a linkage editor}
\end{question}

\begin{question}{H}
Given below are two statements.\\
Statement I: In dynamic linking, all library routines are copied into the executable file at link time.\\
Statement II: In dynamic loading, a routine is not loaded into memory until it is called.
\optslong{Both Statement I and Statement II are true}{Both Statement I and Statement II are false}{Statement I is true but Statement II is false}{Statement I is false but Statement II is true}
\end{question}

\begin{question}{M}
A linkage editor differs from a linking loader in that the linkage editor
\optslong{cannot resolve external references}{performs linking at execution time}{loads the program directly into memory}{produces a linked load module that is saved for later execution}
\end{question}

\begin{question}{M}
In the direct linking loader, the external symbol dictionary (ESD) card contains
\optslong{information about symbols defined in one segment and referenced in another}{the relocation bits}{the address of the first executable instruction}{the object code}
\end{question}

\begin{question}{M}
In SIC/XE, the assembler directives used to define and refer to external symbols across control sections are
\optsshort{EXTDEF and EXTREF}{USING and DROP}{ENTRY and EXTRN}{START and END}
\end{question}

\begin{question}{H}
A program assembled to start at address 0 has an instruction that refers to address 500. If the loader loads the program starting at address 2000, the relocated address is
\optsshort{2500}{1500}{500}{2000}
\end{question}

\begin{question}{M}
In a SIC object program, relocation of text records is indicated using
\optsshort{the symbol table}{a relocation bit mask}{the literal table}{modification records}
\end{question}

\begin{question}{H}
Given below are two statements: one is labelled as Assertion A and the other is labelled as Reason R.\\
Assertion A: An interpreted program generally runs slower than the same program compiled to machine code.\\
Reason R: An interpreter analyses each statement every time it is executed.
\optslong{Both A and R are true and R is the correct explanation of A}{Both A and R are true but R is NOT the correct explanation of A}{A is true but R is false}{A is false but R is true}
\end{question}

\begin{question}{H}
Which of the following statements about compilers and interpreters is correct?
\optslong{An interpreter cannot report errors}{A compiler translates the whole program before execution, while an interpreter translates and executes one statement at a time}{An interpreter produces an object file that can be executed later}{Compiled programs run slower than interpreted programs}
\end{question}

\begin{question}{M}
A compiler that runs on one machine and generates code for a different machine is called a
\optsshort{just-in-time compiler}{cross compiler}{bootstrap compiler}{decompiler}
\end{question}

\begin{question}{M}
Which of the following translates bytecode into native machine code at run time, as in the Java Virtual Machine?
\optsshort{Just-in-time (JIT) compiler}{Assembler}{Linkage editor}{Preprocessor}
\end{question}

\begin{question}{M}
The process of writing a compiler in its own source language and using a simple version of it to compile more complete versions is called
\optsshort{interpretation}{bootstrapping}{cross compilation}{optimisation}
\end{question}

\begin{question}{M}
Match List I with List II.\\
List I (Language level): A. Machine language  B. Assembly language  C. High-level language\\
List II (Characteristic): I. Uses mnemonics with a one-to-one correspondence to machine instructions  II. Machine-independent and must be compiled or interpreted  III. Binary instructions executed directly by the CPU
\optsshort{A-II, B-I, C-III}{A-I, B-III, C-II}{A-III, B-II, C-I}{A-III, B-I, C-II}
\end{question}

\begin{question}{M}
Arrange the following in the order in which they normally process a C program.\\
A. Loader  B. Compiler  C. Preprocessor  D. Linker  E. Assembler
\optsshort{B, C, E, D, A}{C, B, E, D, A}{C, E, B, D, A}{C, B, D, E, A}
\end{question}

\begin{question}{M}
Which of the following is a function of the C preprocessor?
\optslong{Allocating registers}{Macro substitution and file inclusion}{Linking library functions}{Generating object code}
\end{question}

\begin{question}{H}
A debugger temporarily stops the execution of a program at a specified instruction using a
\optsshort{core dump}{watchpoint}{breakpoint}{trace}
\end{question}

\begin{question}{M}
Executing a program one instruction or one statement at a time under a debugger is called
\optsshort{profiling}{single stepping}{dumping}{tracing}
\end{question}

\begin{question}{M}
For a symbolic debugger to show variable names and source lines, the compiler must include
\optslong{a relocation table}{a literal pool}{symbol table and line number information in the object file}{the modification records}
\end{question}

\begin{question}{M}
A file containing the memory image of a process at the time it crashed, used for post-mortem debugging, is called a
\optsshort{map file}{load module}{core dump}{symbol table}
\end{question}

\begin{question}{M}
Given below are two statements.\\
Statement I: Literals and symbol-defining statements are machine-dependent assembler features.\\
Statement II: Instruction formats and addressing modes are machine-independent assembler features.
\optslong{Both Statement I and Statement II are true}{Both Statement I and Statement II are false}{Statement I is true but Statement II is false}{Statement I is false but Statement II is true}
\end{question}

\begin{question}{M}
Which of the following is NOT a system software?
\optsshort{Loader}{Assembler}{Operating system}{Spreadsheet program}
\end{question}

\begin{question}{M}
In program relocation, a program whose code can be loaded and executed at any address without modification is called
\optsshort{reentrant code}{self-relocating code}{position-independent code}{absolute code}
\end{question}

\begin{question}{M}
A program that can be shared by several users at the same time because it does not modify itself is called
\optsshort{relocatable code}{reentrant code}{self-modifying code}{overlay}
\end{question}

\begin{question}{M}
In which technique is only the part of a program needed at a given time kept in memory, with other parts loaded into the same region when needed, under program control?
\optsshort{Overlays}{Paging}{Swapping}{Dynamic linking}
\end{question}

\begin{question}{M}
Arrange the following loader schemes in increasing order of the work done by the loader.\\
A. Absolute loader  B. Compile-and-go loader  C. Direct linking loader
\optsshort{A, C, B}{B, C, A}{B, A, C}{A, B, C}
\end{question}

\begin{question}{L}
Which of the following is an assembler directive used in IBM 360 assembly to inform the assembler which base register to use?
\optsshort{USING}{LTORG}{DC}{EQU}
\end{question}

\begin{question}{L}
In IBM 360 assembly, DC and DS are used respectively to
\optslong{declare a macro and end it}{define storage and define a constant}{define a constant and reserve storage}{load a register and store it}
\end{question}

\begin{question}{M}
Given below are two statements: one is labelled as Assertion A and the other is labelled as Reason R.\\
Assertion A: Shared libraries reduce both disk space and main memory usage.\\
Reason R: With dynamic linking only one copy of the library code needs to be kept in memory and shared by all processes that use it.
\optslong{Both A and R are true and R is the correct explanation of A}{Both A and R are true but R is NOT the correct explanation of A}{A is true but R is false}{A is false but R is true}
\end{question}

\begin{answers}
\ans{1}{3}{Pass 1 assigns addresses using the location counter and records every label in SYMTAB; Pass 2 uses it to generate code.}
\ans{2}{4}{A symbol may be used before it is defined; Pass 1 collects all definitions so that Pass 2 can resolve every reference.}
\ans{3}{4}{OPTAB holds opcodes, SYMTAB labels, LITTAB literals, and LOCCTR is the location counter.}
\ans{4}{2}{Pass 1 defines symbols; Pass 2 assembles instructions, generates data values and writes the object program.}
\ans{5}{4}{LDA is a machine instruction; the others direct the assembler.}
\ans{6}{4}{4096 decimal = 1000 hexadecimal, so the next address is 1003 + 1000 = 2003 (hex).}
\ans{7}{3}{A SIC word is 3 bytes, so RESW 100 reserves 300 bytes.}
\ans{8}{3}{C'EOF' gives one byte per character; X'F1' gives two hexadecimal digits, i.e. one byte.}
\ans{9}{1}{Header, Text, End and Modification records.}
\ans{10}{3}{It keeps a list of incomplete instructions for each undefined symbol and patches them once the symbol is defined.}
\ans{11}{2}{EQU enters a symbol with the given value into the symbol table without generating code.}
\ans{12}{3}{Without LTORG, literals are placed at the end of the program.}
\ans{13}{1}{In-line expansion removes the call overhead, at the cost of larger code.}
\ans{14}{3}{Macro name table, macro definition table and argument list array.}
\ans{15}{3}{Subroutine code is stored once and control is transferred to it at run time; only macros are expanded in-line.}
\ans{16}{2}{AIF is a conditional branch and AGO an unconditional branch during macro expansion.}
\ans{17}{1}{Allocation of space, linking of references, relocation of addresses and physically loading the code.}
\ans{18}{4}{The absolute loader only loads; the programmer fixes addresses and links, and the assembler produces absolute code.}
\ans{19}{3}{Standard loader schemes from Donovan's System Programming.}
\ans{20}{1}{No object file is saved, so the source is retranslated on every run, and memory is wasted on the assembler.}
\ans{21}{4}{Copying library code into the executable is static linking; dynamic linking postpones linking until run time.}
\ans{22}{4}{A linking loader links and loads every time; a linkage editor writes a load module once.}
\ans{23}{1}{ESD lists external symbols (SD, LD and ER entries); RLD records the locations that need relocation or linking.}
\ans{24}{1}{EXTDEF names symbols defined here; EXTREF names symbols used here but defined elsewhere. (ENTRY and EXTRN are the IBM 360 equivalents.)}
\ans{25}{1}{Relocation adds the load address to every address-sensitive field: 2000 + 500 = 2500.}
\ans{26}{2}{SIC uses a bit mask (one bit per word); SIC/XE uses modification records.}
\ans{27}{1}{Repeated analysis of the source at run time is the main overhead of interpretation.}
\ans{28}{2}{A compiler generates target code for the whole program; an interpreter executes as it translates.}
\ans{29}{2}{Cross compilers are used, for example, to build software for embedded targets.}
\ans{30}{1}{A JIT compiler compiles frequently executed bytecode to native code during execution.}
\ans{31}{2}{Bootstrapping is usually explained with T-diagrams.}
\ans{32}{4}{Assembly is a symbolic form of machine language; high-level languages are machine-independent.}
\ans{33}{2}{Preprocessor, compiler, assembler, linker, loader.}
\ans{34}{2}{\#define and \#include are handled before compilation proper.}
\ans{35}{3}{A breakpoint stops at a code location; a watchpoint stops when a variable's value changes.}
\ans{36}{2}{Single stepping lets the programmer examine the state after each step.}
\ans{37}{3}{For example, gcc -g emits debugging information (DWARF).}
\ans{38}{3}{A core dump records memory and register contents for later analysis.}
\ans{39}{2}{It is the other way round: instruction formats and addressing modes are machine-dependent; literals, EQU and expressions are machine-independent.}
\ans{40}{4}{A spreadsheet is application software.}
\ans{41}{3}{Position-independent code uses relative addressing and is widely used for shared libraries.}
\ans{42}{2}{Reentrant (pure) code keeps per-user data separate from the shared instructions.}
\ans{43}{1}{Overlays were used before virtual memory to run programs larger than memory.}
\ans{44}{3}{Compile-and-go has no separate loader; an absolute loader only loads; a direct linking loader also relocates and links.}
\ans{45}{1}{USING declares a base register and its contents; DROP releases it.}
\ans{46}{3}{DC defines a constant; DS reserves storage without initialising it.}
\ans{47}{1}{Static linking duplicates library code in every executable and every process image.}
\end{answers}

\section{Operating System Basics (structure, services, system calls, dual mode, system boot)}

\begin{question}{H}
Match List I with List II.\\
List I (Kernel structure): A. Monolithic  B. Microkernel  C. Layered  D. Hybrid\\
List II (Example): I. Mach  II. Windows NT  III. Traditional UNIX  IV. THE operating system
\optsshort{A-III, B-IV, C-I, D-II}{A-III, B-I, C-II, D-IV}{A-III, B-I, C-IV, D-II}{A-I, B-III, C-IV, D-II}
\end{question}

\begin{question}{H}
Given below are two statements: one is labelled as Assertion A and the other is labelled as Reason R.\\
Assertion A: A microkernel operating system is generally more reliable and secure than a monolithic one.\\
Reason R: Most services run as user-level processes, so a failure in a service does not crash the kernel.
\optslong{Both A and R are true and R is the correct explanation of A}{Both A and R are true but R is NOT the correct explanation of A}{A is true but R is false}{A is false but R is true}
\end{question}

\begin{question}{H}
The main disadvantage of the microkernel approach is
\optslong{difficulty of porting}{performance overhead of increased message passing between user-space services}{lack of modularity}{poor security}
\end{question}

\begin{question}{H}
Which of the following instructions should be privileged (allowed only in kernel mode)?\\
A. Set the value of the timer  B. Read the clock  C. Turn off interrupts  D. Access an I/O device  E. Issue a trap instruction\\
Choose the correct answer from the options given below:
\optsshort{A, B and C only}{A, C, D and E only}{A, C and D only}{B, D and E only}
\end{question}

\begin{question}{H}
In the dual-mode operation described by Silberschatz, the mode bit is
\optslong{1 for kernel mode and 0 for user mode}{0 for kernel mode and 1 for user mode}{0 for both modes}{not stored in hardware}
\end{question}

\begin{question}{H}
A timer is used in an operating system mainly to
\optslong{keep the date and time}{synchronise processes}{measure disk speed}{prevent a user program from getting stuck in an infinite loop and never returning control to the OS}
\end{question}

\begin{question}{H}
Match List I with List II.\\
List I (System call category): A. Process control  B. File management  C. Information maintenance  D. Protection\\
List II (Example UNIX call): I. getpid()  II. chmod()  III. fork()  IV. open()
\optsshort{A-III, B-I, C-IV, D-II}{A-III, B-IV, C-II, D-I}{A-III, B-IV, C-I, D-II}{A-IV, B-III, C-I, D-II}
\end{question}

\begin{question}{H}
Match List I with List II.\\
List I (Windows API): A. CreateProcess()  B. WaitForSingleObject()  C. CreatePipe()  D. ReadFile()\\
List II (UNIX equivalent): I. pipe()  II. wait()  III. fork()  IV. read()
\optsshort{A-III, B-I, C-II, D-IV}{A-III, B-II, C-I, D-IV}{A-III, B-II, C-IV, D-I}{A-II, B-III, C-I, D-IV}
\end{question}

\begin{question}{M}
Which of the following are general methods of passing parameters to a system call?\\
A. In registers  B. In a block or table in memory whose address is passed in a register  C. Pushed onto the stack  D. Through a file on disk\\
Choose the correct answer from the options given below:
\optsshort{A, B, C and D}{B, C and D only}{A and B only}{A, B and C only}
\end{question}

\begin{question}{M}
A trap (or exception) is
\optslong{a type of system program}{a scheduling algorithm}{a hardware interrupt from an I/O device}{a software-generated interrupt caused by an error or by a user request for an OS service}
\end{question}

\begin{question}{M}
The printf() function in a C program ultimately invokes which system call to display output?
\optsshort{read()}{write()}{exec()}{ioctl()}
\end{question}

\begin{question}{M}
Given below are two statements.\\
Statement I: An application programming interface such as POSIX or the Win32 API hides most details of the system call interface from programmers.\\
Statement II: A program written using the POSIX API can, in principle, be compiled on any system that supports POSIX.
\optslong{Both Statement I and Statement II are true}{Both Statement I and Statement II are false}{Statement I is true but Statement II is false}{Statement I is false but Statement II is true}
\end{question}

\begin{question}{H}
The initial program that runs when a computer is powered on and loads the operating system kernel is the
\optslong{interrupt handler}{scheduler}{shell}{bootstrap program (boot loader)}
\end{question}

\begin{question}{M}
The Master Boot Record (MBR) on a BIOS-based disk occupies
\optslong{the first sector (512 bytes) of the disk}{the last sector of the disk}{the first track}{a separate EEPROM chip}
\end{question}

\begin{question}{M}
On a UNIX/Linux system, the first user-level process started by the kernel, which has process ID 1, is
\optsshort{swapper}{cron}{bash}{init (or systemd)}
\end{question}

\begin{question}{H}
Match List I with List II.\\
List I (System type): A. Batch system  B. Time-sharing system  C. Hard real-time system  D. Distributed system\\
List II (Characteristic): I. Collection of independent computers that appears to users as a single system  II. Jobs with similar needs are grouped and run without user interaction  III. CPU is switched among users frequently to give fast response  IV. Tasks must be completed within guaranteed deadlines
\optsshort{A-II, B-III, C-I, D-IV}{A-II, B-III, C-IV, D-I}{A-III, B-II, C-IV, D-I}{A-II, B-IV, C-III, D-I}
\end{question}

\begin{question}{H}
Given below are two statements.\\
Statement I: Multiprogramming increases CPU utilisation by keeping several jobs in memory so that the CPU always has one to execute.\\
Statement II: Time-sharing systems give longer response times than batch systems.
\optslong{Both Statement I and Statement II are true}{Both Statement I and Statement II are false}{Statement I is true but Statement II is false}{Statement I is false but Statement II is true}
\end{question}

\begin{question}{M}
Spooling stands for
\optslong{Sequential Process Ordering On-Line}{System Program Operation On-Line}{Simultaneous Peripheral Operations On-Line}{Shared Peripheral Output On Line}
\end{question}

\begin{question}{M}
The separation of policy from mechanism means that
\optslong{policies are implemented in hardware}{mechanisms change frequently while policies never change}{mechanisms determine how to do something while policies decide what will be done}{the OS decides what to do and how to do it in the same module}
\end{question}

\begin{question}{M}
Arrange the following storage devices in decreasing order of access speed.\\
A. Main memory  B. Registers  C. Hard disk  D. Cache  E. Solid-state disk
\optsshort{B, D, A, C, E}{B, A, D, E, C}{B, D, A, E, C}{D, B, A, C, E}
\end{question}

\begin{question}{M}
Which of the following is a volatile storage?
\optsshort{Main memory (DRAM)}{Flash memory}{Magnetic disk}{Optical disk}
\end{question}

\begin{question}{M}
In symmetric multiprocessing (SMP)
\optslong{one processor controls the others and assigns them tasks}{only one processor can run the OS while others run user programs}{each processor performs all tasks, including operating-system functions, and they share memory}{processors do not share memory}
\end{question}

\begin{question}{M}
In a clustered system using asymmetric clustering
\optslong{one machine is in hot-standby mode, monitoring the active server and taking over if it fails}{all machines run applications and monitor each other}{there is no shared storage}{each node runs a different operating system}
\end{question}

\begin{question}{M}
Which of the following is NOT an operating-system service?
\optslong{Compiling application programs}{I/O operations}{Program execution}{Error detection}
\end{question}

\begin{question}{M}
The command interpreter in UNIX-like systems is called the
\optsshort{loader}{daemon}{kernel}{shell}
\end{question}

\begin{question}{M}
Given below are two statements: one is labelled as Assertion A and the other is labelled as Reason R.\\
Assertion A: Loadable kernel modules allow the kernel to be extended without recompiling it or rebooting.\\
Reason R: Linux is distributed under the GNU General Public License.
\optslong{Both A and R are true and R is the correct explanation of A}{Both A and R are true but R is NOT the correct explanation of A}{A is true but R is false}{A is false but R is true}
\end{question}

\begin{question}{M}
Which of the following is used by an I/O device to signal the CPU that it needs attention?
\optsshort{Polling}{System call}{Trap}{Interrupt}
\end{question}

\begin{question}{M}
The interrupt vector is
\optslong{a queue of pending I/O requests}{the stack of the interrupted process}{a table of addresses of interrupt service routines}{a register that stores the mode bit}
\end{question}

\begin{question}{M}
Which of the following is a non-maskable interrupt typically used for?
\optsshort{Unrecoverable memory errors}{Keyboard input}{Disk I/O completion}{Network packet arrival}
\end{question}

\begin{question}{M}
Match List I with List II.\\
List I (Term): A. Multiprocessing  B. Multiprogramming  C. Multithreading  D. Multitasking\\
List II (Meaning): I. Several programs in memory, with the CPU switched when a job waits  II. More than one CPU executing simultaneously  III. Several threads of control within one process  IV. Rapid switching among jobs to give interactive response
\optsshort{A-I, B-II, C-III, D-IV}{A-II, B-I, C-IV, D-III}{A-II, B-I, C-III, D-IV}{A-II, B-III, C-I, D-IV}
\end{question}

\begin{question}{M}
Which of the following is an example of a real-time operating system?
\optsshort{VxWorks}{Windows 95}{MS-DOS}{Ubuntu Desktop}
\end{question}

\begin{question}{M}
Given below are two statements.\\
Statement I: In a soft real-time system, missing a deadline is treated as a total system failure.\\
Statement II: Soft real-time systems give critical tasks priority over others but do not guarantee deadlines.
\optslong{Both Statement I and Statement II are true}{Both Statement I and Statement II are false}{Statement I is true but Statement II is false}{Statement I is false but Statement II is true}
\end{question}

\begin{question}{M}
Which of the following tools traces the system calls made by a running Linux process?
\optsshort{grep}{gdb}{strace}{ldd}
\end{question}

\begin{question}{M}
Which of the following is responsible for memory protection, i.e. preventing a process from accessing memory outside its address space, in the simplest scheme?
\optsshort{Program counter}{Interrupt vector}{Stack pointer}{Base and limit registers}
\end{question}

\begin{question}{L}
An operating system that allows a single user to run only one program at a time, such as MS-DOS, is called a
\optslong{time-sharing OS}{single-user single-tasking OS}{distributed OS}{multi-user OS}
\end{question}

\begin{question}{M}
The kernel of macOS and iOS (Darwin/XNU) combines
\optslong{the Windows NT kernel and Mach}{the Linux kernel and BSD}{the Mach microkernel and BSD UNIX components}{MINIX and Linux}
\end{question}

\begin{question}{M}
Which of the following events causes a switch from user mode to kernel mode?\\
A. A system call  B. A hardware interrupt  C. A page fault  D. A function call within a user program\\
Choose the correct answer from the options given below:
\optsshort{A and D only}{A, B, C and D}{B, C and D only}{A, B and C only}
\end{question}

\begin{question}{M}
The process of starting a computer by loading the kernel is called
\optsshort{booting}{paging}{spooling}{swapping}
\end{question}

\begin{question}{L}
The portion of the operating system that is always resident in main memory is called the
\optsshort{loader}{shell}{kernel}{compiler}
\end{question}

\begin{question}{M}
Given below are two statements.\\
Statement I: An emulator allows applications compiled for one hardware architecture to run on a different architecture.\\
Statement II: Emulation usually runs faster than running the code natively.
\optslong{Both Statement I and Statement II are true}{Both Statement I and Statement II are false}{Statement I is true but Statement II is false}{Statement I is false but Statement II is true}
\end{question}

\begin{question}{M}
Which of the following is NOT a component of the UNIX system call interface category "device management"?
\optsshort{fork()}{read() on a device}{write() on a device}{ioctl()}
\end{question}

\begin{answers}
\ans{1}{3}{UNIX is monolithic, Mach a microkernel, Dijkstra's THE system layered, Windows NT a hybrid.}
\ans{2}{1}{Moving services out of the kernel limits the damage of faults, at the cost of message-passing overhead.}
\ans{3}{2}{Services communicate through the kernel by message passing, adding system-call and context-switch overhead.}
\ans{4}{3}{Reading the clock is harmless, and a trap is how user programs request kernel services.}
\ans{5}{2}{The hardware switches to kernel mode (0) on a trap or interrupt and back to user mode (1) before passing control to a user program.}
\ans{6}{4}{The timer interrupts after a set period, so the OS regains control.}
\ans{7}{3}{fork is process control, open file management, getpid information maintenance, chmod protection.}
\ans{8}{2}{Standard Windows/UNIX system call correspondence.}
\ans{9}{4}{Registers, a memory block, or the stack are used; files are not.}
\ans{10}{4}{Examples: division by zero, an invalid memory access, or a system call.}
\ans{11}{2}{printf is a library function that buffers output and calls write().}
\ans{12}{1}{APIs give portability and convenience over raw system calls.}
\ans{13}{4}{It is stored in firmware (ROM/EEPROM, BIOS or UEFI) and may load a second-stage loader such as GRUB.}
\ans{14}{1}{The MBR holds the first-stage boot code and the partition table; GPT is the UEFI replacement.}
\ans{15}{4}{init/systemd starts all other services and adopts orphan processes.}
\ans{16}{2}{Standard classification of operating systems.}
\ans{17}{3}{Time-sharing switches the CPU among users frequently precisely to give short response times.}
\ans{18}{3}{Spooling uses the disk as a large buffer, for example for printer output.}
\ans{19}{3}{Separating them allows a policy to be changed without changing the underlying mechanism.}
\ans{20}{3}{Registers, cache, main memory, SSD, magnetic disk.}
\ans{21}{1}{DRAM loses its contents when power is removed.}
\ans{22}{3}{Asymmetric multiprocessing uses a boss-worker relationship.}
\ans{23}{1}{In symmetric clustering, all hosts run applications and monitor each other.}
\ans{24}{1}{Compilation is done by a compiler, which is a system program, not an OS service.}
\ans{25}{4}{Examples are sh, bash, csh and ksh.}
\ans{26}{2}{Both are true, but the licence has nothing to do with why modules can be loaded at run time.}
\ans{27}{4}{Device controllers raise interrupts; the CPU uses the interrupt vector to find the handler.}
\ans{28}{3}{It is indexed by the interrupt number to locate the handler quickly.}
\ans{29}{1}{Non-maskable interrupts cannot be disabled and are reserved for critical events.}
\ans{30}{3}{Standard definitions of the four terms.}
\ans{31}{1}{VxWorks, QNX and FreeRTOS are real-time operating systems.}
\ans{32}{4}{Missing a deadline is a failure only in a hard real-time system.}
\ans{33}{3}{strace lists each system call with its arguments and return value.}
\ans{34}{4}{Every address generated in user mode is checked against the base and limit registers.}
\ans{35}{2}{MS-DOS ran one program at a time for a single user.}
\ans{36}{3}{XNU is a hybrid kernel built on Mach with BSD layers.}
\ans{37}{4}{An ordinary function call stays in user mode.}
\ans{38}{1}{The bootstrap program locates, loads and starts the kernel.}
\ans{39}{3}{The kernel runs at all times; other programs are system or application programs.}
\ans{40}{3}{Every instruction must be translated, so emulation is much slower than native execution.}
\ans{41}{1}{fork is process control; ioctl, read and write can operate on devices.}
\end{answers}

\section{Processes, Threads and Interprocess Communication}

\begin{question}{H}
How many times is "hello" printed by the following C code?\\
fork(); fork(); printf("hello");
\optsshort{4}{8}{3}{2}
\end{question}

\begin{question}{H}
A process executes the code for (i = 0; i < 3; i++) fork(); The total number of processes at the end, including the original, is
\optsshort{7}{8}{6}{3}
\end{question}

\begin{question}{H}
A program executes for (i = 0; i < n; i++) fork(); The number of child processes created is
\optsshort{n}{2n}{2\textsuperscript{n}}{2\textsuperscript{n} $-$ 1}
\end{question}

\begin{question}{H}
How many processes in total (including the original) exist after the following statement executes?\\
fork() \&\& fork() || fork();
\optsshort{8}{3}{4}{5}
\end{question}

\begin{question}{H}
How many processes in total exist after executing the following code?\\
if (fork() == 0) fork();
\optsshort{4}{2}{3}{6}
\end{question}

\begin{question}{H}
The value returned by fork() to the child process and to the parent process is respectively
\optsshort{0 and 0}{the PID of the child and 0}{0 and the PID of the child}{$-$1 and the PID of the parent}
\end{question}

\begin{question}{H}
Match List I with List II.\\
List I (Process state transition): A. Running to Ready  B. Running to Waiting  C. Waiting to Ready  D. Ready to Running\\
List II (Cause): I. The process issues an I/O request  II. Dispatch by the CPU scheduler  III. Interrupt, e.g. time quantum expiry  IV. Completion of the I/O operation
\optsshort{A-II, B-I, C-IV, D-III}{A-I, B-III, C-IV, D-II}{A-III, B-IV, C-I, D-II}{A-III, B-I, C-IV, D-II}
\end{question}

\begin{question}{H}
Which of the following process state transitions is NOT possible in the five-state model?
\optsshort{Running to Ready}{Ready to Running}{Running to Waiting}{Waiting to Running}
\end{question}

\begin{question}{H}
Match List I with List II.\\
List I (Scheduler): A. Long-term scheduler  B. Short-term scheduler  C. Medium-term scheduler\\
List II (Function): I. Selects which ready process gets the CPU next  II. Swaps processes out of and into memory  III. Controls the degree of multiprogramming by admitting jobs
\optsshort{A-III, B-I, C-II}{A-I, B-III, C-II}{A-III, B-II, C-I}{A-II, B-I, C-III}
\end{question}

\begin{question}{M}
The degree of multiprogramming is controlled by the
\optsshort{dispatcher}{short-term scheduler}{medium-term scheduler only}{long-term scheduler}
\end{question}

\begin{question}{H}
Which of the following is NOT stored in the process control block (PCB)?
\optsshort{Program counter}{Source code of the program}{Process state}{CPU register contents}
\end{question}

\begin{question}{M}
A process that has terminated but whose parent has not yet called wait() is called a
\optsshort{orphan process}{daemon process}{blocked process}{zombie process}
\end{question}

\begin{question}{M}
A process whose parent has terminated without waiting for it is called an orphan process. In Linux it is adopted by
\optsshort{init or systemd (PID 1)}{its grandparent}{no process}{the swapper process}
\end{question}

\begin{question}{M}
The exec() family of system calls
\optslong{creates a new child process}{waits for a child to terminate}{replaces the memory image of the calling process with a new program}{duplicates file descriptors}
\end{question}

\begin{question}{M}
The layout of a process in memory, from low to high addresses, is usually
\optsshort{stack, heap, data, text}{data, text, stack, heap}{heap, stack, text, data}{text, data, heap, stack}
\end{question}

\begin{question}{M}
The time taken by the dispatcher to stop one process and start another is called
\optsshort{turnaround time}{response time}{burst time}{dispatch latency}
\end{question}

\begin{question}{H}
Given below are two statements: one is labelled as Assertion A and the other is labelled as Reason R.\\
Assertion A: Context-switch time is pure overhead.\\
Reason R: The hardware saves and restores all process state automatically, without any work by the operating system.
\optslong{Both A and R are true and R is the correct explanation of A}{Both A and R are true but R is NOT the correct explanation of A}{A is true but R is false}{A is false but R is true}
\end{question}

\begin{question}{H}
Which of the following are shared by all the threads of a process?\\
A. Code section  B. Data section  C. Stack  D. Open files  E. Program counter\\
Choose the correct answer from the options given below:
\optsshort{C and E only}{A, B and D only}{A, B and C only}{A, B, C, D and E}
\end{question}

\begin{question}{H}
Match List I with List II.\\
List I (Multithreading model): A. Many-to-one  B. One-to-one  C. Many-to-many\\
List II (Property): I. Each user thread maps to its own kernel thread; used by Linux and Windows  II. Many user threads are multiplexed over an equal or smaller number of kernel threads  III. A blocking system call blocks the whole process and threads cannot run in parallel on multiple processors
\optsshort{A-I, B-III, C-II}{A-II, B-I, C-III}{A-III, B-I, C-II}{A-III, B-II, C-I}
\end{question}

\begin{question}{H}
Given below are two statements.\\
Statement I: User-level threads can be created and switched faster than kernel-level threads.\\
Statement II: In the many-to-one model, threads of the same process can run in parallel on different processors.
\optslong{Both Statement I and Statement II are true}{Both Statement I and Statement II are false}{Statement I is true but Statement II is false}{Statement I is false but Statement II is true}
\end{question}

\begin{question}{H}
According to Amdahl's law, an application that is 75\% parallel and 25\% serial, run on 4 processing cores, has a maximum speedup of approximately
\optsshort{4}{2.28}{1.6}{3}
\end{question}

\begin{question}{H}
According to Amdahl's law, if 25\% of a program is serial, the maximum speedup obtainable with an unlimited number of processors is
\optsshort{2}{unlimited}{4}{25}
\end{question}

\begin{question}{M}
Given below are two statements.\\
Statement I: Concurrency is possible on a single-core system by interleaving the execution of threads.\\
Statement II: Parallelism requires more than one processing core.
\optslong{Both Statement I and Statement II are true}{Both Statement I and Statement II are false}{Statement I is true but Statement II is false}{Statement I is false but Statement II is true}
\end{question}

\begin{question}{M}
Distributing subsets of the same data across cores and performing the same operation on each is called
\optsshort{multiprogramming}{data parallelism}{pipelining}{task parallelism}
\end{question}

\begin{question}{M}
Match List I with List II.\\
List I (Implicit threading approach): A. Thread pool  B. OpenMP  C. Fork-join  D. Grand Central Dispatch\\
List II (Description): I. Compiler directives such as \#pragma omp parallel  II. Apple's dispatch queues of blocks  III. A set of pre-created threads waiting for work  IV. A parent task splits work into subtasks and waits for them to finish
\optsshort{A-I, B-III, C-IV, D-II}{A-III, B-IV, C-I, D-II}{A-III, B-I, C-IV, D-II}{A-IV, B-I, C-III, D-II}
\end{question}

\begin{question}{M}
Given below are two statements.\\
Statement I: Pthreads is a thread library implementation available only on Linux.\\
Statement II: pthread\_\allowbreak{}join() makes the calling thread wait for the specified thread to terminate.
\optslong{Both Statement I and Statement II are true}{Both Statement I and Statement II are false}{Statement I is true but Statement II is false}{Statement I is false but Statement II is true}
\end{question}

\begin{question}{M}
In deferred thread cancellation
\optslong{the target thread periodically checks whether it should terminate, at cancellation points}{the target thread is terminated immediately}{the whole process is terminated}{the thread cannot be cancelled}
\end{question}

\begin{question}{M}
Which IPC mechanism is generally faster once established?
\optsshort{Message passing}{Pipes}{Shared memory}{Sockets}
\end{question}

\begin{question}{H}
Given below are two statements.\\
Statement I: Message passing is easier to implement than shared memory in a distributed system.\\
Statement II: Message passing is useful for exchanging smaller amounts of data because no conflicts need be avoided.
\optslong{Both Statement I and Statement II are true}{Both Statement I and Statement II are false}{Statement I is true but Statement II is false}{Statement I is false but Statement II is true}
\end{question}

\begin{question}{H}
In the bounded-buffer producer-consumer solution using a circular array of size n with in and out pointers (and no count variable), the buffer is full when
\optsshort{in == out}{(in + 1) \% n == out}{in == n}{(out + 1) \% n == in}
\end{question}

\begin{question}{M}
When both send() and receive() are blocking, the arrangement is called a
\optsshort{pipe}{rendezvous}{mailbox}{signal}
\end{question}

\begin{question}{M}
Match List I with List II.\\
List I (Buffering capacity of a message queue): A. Zero capacity  B. Bounded capacity  C. Unbounded capacity\\
List II (Behaviour of the sender): I. Never blocks  II. Must block until the receiver receives the message  III. Blocks only when the queue is full
\optsshort{A-III, B-II, C-I}{A-II, B-III, C-I}{A-I, B-III, C-II}{A-II, B-I, C-III}
\end{question}

\begin{question}{M}
In indirect communication, messages are sent to and received from
\optsshort{stacks}{process identifiers}{mailboxes (ports)}{registers}
\end{question}

\begin{question}{H}
Given below are two statements.\\
Statement I: An ordinary (anonymous) pipe in UNIX is unidirectional and requires a parent-child relationship between the communicating processes.\\
Statement II: A named pipe (FIFO) can be used only between a parent and its child process.
\optslong{Both Statement I and Statement II are true}{Both Statement I and Statement II are false}{Statement I is true but Statement II is false}{Statement I is false but Statement II is true}
\end{question}

\begin{question}{M}
After the call pipe(fd) in UNIX
\optslong{fd[0] is the read end and fd[1] is the write end}{both ends can be used for reading and writing}{fd[0] is the write end and fd[1] is the read end}{fd[0] and fd[1] are the same descriptor}
\end{question}

\begin{question}{M}
In a remote procedure call (RPC), packaging the parameters into a form that can be transmitted over the network is called
\optsshort{polling}{binding}{stubbing}{marshalling}
\end{question}

\begin{question}{M}
Which of the following signals cannot be caught or ignored by a UNIX process?
\optsshort{SIGTERM}{SIGINT}{SIGALRM}{SIGKILL}
\end{question}

\begin{question}{M}
A socket is identified by
\optslong{a process ID and a port}{a file descriptor only}{an IP address concatenated with a port number}{a MAC address}
\end{question}

\begin{question}{M}
Thread-local storage (TLS) is
\optslong{storage for thread IDs only}{data that each thread keeps its own copy of}{data shared by all threads of a process}{the stack of the kernel}
\end{question}

\begin{question}{M}
In the scheduler activation scheme, the kernel informs the user-level thread library about events such as a thread about to block by means of
\optsshort{an upcall}{a system call}{a page fault}{a signal to every thread}
\end{question}

\begin{question}{M}
Which of the following benefits of multithreading is illustrated by a web browser that remains interactive while an image is loading?
\optsshort{Responsiveness}{Resource sharing}{Economy}{Scalability}
\end{question}

\begin{question}{M}
Creating a thread is cheaper than creating a process because
\optslong{threads do not need a stack}{threads share the address space and resources of the process}{threads cannot do I/O}{threads do not run on the CPU}
\end{question}

\begin{question}{M}
Arrange the following in the order in which they occur when a process is created in UNIX and runs a new program.\\
A. The child calls exec()  B. The parent calls fork()  C. The parent's wait() returns the child's exit status  D. The child calls exit()
\optsshort{A, B, C, D}{B, C, A, D}{B, A, D, C}{B, A, C, D}
\end{question}

\begin{question}{M}
The medium-term scheduler is used mainly to
\optslong{select the next process to run}{select jobs from the job pool}{handle interrupts}{remove processes from memory to reduce the degree of multiprogramming}
\end{question}

\begin{question}{M}
A process is said to be I/O-bound if it
\optslong{never performs I/O}{spends more time doing computation than I/O}{is always in the ready queue}{spends more time doing I/O than computation}
\end{question}

\begin{question}{M}
When a process terminates, causing all its children to be terminated as well, this is called
\optsshort{orphaning}{cascading termination}{process aging}{deferred cancellation}
\end{question}

\begin{question}{M}
Which Linux system call creates a new task and, through flags such as CLONE\_\allowbreak{}VM and CLONE\_\allowbreak{}FILES, decides how much is shared with the parent?
\optsshort{clone()}{exec()}{spawn()}{fork()}
\end{question}

\begin{question}{M}
Which of the following is a lightweight process?
\optsshort{An orphan}{A zombie}{A thread}{A daemon}
\end{question}

\begin{question}{L}
A process running in the background that provides a service and is not attached to a terminal is called a
\optsshort{shell}{orphan}{zombie}{daemon}
\end{question}

\begin{question}{M}
How many processes in total exist after the following code runs?\\
if (fork()) fork(); fork();
\optsshort{6}{4}{8}{5}
\end{question}

\begin{question}{M}
How many processes in total, including the parent, are created by the following code?\\
for (i = 0; i < 4; i++) if (fork() == 0) break;
\optsshort{4}{5}{16}{8}
\end{question}

\begin{answers}
\ans{1}{1}{Each fork doubles the number of processes: 2 $\times$ 2 = 4 processes, each prints once.}
\ans{2}{2}{n forks in sequence give 2\textsuperscript{n} = 8 processes, i.e. 7 child processes.}
\ans{3}{4}{There are 2\textsuperscript{n} processes in total, one of which is the original parent.}
\ans{4}{4}{The parent's first two forks succeed and short-circuit the OR; the two children that get 0 each execute the third fork: 1 + 4 = 5.}
\ans{5}{3}{Only the child of the first fork executes the second fork, giving parent, child and grandchild.}
\ans{6}{3}{fork returns $-$1 to the parent if creation fails.}
\ans{7}{4}{Standard five-state process model transitions.}
\ans{8}{4}{A process whose I/O completes must first become Ready; only the dispatcher moves a process to Running.}
\ans{9}{1}{The long-term scheduler admits jobs, the CPU scheduler dispatches, the medium-term scheduler swaps.}
\ans{10}{4}{The long-term (job) scheduler decides how many processes are in memory.}
\ans{11}{2}{The PCB holds state, PC, registers, scheduling, memory-management, accounting and I/O status information.}
\ans{12}{4}{Its entry remains in the process table until the parent collects its exit status.}
\ans{13}{1}{init/systemd periodically calls wait() to release orphans.}
\ans{14}{3}{After a successful exec the original program no longer exists in that process.}
\ans{15}{4}{The heap grows upward and the stack grows downward toward each other.}
\ans{16}{4}{Dispatch latency includes the context switch and the switch to user mode.}
\ans{17}{3}{The kernel itself saves the context in the PCB and restores the next one; this work does nothing for any user process.}
\ans{18}{2}{Each thread has its own stack, register set and program counter.}
\ans{19}{3}{Standard properties of the three threading models.}
\ans{20}{3}{Only one thread at a time can access the kernel in the many-to-one model, so the threads cannot run in parallel.}
\ans{21}{2}{Speedup $\leq$ 1 / (0.25 + 0.75 / 4) = 1 / 0.4375 $\approx$ 2.28.}
\ans{22}{3}{As N tends to infinity the speedup approaches 1/S = 1/0.25 = 4.}
\ans{23}{1}{Concurrency means making progress on several tasks; parallelism means running them at the same instant.}
\ans{24}{2}{Task parallelism instead distributes different operations.}
\ans{25}{3}{Common implicit threading techniques.}
\ans{26}{4}{Pthreads is a POSIX specification, implemented on Linux, macOS, Solaris and other UNIX-like systems.}
\ans{27}{1}{Deferred cancellation is the Pthreads default and lets the thread release resources safely.}
\ans{28}{3}{After the region is set up, accesses need no kernel intervention; message passing needs a system call per message.}
\ans{29}{1}{Both are standard comparisons of the two IPC models.}
\ans{30}{2}{This scheme can hold at most n $-$ 1 items, because in == out is used to mean "empty".}
\ans{31}{2}{Sender and receiver meet at the point of message transfer.}
\ans{32}{2}{Zero-capacity links need a rendezvous; unbounded queues never block the sender.}
\ans{33}{3}{In direct communication, processes name each other explicitly.}
\ans{34}{3}{Named pipes appear in the file system and can be used by unrelated processes.}
\ans{35}{1}{Data written to fd[1] is read from fd[0].}
\ans{36}{4}{The client stub marshals the parameters; XDR is a machine-independent representation.}
\ans{37}{4}{SIGKILL (and SIGSTOP) always take effect.}
\ans{38}{3}{A connection is identified by the pair of sockets at its two ends.}
\ans{39}{2}{TLS persists across function calls, unlike local variables, but is private to a thread.}
\ans{40}{1}{Upcalls are handled by the thread library's upcall handler on a lightweight process.}
\ans{41}{1}{Another thread continues to respond to the user.}
\ans{42}{2}{There is no need to allocate a new address space or duplicate resources.}
\ans{43}{3}{The parent forks, the child execs and finally exits, and only then can the parent's wait() return with the exit status.}
\ans{44}{4}{Swapping out partially executed processes improves the process mix or frees memory.}
\ans{45}{4}{The long-term scheduler should admit a good mix of I/O-bound and CPU-bound processes.}
\ans{46}{2}{Some systems do not allow a child to exist if its parent has terminated.}
\ans{47}{1}{Linux does not distinguish processes from threads; clone() with sharing flags creates threads.}
\ans{48}{3}{A thread is often called a lightweight process.}
\ans{49}{4}{Examples are sshd, cron and httpd.}
\ans{50}{1}{The parent creates two children before the last fork (3 processes); each of them forks once more: 3 $\times$ 2 = 6.}
\ans{51}{2}{Each child breaks out of the loop immediately, so only the parent forks four times: 1 + 4 = 5.}
\end{answers}

\section{CPU Scheduling (criteria, algorithms, Gantt-chart numericals, multiprocessor and real-time scheduling)}

\begin{question}{H}
Three processes P1, P2 and P3 arrive at time 0 in that order with CPU burst times of 24, 3 and 3 ms. The average waiting time under FCFS scheduling is
\optsshort{17 ms}{5.67 ms}{3 ms}{27 ms}
\end{question}

\begin{question}{H}
Processes P1, P2 and P3 with burst times 24, 3 and 3 ms all arrive at time 0. If they arrive in the order P2, P3, P1, the average waiting time under FCFS is
\optsshort{17 ms}{3 ms}{10 ms}{5 ms}
\end{question}

\begin{question}{H}
Processes P1, P2 and P3 with burst times 24, 3 and 3 ms arrive at time 0 in that order. The average waiting time under round-robin scheduling with a time quantum of 4 ms is
\optsshort{6 ms}{5.67 ms}{17 ms}{4 ms}
\end{question}

\begin{question}{H}
Consider processes with (arrival time, burst time): P1 (0, 7), P2 (2, 4), P3 (4, 1), P4 (5, 4). The average waiting time under non-preemptive SJF scheduling is
\optsshort{4}{6}{3}{5}
\end{question}

\begin{question}{H}
Consider processes with (arrival time, burst time): P1 (0, 7), P2 (2, 4), P3 (4, 1), P4 (5, 4). The average waiting time under preemptive SJF (shortest remaining time first) scheduling is
\optsshort{4}{7}{5}{3}
\end{question}

\begin{question}{H}
Consider processes with (arrival time, burst time): P1 (0, 7), P2 (2, 4), P3 (4, 1), P4 (5, 4). The average turnaround time under SRTF scheduling is
\optsshort{8}{7}{6}{9}
\end{question}

\begin{question}{H}
Consider processes with (arrival time, burst time): P1 (0, 8), P2 (1, 4), P3 (2, 9), P4 (3, 5). The average waiting time under SRTF scheduling is
\optsshort{13}{7.75}{6.5}{8.75}
\end{question}

\begin{question}{H}
Consider processes with (arrival time, burst time): P1 (0, 8), P2 (1, 4), P3 (2, 9), P4 (3, 5). The average waiting time under non-preemptive SJF scheduling is
\optsshort{8.75}{7.75}{6.5}{9.25}
\end{question}

\begin{question}{M}
Consider processes with (arrival time, burst time): P1 (0, 8), P2 (1, 4), P3 (2, 9), P4 (3, 5). How many context switches occur under SRTF scheduling, not counting the initial dispatch at time 0?
\optsshort{4}{3}{5}{6}
\end{question}

\begin{question}{H}
Five processes arrive at time 0 with (burst time, priority): P1 (10, 3), P2 (1, 1), P3 (2, 4), P4 (1, 5), P5 (5, 2). A smaller number means a higher priority. The average waiting time under non-preemptive priority scheduling is
\optsshort{3.2}{12}{8.2}{9.6}
\end{question}

\begin{question}{H}
Five processes arrive at time 0 with burst times P1 = 10, P2 = 1, P3 = 2, P4 = 1 and P5 = 5. The average waiting time under SJF scheduling is
\optsshort{9.6}{8.2}{3.2}{5.4}
\end{question}

\begin{question}{M}
Five processes arrive at time 0 in the order P1 to P5 with burst times 10, 1, 2, 1 and 5. The average waiting time under round robin with a time quantum of 1 is
\optsshort{8.2}{9.6}{5.4}{3.2}
\end{question}

\begin{question}{H}
Consider processes with (arrival time, burst time, priority): P1 (0, 5, 3), P2 (1, 3, 1), P3 (2, 8, 4), P4 (3, 6, 2). A smaller number means higher priority. The average waiting time under preemptive priority scheduling is
\optsshort{5.5}{6}{11}{5.25}
\end{question}

\begin{question}{H}
Consider five jobs with (arrival time, service time): A (0, 3), B (2, 6), C (4, 4), D (6, 5), E (8, 2). Under highest response ratio next (HRRN) scheduling, the order of completion is
\optsshort{A, B, E, C, D}{A, C, B, E, D}{A, B, C, E, D}{A, B, C, D, E}
\end{question}

\begin{question}{M}
Consider five jobs with (arrival time, service time): A (0, 3), B (2, 6), C (4, 4), D (6, 5), E (8, 2). The average turnaround time under HRRN scheduling is
\optsshort{8.6}{7.6}{7.2}{8}
\end{question}

\begin{question}{M}
Consider five jobs with (arrival time, service time): A (0, 3), B (2, 6), C (4, 4), D (6, 5), E (8, 2). The average turnaround time under SRTF scheduling is
\optsshort{7.6}{8}{8.6}{7.2}
\end{question}

\begin{question}{M}
The response ratio used in highest response ratio next (HRRN) scheduling is
\optslong{(waiting time + service time) / service time}{service time / waiting time}{waiting time / service time}{(waiting time + service time) / waiting time}
\end{question}

\begin{question}{M}
Six processes with (arrival time, burst time) P1 (0, 4), P2 (1, 5), P3 (2, 2), P4 (3, 1), P5 (4, 6), P6 (6, 3) are scheduled using round robin with a quantum of 2. A newly arrived process joins the ready queue before a process whose quantum has just expired. The completion time of P2 is
\optsshort{15}{21}{18}{17}
\end{question}

\begin{question}{H}
Which scheduling algorithm gives the minimum average waiting time for a given set of processes that are all available at the same time?
\optsshort{FCFS}{Priority}{Round robin}{SJF}
\end{question}

\begin{question}{H}
Given below are two statements.\\
Statement I: SJF scheduling is optimal but difficult to implement because the length of the next CPU burst is not known in advance.\\
Statement II: The next CPU burst is usually predicted as an exponential average of the measured lengths of previous bursts.
\optslong{Both Statement I and Statement II are true}{Both Statement I and Statement II are false}{Statement I is true but Statement II is false}{Statement I is false but Statement II is true}
\end{question}

\begin{question}{H}
Using exponential averaging with $\alpha$ = 0.5 and an initial prediction $\tau$0 = 10 ms, the predicted length of the next CPU burst after observed bursts of 6 ms and then 4 ms is
\optsshort{8 ms}{6 ms}{7 ms}{5 ms}
\end{question}

\begin{question}{M}
In exponential averaging $\tau$(n+1) = $\alpha$ t(n) + (1 $-$ $\alpha$) $\tau$(n), the value $\alpha$ = 0 means that
\optslong{only the most recent burst matters}{all bursts are given equal weight}{recent history has no effect on the prediction}{the prediction is always zero}
\end{question}

\begin{question}{H}
Match List I with List II.\\
List I (Scheduling algorithm): A. FCFS  B. Round robin  C. SJF  D. Priority\\
List II (Problem/feature): I. Convoy effect  II. Needs knowledge of the next CPU burst length  III. Indefinite blocking (starvation), solved by aging  IV. Performance depends heavily on the size of the time quantum
\optsshort{A-I, B-IV, C-III, D-II}{A-I, B-IV, C-II, D-III}{A-I, B-II, C-IV, D-III}{A-IV, B-I, C-II, D-III}
\end{question}

\begin{question}{H}
Which of the following scheduling algorithms are non-preemptive?\\
A. FCFS  B. Round robin  C. SRTF  D. HRRN\\
Choose the correct answer from the options given below:
\optsshort{A, B and D only}{B and C only}{A and D only}{A, C and D only}
\end{question}

\begin{question}{H}
Aging is a technique used to
\optslong{decrease the quantum of long jobs}{reduce context-switch overhead}{terminate old processes}{gradually increase the priority of processes that wait in the system for a long time}
\end{question}

\begin{question}{H}
If the time quantum in round-robin scheduling is made very large, the algorithm behaves like
\optsshort{multilevel queue scheduling}{FCFS}{priority scheduling}{SJF}
\end{question}

\begin{question}{M}
If the time quantum in round-robin scheduling is very small
\optslong{throughput increases indefinitely}{starvation occurs}{the algorithm becomes FCFS}{context-switch overhead becomes very high}
\end{question}

\begin{question}{M}
In round-robin scheduling with 5 processes in the ready queue and a time quantum of 20 ms, the maximum time a process waits before getting its next quantum is
\optsshort{80 ms}{120 ms}{20 ms}{100 ms}
\end{question}

\begin{question}{M}
In round robin, each time quantum is 10 ms and each context switch takes 1 ms. If all processes use their full quantum, the CPU efficiency is approximately
\optsshort{100\%}{90.9\%}{81.8\%}{99\%}
\end{question}

\begin{question}{H}
Match List I with List II.\\
List I (Scheduling criterion): A. Throughput  B. Turnaround time  C. Waiting time  D. Response time\\
List II (Definition): I. Time from submission until the first response is produced  II. Number of processes completed per unit time  III. Total time spent in the ready queue  IV. Time from submission to completion
\optsshort{A-II, B-III, C-IV, D-I}{A-II, B-IV, C-III, D-I}{A-II, B-IV, C-I, D-III}{A-IV, B-II, C-III, D-I}
\end{question}

\begin{question}{H}
Which of the following scheduling criteria should be maximised?\\
A. CPU utilisation  B. Throughput  C. Turnaround time  D. Waiting time  E. Response time\\
Choose the correct answer from the options given below:
\optsshort{C, D and E only}{A and B only}{A, B and C only}{B and E only}
\end{question}

\begin{question}{M}
Turnaround time is equal to
\optslong{arrival time + burst time}{waiting time $-$ burst time}{completion time $-$ burst time}{completion time $-$ arrival time}
\end{question}

\begin{question}{M}
CPU scheduling decisions are preemptive when a process switches
\optslong{from running to ready state}{from running to terminated state}{from running to waiting state}{from new to ready state}
\end{question}

\begin{question}{H}
Given below are two statements.\\
Statement I: In multilevel queue scheduling, processes move between queues according to their CPU-burst behaviour.\\
Statement II: In multilevel feedback queue scheduling, a process is permanently assigned to one queue when it enters the system.
\optslong{Both Statement I and Statement II are true}{Both Statement I and Statement II are false}{Statement I is true but Statement II is false}{Statement I is false but Statement II is true}
\end{question}

\begin{question}{M}
In a multilevel feedback queue scheduler, a process that uses too much CPU time is
\optslong{terminated}{moved to a lower-priority queue}{moved to the waiting queue}{moved to a higher-priority queue}
\end{question}

\begin{question}{M}
Which of the following is NOT a parameter needed to define a multilevel feedback queue scheduler?
\optslong{The scheduling algorithm for each queue}{The method used to demote or upgrade a process}{The number of queues}{The size of main memory}
\end{question}

\begin{question}{M}
The convoy effect, in which many short processes wait behind one long process, is associated with
\optsshort{round robin}{FCFS}{SRTF}{SJF}
\end{question}

\begin{question}{M}
In a system, each process spends 80\% of its time waiting for I/O. With 4 processes in memory, the approximate CPU utilisation is
\optsshort{80\%}{59\%}{41\%}{20\%}
\end{question}

\begin{question}{H}
Given below are two statements.\\
Statement I: Rate-monotonic scheduling assigns static priorities, with a shorter period meaning a higher priority.\\
Statement II: Earliest-deadline-first scheduling assigns priorities dynamically according to deadlines.
\optslong{Both Statement I and Statement II are true}{Both Statement I and Statement II are false}{Statement I is true but Statement II is false}{Statement I is false but Statement II is true}
\end{question}

\begin{question}{H}
The worst-case CPU utilisation bound that guarantees schedulability of two processes under rate-monotonic scheduling is approximately
\optsshort{100\%}{83\%}{69\%}{78\%}
\end{question}

\begin{question}{M}
Two periodic tasks have (period, execution time) of (50, 20) and (100, 35). Their total CPU utilisation is
\optsshort{75\%}{67.5\%}{85\%}{55\%}
\end{question}

\begin{question}{H}
Given below are two statements: one is labelled as Assertion A and the other is labelled as Reason R.\\
Assertion A: Two tasks with (period, execution time) of (50, 25) and (80, 35) can be scheduled by EDF but may miss deadlines under rate-monotonic scheduling.\\
Reason R: Their utilisation of about 94\% is at most 100\%, but exceeds the rate-monotonic bound of about 83\% for two tasks.
\optslong{Both A and R are true and R is the correct explanation of A}{Both A and R are true but R is NOT the correct explanation of A}{A is true but R is false}{A is false but R is true}
\end{question}

\begin{question}{M}
The upper bound on CPU utilisation for a set of periodic tasks to be schedulable under EDF is
\optsshort{83\%}{100\%}{There is no bound}{69\%}
\end{question}

\begin{question}{H}
Priority inversion occurs when
\optslong{two processes have the same priority}{a low-priority process holds a resource needed by a high-priority process, and a medium-priority process preempts the low-priority one}{a high-priority process always runs first}{the scheduler reverses the priority queue}
\end{question}

\begin{question}{M}
Given below are two statements.\\
Statement I: Processor affinity means that a process prefers to keep running on the processor on which it is currently executing.\\
Statement II: Processor affinity increases cache misses because the process keeps moving between processors.
\optslong{Both Statement I and Statement II are true}{Both Statement I and Statement II are false}{Statement I is true but Statement II is false}{Statement I is false but Statement II is true}
\end{question}

\begin{question}{M}
In multiprocessor load balancing, a periodic task that moves processes from overloaded to idle processors is called
\optsshort{push migration}{gang scheduling}{pull migration}{processor affinity}
\end{question}

\begin{question}{M}
Under process-contention scope (PCS), threads
\optslong{compete for the CPU with other threads of the same process}{are not scheduled at all}{compete for the CPU with all threads in the system}{are scheduled only by the kernel}
\end{question}

\begin{question}{M}
Which of the following real-time scheduling policies are defined by POSIX?
\optsshort{SCHED\_\allowbreak{}LOTTERY and SCHED\_\allowbreak{}FAIR}{SCHED\_\allowbreak{}FIFO and SCHED\_\allowbreak{}RR}{SCHED\_\allowbreak{}EDF and SCHED\_\allowbreak{}RM}{SCHED\_\allowbreak{}SJF and SCHED\_\allowbreak{}HRRN}
\end{question}

\begin{question}{M}
In lottery scheduling, a process with more tickets
\optslong{is never preempted}{gets a longer quantum}{has a proportionally higher chance of getting the CPU}{always runs first}
\end{question}

\begin{question}{M}
Which component gives control of the CPU to the process selected by the short-term scheduler?
\optsshort{Loader}{Dispatcher}{Long-term scheduler}{Interrupt handler}
\end{question}

\begin{question}{M}
An I/O-bound process typically has
\optsshort{exactly one CPU burst}{many short CPU bursts}{no CPU bursts}{a few very long CPU bursts}
\end{question}

\begin{question}{H}
Arrange the following algorithms in increasing order of average waiting time for processes with (arrival, burst): P1 (0, 7), P2 (2, 4), P3 (4, 1), P4 (5, 4).\\
A. FCFS  B. SJF  C. SRTF
\optsshort{C, A, B}{B, C, A}{A, B, C}{C, B, A}
\end{question}

\begin{question}{M}
Which scheduling algorithm can lead to starvation?\\
A. FCFS  B. SJF  C. Priority  D. Round robin\\
Choose the correct answer from the options given below:
\optsshort{A and D only}{B and C only}{A, B, C and D}{A, B and C only}
\end{question}

\begin{question}{L}
A scheduling algorithm that assigns each process a fixed share of CPU time, e.g. 50\% to one user's processes, is called
\optsshort{fair-share scheduling}{deadline scheduling}{SJF scheduling}{FCFS scheduling}
\end{question}

\begin{question}{M}
Simultaneous multithreading (hyper-threading) allows
\optslong{a process to run on two machines}{one core to appear as multiple logical processors to the operating system}{multiple cores to share one cache line}{a thread to run without a stack}
\end{question}

\begin{answers}
\ans{1}{1}{Waiting times are 0, 24 and 27 ms; the average is 51/3 = 17 ms.}
\ans{2}{2}{Waiting times are 0, 3 and 6 ms; the average is 3 ms, showing how much FCFS depends on arrival order.}
\ans{3}{2}{Gantt chart: P1 0--4, P2 4--7, P3 7--10, P1 10--30. Waiting times are 6, 4 and 7; the average is 17/3 $\approx$ 5.67 ms.}
\ans{4}{1}{Order P1 (0--7), P3 (7--8), P2 (8--12), P4 (12--16). Waiting times 0, 6, 3, 7; average 16/4 = 4.}
\ans{5}{4}{Gantt: P1 0--2, P2 2--4, P3 4--5, P2 5--7, P4 7--11, P1 11--16. Waiting times 9, 1, 0, 2; average 12/4 = 3.}
\ans{6}{2}{Completion times are 16, 7, 5 and 11; turnaround times 16, 5, 1, 6; average 28/4 = 7.}
\ans{7}{3}{Gantt: P1 0--1, P2 1--5, P4 5--10, P1 10--17, P3 17--26. Waiting times 9, 0, 15, 2; average 26/4 = 6.5.}
\ans{8}{2}{Order P1 0--8, P2 8--12, P4 12--17, P3 17--26. Waiting times 0, 7, 15, 9; average 31/4 = 7.75.}
\ans{9}{1}{The CPU passes P1 $\rightarrow$ P2 $\rightarrow$ P4 $\rightarrow$ P1 $\rightarrow$ P3, which is 4 switches.}
\ans{10}{3}{Order P2, P5, P1, P3, P4 with waiting times 0, 1, 6, 16, 18; average 41/5 = 8.2.}
\ans{11}{3}{Order P2, P4, P3, P5, P1 with waiting times 0, 1, 2, 4, 9; average 16/5 = 3.2.}
\ans{12}{3}{Completion times are 19, 2, 7, 4 and 14; waiting times 9, 1, 5, 3, 9; average 27/5 = 5.4.}
\ans{13}{1}{Gantt: P1 0--1, P2 1--4, P4 4--10, P1 10--14, P3 14--22. Waiting times 9, 0, 12, 1; average 22/4 = 5.5.}
\ans{14}{3}{At t = 9 the ratios are C 2.25, D 1.6, E 1.5, so C runs; at t = 13 they are D 2.4 and E 3.5, so E runs before D.}
\ans{15}{4}{Completion times 3, 9, 13, 20, 15; turnaround times 3, 7, 9, 14, 7; average 40/5 = 8.}
\ans{16}{4}{Gantt: A 0--3, B 3--4, C 4--8, E 8--10, B 10--15, D 15--20. Turnaround 3, 13, 4, 14, 2; average 36/5 = 7.2.}
\ans{17}{1}{It favours short jobs, but the ratio of a long job grows as it waits, so there is no starvation.}
\ans{18}{3}{Gantt: P1 0--2, P2 2--4, P3 4--6, P1 6--8, P4 8--9, P5 9--11, P2 11--13, P6 13--15, P5 15--17, P2 17--18, P6 18--19, P5 19--21.}
\ans{19}{4}{SJF is provably optimal for average waiting time; its preemptive version SRTF is optimal when arrivals differ.}
\ans{20}{1}{$\tau$(n+1) = $\alpha$ t(n) + (1 $-$ $\alpha$) $\tau$(n).}
\ans{21}{2}{$\tau$1 = 0.5 $\times$ 6 + 0.5 $\times$ 10 = 8; $\tau$2 = 0.5 $\times$ 4 + 0.5 $\times$ 8 = 6 ms.}
\ans{22}{3}{With $\alpha$ = 0 the prediction never changes; with $\alpha$ = 1 only the last burst counts.}
\ans{23}{2}{Standard characteristics of each algorithm.}
\ans{24}{3}{Round robin and SRTF preempt the running process.}
\ans{25}{4}{Aging prevents starvation in priority scheduling.}
\ans{26}{2}{Every process finishes its burst within one quantum, so RR reduces to FCFS.}
\ans{27}{4}{Most CPU time is then spent switching (processor sharing in the limit).}
\ans{28}{1}{A process waits at most (n $-$ 1) $\times$ q = 4 $\times$ 20 = 80 ms.}
\ans{29}{2}{Useful time / total time = 10 / (10 + 1) $\approx$ 90.9\%.}
\ans{30}{2}{Standard scheduling criteria definitions.}
\ans{31}{2}{Utilisation and throughput are maximised; the time measures are minimised.}
\ans{32}{4}{Waiting time = turnaround time $-$ burst time.}
\ans{33}{1}{Switching from running to waiting or terminating is voluntary (non-preemptive); running to ready and waiting to ready allow preemption.}
\ans{34}{2}{The two descriptions are swapped: multilevel queues assign processes permanently, while feedback queues let them move.}
\ans{35}{2}{I/O-bound and interactive processes stay in the higher-priority queues.}
\ans{36}{4}{Memory size is unrelated to the definition of an MLFQ scheduler.}
\ans{37}{2}{It lowers CPU and device utilisation.}
\ans{38}{2}{Utilisation = 1 $-$ p\textsuperscript{n} = 1 $-$ 0.8\textsuperscript{4} = 1 $-$ 0.4096 $\approx$ 59\%.}
\ans{39}{1}{RM is optimal among static-priority algorithms; EDF is dynamic.}
\ans{40}{2}{Bound = n(2\textsuperscript{1/n} $-$ 1) = 2($\surd$2 $-$ 1) $\approx$ 0.828 for n = 2; it tends to ln 2 $\approx$ 69\% as n grows.}
\ans{41}{1}{20/50 + 35/100 = 0.40 + 0.35 = 0.75, which is below the RM bound of 0.828, so they are schedulable under RM.}
\ans{42}{1}{EDF can theoretically achieve 100\% utilisation; above the RM bound, schedulability under RM is not guaranteed, and this set does miss a deadline under RM.}
\ans{43}{2}{EDF is optimal for preemptive uniprocessor scheduling of periodic tasks.}
\ans{44}{2}{It is solved by the priority-inheritance protocol (famous in the Mars Pathfinder case).}
\ans{45}{3}{Affinity avoids migration, so the process can reuse the contents of its processor's cache.}
\ans{46}{1}{In pull migration, an idle processor pulls a waiting task from a busy one.}
\ans{47}{1}{System-contention scope (SCS) competes among all threads in the system; one-to-one systems such as Linux use only SCS.}
\ans{48}{2}{SCHED\_\allowbreak{}OTHER is the implementation-defined non-real-time policy.}
\ans{49}{3}{It is a probabilistic proportional-share scheduler.}
\ans{50}{2}{It switches context, switches to user mode and jumps to the proper location.}
\ans{51}{2}{CPU-bound processes have a few long CPU bursts.}
\ans{52}{4}{SRTF gives 3, SJF gives 4 and FCFS gives (0 + 5 + 7 + 7)/4 = 4.75.}
\ans{53}{2}{FCFS and RR serve every process in turn; SJF and priority can postpone a process indefinitely.}
\ans{54}{1}{Fair-share scheduling allocates CPU among users or groups rather than individual processes.}
\ans{55}{2}{Each hardware thread has its own architectural state; the OS schedules them as separate CPUs.}
\end{answers}

\section{Process Synchronization (critical section, Peterson's solution, semaphores, monitors, classical problems)}

\begin{question}{H}
Which three requirements must a solution to the critical-section problem satisfy?
\optslong{Progress, starvation and deadlock}{Mutual exclusion, hold and wait, and no preemption}{Atomicity, consistency and isolation}{Mutual exclusion, progress and bounded waiting}
\end{question}

\begin{question}{H}
A counting semaphore S is initialised to 10. Then 6 wait (P) operations and 4 signal (V) operations are performed on it. The final value of S is
\optsshort{20}{8}{4}{0}
\end{question}

\begin{question}{H}
A counting semaphore is initialised to 7. Then 20 P (wait) operations and 15 V (signal) operations are completed on it. The final value of the semaphore is
\optsshort{0}{12}{5}{2}
\end{question}

\begin{question}{H}
In the semaphore implementation in which processes block on a waiting queue, a semaphore value of $-$3 means that
\optslong{3 signal operations are pending}{the semaphore is corrupted}{3 processes are waiting on the semaphore}{3 resources are available}
\end{question}

\begin{question}{H}
Given below are two statements.\\
Statement I: Peterson's solution satisfies mutual exclusion, progress and bounded waiting for two processes.\\
Statement II: Peterson's solution is guaranteed to work on modern architectures that reorder memory operations, without memory barriers.
\optslong{Both Statement I and Statement II are true}{Both Statement I and Statement II are false}{Statement I is true but Statement II is false}{Statement I is false but Statement II is true}
\end{question}

\begin{question}{H}
In Peterson's solution for two processes, the shared variables are
\optslong{a single integer lock}{boolean flag[2] and int turn}{a monitor and a condition variable}{a binary semaphore and a counting semaphore}
\end{question}

\begin{question}{H}
A simple strict-alternation solution using a single shared variable turn (each process waits until turn equals its own number) violates which requirement?
\optsshort{Bounded waiting}{Progress}{Atomicity}{Mutual exclusion}
\end{question}

\begin{question}{M}
Using a single shared lock variable that is tested and then set in two separate (non-atomic) steps can violate
\optsshort{mutual exclusion}{progress}{bounded waiting}{portability}
\end{question}

\begin{question}{H}
Given below are two statements: one is labelled as Assertion A and the other is labelled as Reason R.\\
Assertion A: Disabling interrupts is not a suitable way to implement critical sections on a multiprocessor.\\
Reason R: Disabling interrupts on one processor does not stop processes on other processors from accessing shared data, and it is time-consuming to send the message to all processors.
\optslong{Both A and R are true and R is the correct explanation of A}{Both A and R are true but R is NOT the correct explanation of A}{A is true but R is false}{A is false but R is true}
\end{question}

\begin{question}{H}
Match List I with List II.\\
List I (Synchronization tool): A. Spinlock  B. Counting semaphore  C. Binary semaphore  D. Monitor\\
List II (Description): I. Behaves like a mutex lock with values 0 and 1  II. High-level construct in which only one process is active at a time  III. Busy-waits until the lock becomes available  IV. Controls access to a resource with a finite number of instances
\optsshort{A-III, B-IV, C-I, D-II}{A-III, B-I, C-IV, D-II}{A-III, B-IV, C-II, D-I}{A-IV, B-III, C-I, D-II}
\end{question}

\begin{question}{M}
Spinlocks are most appropriate when
\optslong{the lock is held for a short time on a multiprocessor}{the process needs to perform I/O while holding the lock}{the lock is held for a long time on a uniprocessor}{there is only one process}
\end{question}

\begin{question}{M}
The test\_\allowbreak{}and\_\allowbreak{}set() hardware instruction, used in the simple lock algorithm, guarantees mutual exclusion but does NOT satisfy
\optsshort{mutual exclusion}{bounded waiting}{atomicity}{progress}
\end{question}

\begin{question}{M}
The compare\_\allowbreak{}and\_\allowbreak{}swap(value, expected, new\_\allowbreak{}value) instruction sets value to new\_\allowbreak{}value only if
\optslong{value is zero}{the lock is held}{new\_\allowbreak{}value is greater than value}{value equals expected}
\end{question}

\begin{question}{H}
In the bounded-buffer producer-consumer problem with n buffers, the semaphores mutex, empty and full are initialised to
\optsshort{1, n and 0}{n, 1 and 0}{0, n and 1}{1, 0 and n}
\end{question}

\begin{question}{H}
Given below are two statements: one is labelled as Assertion A and the other is labelled as Reason R.\\
Assertion A: In the producer-consumer solution, swapping the order of the two signal() operations at the end of the producer causes deadlock.\\
Reason R: A producer that blocks on wait(empty) while holding mutex prevents the consumer from ever entering to free a slot.
\optslong{Both A and R are true and R is the correct explanation of A}{Both A and R are true but R is NOT the correct explanation of A}{A is true but R is false}{A is false but R is true}
\end{question}

\begin{question}{H}
In the first readers-writers problem (readers have priority), which process may starve?
\optsshort{Readers}{Writers}{Both readers and writers}{Neither}
\end{question}

\begin{question}{M}
In the first readers-writers solution, the variables rw\_\allowbreak{}mutex, mutex and read\_\allowbreak{}count are initialised to
\optsshort{0, 1 and 1}{1, 1 and 1}{1, 1 and 0}{1, 0 and 0}
\end{question}

\begin{question}{H}
In the dining-philosophers problem with five philosophers, if each philosopher picks up the left chopstick first and then the right one
\optslong{the problem is solved correctly}{only one philosopher can eat at a time}{deadlock can occur when all pick up their left chopsticks at the same time}{starvation is impossible}
\end{question}

\begin{question}{H}
Which of the following are remedies for deadlock in the dining-philosophers problem?\\
A. Allow at most four philosophers at the table at the same time  B. Allow a philosopher to pick up chopsticks only if both are available  C. Odd-numbered philosophers pick up the left chopstick first and even-numbered the right one first  D. Increase the thinking time of every philosopher\\
Choose the correct answer from the options given below:
\optsshort{A, B and C only}{B, C and D only}{A and D only}{A, B, C and D}
\end{question}

\begin{question}{H}
Given below are two statements.\\
Statement I: Within a monitor, only one process can be active at a time.\\
Statement II: If x.signal() is executed when no process is waiting on x, the signal is remembered and the next x.wait() returns immediately.
\optslong{Both Statement I and Statement II are true}{Both Statement I and Statement II are false}{Statement I is true but Statement II is false}{Statement I is false but Statement II is true}
\end{question}

\begin{question}{M}
Two processes P1 and P2 share a variable x, initially 0. P1 executes x = x + 1 and P2 executes x = x $-$ 1, each implemented as load, modify and store, with arbitrary interleaving. The set of possible final values of x is
\optsshort{\{$-$1, 1\}}{\{0\}}{\{$-$2, 0, 2\}}{\{$-$1, 0, 1\}}
\end{question}

\begin{question}{M}
Two processes share a variable x, initially 5. Each executes x = x + 1 once, implemented as load, increment and store. The possible final values of x are
\optsshort{5, 6 or 7}{6 or 7}{only 7}{only 6}
\end{question}

\begin{question}{H}
To make sure that statement S1 in process P1 executes before statement S2 in process P2, using a semaphore synch, the initial value of synch and the code are
\optslong{synch = 1; P1: S1; signal(synch);  P2: wait(synch); S2;}{synch = 0; P1: S1; signal(synch);  P2: wait(synch); S2;}{synch = 1; P1: wait(synch); S1;  P2: signal(synch); S2;}{synch = 0; P1: wait(synch); S1;  P2: S2; signal(synch);}
\end{question}

\begin{question}{M}
A situation in which several processes access and manipulate the same data concurrently and the outcome depends on the order of access is called a
\optsshort{deadlock}{race condition}{livelock}{starvation}
\end{question}

\begin{question}{M}
Given below are two statements.\\
Statement I: A binary semaphore can be used as a mutex lock.\\
Statement II: A counting semaphore can be implemented using binary semaphores.
\optslong{Both Statement I and Statement II are true}{Both Statement I and Statement II are false}{Statement I is true but Statement II is false}{Statement I is false but Statement II is true}
\end{question}

\begin{question}{M}
In the monitor signalling discipline "signal and wait"
\optslong{the signal is ignored}{the signalling process waits until the signalled process leaves the monitor or waits on another condition}{both processes continue together}{the signalling process continues and the signalled process waits}
\end{question}

\begin{question}{M}
Which classical synchronization problem models a shop where a single worker sleeps when there are no customers and customers leave if all waiting chairs are full?
\optsshort{Cigarette smokers}{Sleeping barber}{Readers-writers}{Dining philosophers}
\end{question}

\begin{question}{M}
Which of the following is a liveness failure?\\
A. Deadlock  B. Starvation  C. Priority inversion\\
Choose the correct answer from the options given below:
\optsshort{B and C only}{A, B and C}{A only}{A and B only}
\end{question}

\begin{question}{M}
A semaphore S used for mutual exclusion should be initialised to
\optsshort{the number of processes}{$-$1}{0}{1}
\end{question}

\begin{question}{H}
Given below are two statements: one is labelled as Assertion A and the other is labelled as Reason R.\\
Assertion A: wait() and signal() operations on a semaphore must be executed atomically.\\
Reason R: If two processes modify the semaphore value at the same time, the update may be lost.
\optslong{Both A and R are true and R is the correct explanation of A}{Both A and R are true but R is NOT the correct explanation of A}{A is true but R is false}{A is false but R is true}
\end{question}

\begin{question}{M}
Busy waiting wastes CPU cycles because
\optslong{the process loops continuously checking a condition}{the process is blocked}{the process is swapped out}{the process makes a system call}
\end{question}

\begin{question}{M}
Which of the following hardware features is used to guarantee that changes made by one processor become visible to other processors before later operations?
\optsshort{Translation lookaside buffer}{Memory barrier (fence)}{Base register}{Interrupt vector}
\end{question}

\begin{question}{M}
A semaphore whose value can range over an unrestricted domain is called a
\optsshort{spinlock}{binary semaphore}{mutex}{counting semaphore}
\end{question}

\begin{question}{M}
Three processes each need the shared resource R, which has two instances, protected by a counting semaphore. The semaphore should be initialised to
\optsshort{1}{3}{0}{2}
\end{question}

\begin{question}{M}
Given below are two statements.\\
Statement I: A deadlock-free solution can still lead to starvation.\\
Statement II: A deadlock-free solution is necessarily starvation-free.
\optslong{Both Statement I and Statement II are true}{Both Statement I and Statement II are false}{Statement I is true but Statement II is false}{Statement I is false but Statement II is true}
\end{question}

\begin{question}{M}
Which of the following is NOT a classical problem of synchronization?
\optsshort{Travelling salesman problem}{Dining-philosophers problem}{Bounded-buffer problem}{Readers-writers problem}
\end{question}

\begin{question}{M}
Lock-free algorithms based on compare-and-swap avoid
\optslong{shared memory}{all forms of contention}{the need for atomic instructions}{the use of locks, so a thread cannot block other threads by holding a lock}
\end{question}

\begin{question}{M}
The wait() and signal() semaphore operations were originally called P and V by
\optsshort{Dijkstra}{Peterson}{Lamport}{Hoare}
\end{question}

\begin{question}{M}
Monitors were proposed as a synchronization construct by
\optsshort{Lamport and Bakery}{Coffman and Denning}{Hoare and Brinch Hansen}{Dijkstra and Peterson}
\end{question}

\begin{question}{M}
Lamport's bakery algorithm solves the critical-section problem for
\optslong{n processes using a numbered ticket for each}{only reader processes}{exactly two processes}{only multiprocessor systems with hardware support}
\end{question}

\begin{question}{H}
Given below are two statements.\\
Statement I: A mutex lock must be released only by the thread that acquired it.\\
Statement II: A semaphore can be signalled by a thread other than the one that performed wait on it.
\optslong{Both Statement I and Statement II are true}{Both Statement I and Statement II are false}{Statement I is true but Statement II is false}{Statement I is false but Statement II is true}
\end{question}

\begin{question}{M}
In a producer-consumer solution with 1 producer, 1 consumer and a buffer of 5 slots, the producer has inserted 7 items and the consumer has removed 4. The current values of the semaphores full and empty are
\optsshort{4 and 1}{3 and 2}{7 and 4}{2 and 3}
\end{question}

\begin{question}{L}
The cigarette smokers problem was proposed by
\optsshort{Hoare}{Knuth}{Dijkstra}{Patil}
\end{question}

\begin{question}{M}
Which of the following conditions can make a solution that satisfies mutual exclusion still be unacceptable?
\optslong{It does not depend on CPU speed}{It forces processes to enter strictly alternately even when one does not want to enter}{It guarantees bounded waiting}{It guarantees progress}
\end{question}

\begin{question}{M}
A condition variable in a monitor is used with the operations
\optsshort{lock and unlock}{P and V}{wait and signal}{send and receive}
\end{question}

\begin{answers}
\ans{1}{4}{These are the standard requirements; the solution must not assume anything about the relative speeds of processes.}
\ans{2}{2}{10 $-$ 6 + 4 = 8.}
\ans{3}{4}{7 $-$ 20 + 15 = 2.}
\ans{4}{3}{The magnitude of a negative value equals the number of processes blocked on it.}
\ans{5}{3}{Instruction reordering by compilers and processors can break it unless memory barriers are used.}
\ans{6}{2}{flag[i] shows that Pi wants to enter; turn shows whose turn it is.}
\ans{7}{2}{A process outside its critical section can block the other one from entering if it is not its turn.}
\ans{8}{1}{Both processes can see the lock free before either sets it.}
\ans{9}{1}{It works only on a uniprocessor, and even there it is unwise for user programs.}
\ans{10}{1}{Standard descriptions of synchronization primitives.}
\ans{11}{1}{Spinning avoids a context switch, which pays off only when the wait is short and another CPU can release the lock.}
\ans{12}{2}{A process may keep losing the race; a more complex algorithm with a waiting array is needed for bounded waiting.}
\ans{13}{4}{It always returns the original value and is executed atomically.}
\ans{14}{1}{empty counts empty slots, full counts full slots and mutex protects the buffer.}
\ans{15}{4}{The order of signal operations is harmless; it is swapping wait(mutex) and wait(empty) that can deadlock, as R describes.}
\ans{16}{2}{A continuous stream of readers can keep a writer waiting forever.}
\ans{17}{3}{The first reader locks rw\_\allowbreak{}mutex and the last reader releases it.}
\ans{18}{3}{Each holds one chopstick and waits for another: a circular wait.}
\ans{19}{1}{Changing the thinking time does not remove the possibility of circular wait.}
\ans{20}{3}{A condition-variable signal with no waiting process has no effect; unlike a semaphore, it is not remembered.}
\ans{21}{4}{Without interleaving the result is 0; if one store overwrites the other, the result is 1 or $-$1.}
\ans{22}{2}{A lost update gives 6; serial execution gives 7.}
\ans{23}{2}{P2 blocks until P1 has finished S1 and signalled.}
\ans{24}{2}{Synchronization is needed to prevent race conditions.}
\ans{25}{1}{Both are standard results.}
\ans{26}{2}{In "signal and continue" the signalling process continues instead.}
\ans{27}{2}{The sleeping barber problem uses semaphores for customers, the barber and a mutex.}
\ans{28}{2}{Each of them can prevent a process from making progress.}
\ans{29}{4}{One process can enter; the next one blocks until the first signals.}
\ans{30}{1}{The semaphore operations are themselves critical sections.}
\ans{31}{1}{A spinning process uses the CPU without doing useful work.}
\ans{32}{2}{Memory barriers prevent reordering across the barrier.}
\ans{33}{4}{Binary semaphores take only the values 0 and 1.}
\ans{34}{4}{The initial value equals the number of available instances.}
\ans{35}{3}{Deadlock-freedom guarantees that someone progresses, not that every process does.}
\ans{36}{1}{The travelling salesman problem is an optimisation problem.}
\ans{37}{4}{They rely on atomic CAS retries instead of locks.}
\ans{38}{1}{P and V come from the Dutch words proberen (to test) and verhogen (to increment).}
\ans{39}{3}{Hoare and Brinch Hansen developed the monitor concept in the 1970s.}
\ans{40}{1}{Each process takes a number; the smallest number (ties broken by process ID) enters first.}
\ans{41}{1}{This is why semaphores can be used for ordering (signalling) between threads.}
\ans{42}{2}{There are 7 $-$ 4 = 3 items in the buffer, so full = 3 and empty = 5 $-$ 3 = 2.}
\ans{43}{4}{Suhas Patil posed it to show the limits of semaphores without conditional constructs.}
\ans{44}{2}{Strict alternation violates the progress requirement.}
\ans{45}{3}{x.wait() suspends the caller until another process calls x.signal().}
\end{answers}

\section{Deadlocks (characterization, resource-allocation graphs, prevention, avoidance and Banker's algorithm, detection, recovery)}

\begin{question}{H}
Which of the following conditions must hold simultaneously for a deadlock to occur?\\
A. Mutual exclusion  B. Hold and wait  C. No preemption  D. Circular wait  E. Starvation\\
Choose the correct answer from the options given below:
\optsshort{A, B and C only}{B, C, D and E only}{A, B, C and D only}{A, B, C, D and E}
\end{question}

\begin{question}{H}
A system has 3 processes sharing 4 instances of a resource type. Each process needs at most 2 instances. Which of the following is true?
\optsshort{The system is deadlock-free}{At most one process can run}{Deadlock can occur}{Deadlock occurs always}
\end{question}

\begin{question}{H}
A system has 6 processes, each of which may need up to 3 tape drives. The minimum number of tape drives that guarantees freedom from deadlock is
\optsshort{18}{12}{13}{19}
\end{question}

\begin{question}{H}
A system has 6 identical resources and n processes competing for them. Each process can request at most 2 resources. The maximum value of n for which the system is guaranteed to be deadlock-free is
\optsshort{3}{4}{5}{6}
\end{question}

\begin{question}{H}
Five processes need at most 2, 3, 4, 5 and 6 instances of a resource respectively. The minimum number of instances that guarantees no deadlock is
\optsshort{16}{21}{15}{20}
\end{question}

\begin{question}{H}
Given below are two statements.\\
Statement I: If a resource-allocation graph contains no cycle, there is no deadlock.\\
Statement II: If a resource-allocation graph contains a cycle and every resource type has several instances, a deadlock must exist.
\optslong{Both Statement I and Statement II are true}{Both Statement I and Statement II are false}{Statement I is true but Statement II is false}{Statement I is false but Statement II is true}
\end{question}

\begin{question}{H}
If every resource type has exactly one instance, a cycle in the resource-allocation graph is
\optslong{both necessary and sufficient for deadlock}{neither necessary nor sufficient}{necessary but not sufficient for deadlock}{sufficient but not necessary}
\end{question}

\begin{question}{M}
In a resource-allocation graph, an edge from a process Pi to a resource type Rj is called a
\optsshort{claim edge}{assignment edge}{wait-for edge}{request edge}
\end{question}

\begin{question}{H}
Match List I with List II.\\
List I (Deadlock prevention): A. Deny mutual exclusion  B. Deny hold and wait  C. Deny no preemption  D. Deny circular wait\\
List II (Technique): I. Impose a total ordering on resource types and request them in increasing order  II. Request all resources before starting execution  III. Take resources away from a process that waits for another resource  IV. Make resources shareable, e.g. by spooling a printer
\optsshort{A-IV, B-II, C-III, D-I}{A-IV, B-III, C-II, D-I}{A-II, B-IV, C-III, D-I}{A-IV, B-II, C-I, D-III}
\end{question}

\begin{question}{H}
Given below are two statements: one is labelled as Assertion A and the other is labelled as Reason R.\\
Assertion A: Requesting all resources before execution prevents deadlock but can cause low resource utilisation and starvation.\\
Reason R: Resources may be allocated long before they are used, and a process needing many popular resources may wait indefinitely.
\optslong{Both A and R are true and R is the correct explanation of A}{Both A and R are true but R is NOT the correct explanation of A}{A is true but R is false}{A is false but R is true}
\end{question}

\begin{question}{H}
Which of the following statements about safe and unsafe states is correct?
\optslong{An unsafe state is always a deadlocked state}{A safe state can never lead to deadlock if the OS follows the avoidance algorithm, while an unsafe state may lead to deadlock}{A deadlocked state is a safe state}{Safe and unsafe states are the same}
\end{question}

\begin{question}{H}
In the Banker's algorithm, the Need matrix is calculated as
\optsshort{Max $-$ Allocation}{Allocation $-$ Max}{Available $-$ Allocation}{Max + Allocation}
\end{question}

\begin{question}{H}
Five processes P0 to P4 share resources A, B and C (10, 5 and 7 instances). Allocation: P0 (0,1,0), P1 (2,0,0), P2 (3,0,2), P3 (2,1,1), P4 (0,0,2). Max: P0 (7,5,3), P1 (3,2,2), P2 (9,0,2), P3 (2,2,2), P4 (4,3,3). Available is (3,3,2). The Need of P4 is
\optsshort{(0,0,2)}{(4,2,1)}{(4,3,1)}{(4,3,3)}
\end{question}

\begin{question}{H}
For the system with Allocation P0 (0,1,0), P1 (2,0,0), P2 (3,0,2), P3 (2,1,1), P4 (0,0,2), Max P0 (7,5,3), P1 (3,2,2), P2 (9,0,2), P3 (2,2,2), P4 (4,3,3) and Available (3,3,2), which of the following is a safe sequence?
\optsshort{P4, P1, P3, P0, P2}{P0, P1, P2, P3, P4}{P1, P3, P4, P2, P0}{P2, P1, P3, P4, P0}
\end{question}

\begin{question}{M}
In the same Banker's example (Available (3,3,2); after granting P1's request for (1,0,2) the state is still safe), if P0 then requests (0,2,0), the request
\optslong{causes immediate deadlock}{is rejected because it exceeds P0's maximum claim}{must wait because the resulting state would be unsafe}{is granted because resources are available and the state remains safe}
\end{question}

\begin{question}{H}
A system has 12 instances of one resource type. P0 has a maximum need of 10 and holds 5, P1 a maximum of 4 and holds 2, P2 a maximum of 9 and holds 2. Which statement is correct?
\optslong{The state is safe with the safe sequence P0, P1, P2}{The system is deadlocked}{The state is safe with the safe sequence P1, P0, P2}{The state is unsafe}
\end{question}

\begin{question}{M}
In the system with 12 instances where P0 (max 10, holds 5), P1 (max 4, holds 2) and P2 (max 9, holds 2), if P2 is given one more instance
\optslong{the state becomes unsafe}{the state remains safe}{P0 can finish first}{the system immediately deadlocks}
\end{question}

\begin{question}{H}
The time complexity of the safety algorithm of the Banker's algorithm for n processes and m resource types is
\optsshort{O(m $\times$ n)}{O(n\textsuperscript{3})}{O(m\textsuperscript{2} $\times$ n)}{O(m $\times$ n\textsuperscript{2})}
\end{question}

\begin{question}{M}
Given below are two statements.\\
Statement I: The Banker's algorithm requires each process to declare in advance the maximum number of instances of each resource type it may need.\\
Statement II: The Banker's algorithm can be applied only to resource types with a single instance.
\optslong{Both Statement I and Statement II are true}{Both Statement I and Statement II are false}{Statement I is true but Statement II is false}{Statement I is false but Statement II is true}
\end{question}

\begin{question}{M}
For deadlock detection when every resource type has a single instance, the OS maintains a
\optslong{wait-for graph}{safe sequence table}{resource-allocation graph with claim edges}{Need matrix}
\end{question}

\begin{question}{M}
Detecting a cycle in a wait-for graph of n vertices requires
\optsshort{O(n\textsuperscript{2}) operations}{O(n\textsuperscript{3}) operations}{O(n) operations}{O(log n) operations}
\end{question}

\begin{question}{H}
Match List I with List II.\\
List I (Strategy): A. Deadlock prevention  B. Deadlock avoidance  C. Deadlock detection  D. Ignoring the problem\\
List II (Method): I. Ostrich algorithm used by most general-purpose OSes  II. Banker's algorithm  III. Ensure at least one necessary condition cannot hold  IV. Allow deadlocks, then find and recover from them
\optsshort{A-III, B-IV, C-II, D-I}{A-II, B-III, C-IV, D-I}{A-III, B-II, C-I, D-IV}{A-III, B-II, C-IV, D-I}
\end{question}

\begin{question}{M}
Which of the following are issues that must be addressed when recovering from deadlock by resource preemption?\\
A. Selecting a victim  B. Rollback  C. Starvation  D. Increasing the number of resource types\\
Choose the correct answer from the options given below:
\optsshort{A, B, C and D}{A, B and C only}{A and B only}{B and D only}
\end{question}

\begin{question}{M}
Aborting all deadlocked processes to recover from deadlock
\optslong{is cheap because no computation is lost}{requires the Banker's algorithm}{breaks the deadlock with certainty but at great expense}{never breaks the deadlock}
\end{question}

\begin{question}{M}
A situation in which two processes keep changing their state in response to each other without making progress is called
\optsshort{livelock}{deadlock}{starvation}{thrashing}
\end{question}

\begin{question}{M}
Ordering resource types R1 < R2 < R3 and requiring every process to request resources in increasing order prevents deadlock because
\optslong{mutual exclusion is removed}{preemption is allowed}{processes never hold resources}{a circular wait becomes impossible}
\end{question}

\begin{question}{M}
Given below are two statements.\\
Statement I: Mutual exclusion cannot be denied for intrinsically non-shareable resources such as a mutex lock.\\
Statement II: Read-only files are shareable resources and are not involved in deadlocks.
\optslong{Both Statement I and Statement II are true}{Both Statement I and Statement II are false}{Statement I is true but Statement II is false}{Statement I is false but Statement II is true}
\end{question}

\begin{question}{M}
The wait-die and wound-wait schemes prevent deadlock using
\optslong{timestamps of processes or transactions}{the Banker's algorithm}{resource ordering}{wait-for graphs}
\end{question}

\begin{question}{M}
In the wait-die scheme, if an older process requests a resource held by a younger process, the older process
\optsshort{is terminated}{waits}{is rolled back}{preempts the resource}
\end{question}

\begin{question}{M}
In the wound-wait scheme, if an older process requests a resource held by a younger process
\optslong{the younger process is rolled back (wounded)}{the request is ignored}{the older process waits}{both processes are rolled back}
\end{question}

\begin{question}{H}
A system with Available = (1, 0, 0) for resource types A, B and C has three processes. Allocation: P0 (0,1,0), P1 (2,0,0), P2 (3,0,3). Request: P0 (0,0,0), P1 (2,0,2), P2 (0,0,0). Using the deadlock-detection algorithm, which statement is correct?
\optslong{P0 and P1 are deadlocked}{The system is deadlocked}{The system is not deadlocked; the sequence P0, P2, P1 completes}{Only P1 is deadlocked}
\end{question}

\begin{question}{M}
Deadlock avoidance is generally less restrictive than deadlock prevention because
\optslong{it never refuses a request}{it allows circular waits that already exist}{it denies a request only if granting it could lead to an unsafe state, rather than ruling out a necessary condition altogether}{it does not need any information in advance}
\end{question}

\begin{question}{M}
How often should the deadlock-detection algorithm be invoked?
\optslong{Only when the system is shut down}{Never, because detection is impossible}{After every instruction}{It depends on how often deadlocks occur and how many processes they affect}
\end{question}

\begin{question}{M}
Given below are two statements: one is labelled as Assertion A and the other is labelled as Reason R.\\
Assertion A: Most general-purpose operating systems, including Linux and Windows, ignore the deadlock problem.\\
Reason R: Deadlocks occur infrequently, and prevention, avoidance or detection would be costly.
\optslong{Both A and R are true and R is the correct explanation of A}{Both A and R are true but R is NOT the correct explanation of A}{A is true but R is false}{A is false but R is true}
\end{question}

\begin{question}{M}
In a system with 4 processes and 5 resources of a single type, where each process needs at most 2 resources, the system
\optslong{is deadlock-free}{may deadlock}{always deadlocks}{needs at least 8 resources to avoid deadlock}
\end{question}

\begin{question}{M}
Which deadlock condition is broken by allowing the OS to take a resource away from a process that holds it?
\optsshort{Mutual exclusion}{Circular wait}{No preemption}{Hold and wait}
\end{question}

\begin{question}{M}
Which of the following is a necessary condition for a cycle in a wait-for graph to occur?
\optslong{Each process waits for a resource held by the next process in the cycle}{All processes have the same priority}{All resources have multiple instances}{The system uses preemptive scheduling}
\end{question}

\begin{question}{M}
Given below are two statements.\\
Statement I: The resource-allocation-graph algorithm for deadlock avoidance uses claim edges.\\
Statement II: A request is granted only if converting the claim edge to an assignment edge does not create a cycle.
\optslong{Both Statement I and Statement II are true}{Both Statement I and Statement II are false}{Statement I is true but Statement II is false}{Statement I is false but Statement II is true}
\end{question}

\begin{answers}
\ans{1}{3}{These are the four Coffman conditions; starvation is not one of them.}
\ans{2}{1}{Deadlock-free if R $\geq$ n(k $-$ 1) + 1 = 3 $\times$ 1 + 1 = 4; with 4 instances some process can always finish.}
\ans{3}{3}{n(k $-$ 1) + 1 = 6 $\times$ 2 + 1 = 13.}
\ans{4}{3}{n(2 $-$ 1) + 1 $\leq$ 6 gives n $\leq$ 5.}
\ans{5}{1}{$\Sigma$(max $-$ 1) + 1 = (1 + 2 + 3 + 4 + 5) + 1 = 16.}
\ans{6}{3}{With several instances per type, a cycle is necessary but not sufficient for deadlock.}
\ans{7}{1}{With single instances, a cycle implies a deadlock.}
\ans{8}{4}{An edge from Rj to Pi is an assignment edge; claim edges (dashed) are used in avoidance.}
\ans{9}{1}{Each prevention method breaks one of the four conditions.}
\ans{10}{1}{These are the two drawbacks of denying hold and wait.}
\ans{11}{2}{Deadlocked states are a subset of unsafe states.}
\ans{12}{1}{Need[i][j] is how many more instances of Rj process Pi may still request.}
\ans{13}{3}{Need = Max $-$ Allocation = (4,3,3) $-$ (0,0,2) = (4,3,1).}
\ans{14}{3}{Needs are P0 (7,4,3), P1 (1,2,2), P2 (6,0,0), P3 (0,1,1), P4 (4,3,1). Only P1 or P3 can go first; P1, P3, P4, P2, P0 succeeds.}
\ans{15}{3}{Available would become (2,1,0), which is not enough for any process's remaining need, so the state would be unsafe.}
\ans{16}{3}{3 are free; P1 can finish (5 free), then P0 (10 free), then P2.}
\ans{17}{1}{Only 2 would remain free: P1 can finish (4 free), but then neither P0 (needs 5) nor P2 (needs 6) can be guaranteed.}
\ans{18}{4}{In the worst case each of n passes scans n processes and compares m resources.}
\ans{19}{3}{The Banker's algorithm handles multiple instances; the resource-allocation-graph algorithm is the one limited to single instances.}
\ans{20}{1}{A wait-for graph is obtained by collapsing the resource nodes; a cycle means deadlock.}
\ans{21}{1}{Cycle detection in a graph with up to n\textsuperscript{2} edges takes O(n\textsuperscript{2}).}
\ans{22}{4}{The four approaches to handling deadlocks.}
\ans{23}{2}{The same victim must not always be picked, or it will starve.}
\ans{24}{3}{Aborting one process at a time is cheaper per abort but requires repeated detection.}
\ans{25}{1}{In a livelock the processes are not blocked, but they still make no progress.}
\ans{26}{4}{A cycle would require some process to request a lower-numbered resource while holding a higher one.}
\ans{27}{1}{Shareable resources do not require mutually exclusive access.}
\ans{28}{1}{In wait-die an older process waits and a younger one dies; in wound-wait an older process preempts a younger one.}
\ans{29}{2}{Older processes may wait for younger ones; younger processes requesting resources held by older ones are rolled back (die).}
\ans{30}{1}{A younger process requesting a resource held by an older one waits.}
\ans{31}{3}{P0 finishes: work (1,1,0); P2 finishes: (4,1,3); P1's request (2,0,2) can then be met.}
\ans{32}{3}{Prevention constrains how requests may be made; avoidance decides dynamically.}
\ans{33}{4}{It may run on every request that cannot be granted, periodically, or when CPU utilisation drops below a threshold.}
\ans{34}{1}{Applications are expected to handle deadlocks themselves.}
\ans{35}{1}{4 $\times$ (2 $-$ 1) + 1 = 5, so at least one process can always obtain both resources.}
\ans{36}{3}{This works for resources whose state can be saved, such as CPU registers and memory.}
\ans{37}{1}{An edge Pi $\rightarrow$ Pj means Pi waits for a resource that Pj holds.}
\ans{38}{1}{This works for single-instance resource types.}
\end{answers}

\section{Memory Management (address binding, contiguous allocation, paging, TLB, segmentation)}

\begin{question}{H}
Consider a logical address space of 8 pages of 1024 words each, mapped onto a physical memory of 32 frames. The number of bits in the logical address and in the physical address are respectively
\optsshort{15 and 13}{13 and 10}{13 and 15}{10 and 15}
\end{question}

\begin{question}{H}
A system uses 32-bit logical addresses, a page size of 4 KB and page-table entries of 4 bytes. The size of a single-level page table for one process is
\optsshort{4 MB}{4 KB}{1 MB}{2 MB}
\end{question}

\begin{question}{H}
A computer has 64 MB of physical memory and a frame size of 4 KB. The number of bits needed for the frame number is
\optsshort{12}{26}{14}{16}
\end{question}

\begin{question}{H}
With a page size of 1024 bytes, the logical address 10000 (decimal) corresponds to
\optsshort{page 8, offset 784}{page 9, offset 1000}{page 10, offset 0}{page 9, offset 784}
\end{question}

\begin{question}{M}
With a 4 KB page size, the 16-bit virtual address 0x3A7F has page number and offset
\optsshort{0x3A and 0x7F}{0xA and 0x37F}{7 and 0xA3F}{3 and 0xA7F}
\end{question}

\begin{question}{H}
In a paging system with a TLB, the TLB access time is 20 ns, the memory access time is 100 ns and the TLB hit ratio is 80\%. The effective memory access time is
\optsshort{120 ns}{220 ns}{140 ns}{200 ns}
\end{question}

\begin{question}{H}
In a paging system the TLB lookup takes 20 ns and a memory access 100 ns. If the hit ratio is 98\%, the effective access time is
\optsshort{122 ns}{220 ns}{140 ns}{120 ns}
\end{question}

\begin{question}{H}
The TLB lookup takes 20 ns and a memory access 100 ns. The minimum TLB hit ratio needed to keep the effective memory access time at 130 ns is
\optsshort{95\%}{90\%}{80\%}{85\%}
\end{question}

\begin{question}{M}
Without a TLB, a paging system with a single-level page table kept in memory, where a memory access takes 100 ns, has an effective access time of
\optsshort{200 ns}{100 ns}{300 ns}{150 ns}
\end{question}

\begin{question}{H}
A TLB has 64 entries and the page size is 4 KB. The TLB reach is
\optsshort{64 KB}{4 MB}{16 MB}{256 KB}
\end{question}

\begin{question}{H}
Given below are two statements.\\
Statement I: Paging eliminates internal fragmentation but can suffer from external fragmentation.\\
Statement II: Segmentation eliminates internal fragmentation but can suffer from external fragmentation.
\optslong{Both Statement I and Statement II are true}{Both Statement I and Statement II are false}{Statement I is true but Statement II is false}{Statement I is false but Statement II is true}
\end{question}

\begin{question}{H}
Match List I with List II.\\
List I (Address binding time): A. Compile time  B. Load time  C. Execution time\\
List II (Result): I. Requires hardware support such as a relocation register (MMU)  II. Absolute code is generated  III. Relocatable code is generated
\optsshort{A-II, B-III, C-I}{A-III, B-II, C-I}{A-I, B-III, C-II}{A-II, B-I, C-III}
\end{question}

\begin{question}{H}
With a relocation register value of 14000, the logical address 346 is mapped to the physical address
\optsshort{346}{14346}{13654}{14000}
\end{question}

\begin{question}{H}
Memory partitions of 100 KB, 500 KB, 200 KB, 300 KB and 600 KB (in order) must hold processes of 212 KB, 417 KB, 112 KB and 426 KB (in order). Which placement algorithm allocates all four processes?
\optslong{First fit}{Next fit starting at the 600 KB partition}{Best fit}{Worst fit}
\end{question}

\begin{question}{H}
Memory partitions of 100 KB, 500 KB, 200 KB, 300 KB and 600 KB (in order) must hold processes of 212 KB, 417 KB, 112 KB and 426 KB (in order). Using first fit, the 112 KB process is placed in
\optslong{the 183 KB remainder of the 600 KB partition}{the 200 KB partition}{the 288 KB remainder of the 500 KB partition}{the 300 KB partition}
\end{question}

\begin{question}{M}
Given below are two statements.\\
Statement I: Best fit searches the entire list unless the list is ordered by size, and produces the smallest leftover hole.\\
Statement II: Worst fit allocates the largest hole, producing the largest leftover hole.
\optslong{Both Statement I and Statement II are true}{Both Statement I and Statement II are false}{Statement I is true but Statement II is false}{Statement I is false but Statement II is true}
\end{question}

\begin{question}{M}
According to the 50-percent rule for first-fit allocation, if N blocks are allocated, about 0.5N blocks are lost to fragmentation. The fraction of memory that may be unusable is about
\optsshort{one-third}{one-quarter}{one-half}{two-thirds}
\end{question}

\begin{question}{M}
Compaction as a solution to external fragmentation is possible only if
\optslong{relocation is dynamic and done at execution time}{relocation is static and done at assembly time}{paging is used}{segments are of equal size}
\end{question}

\begin{question}{H}
A segment table has entries (base, limit): segment 0 (219, 600), segment 1 (2300, 14), segment 2 (90, 100), segment 3 (1327, 580), segment 4 (1952, 96). The physical address for the logical address (0, 430) is
\optsshort{649}{819}{430}{2730}
\end{question}

\begin{question}{H}
For the segment table with entries (base, limit): segment 0 (219, 600), segment 1 (2300, 14), segment 2 (90, 100), segment 3 (1327, 580), segment 4 (1952, 96), which logical address causes an addressing-error trap?
\optsshort{(0, 430)}{(2, 500)}{(1, 10)}{(3, 400)}
\end{question}

\begin{question}{H}
Match List I with List II.\\
List I (Page-table structure): A. Hierarchical page table  B. Hashed page table  C. Inverted page table\\
List II (Characteristic): I. One entry per physical frame of memory  II. The page table itself is paged, e.g. with an outer and an inner table  III. Virtual page numbers are hashed into a table with linked chains, common for address spaces larger than 32 bits
\optsshort{A-III, B-II, C-I}{A-I, B-III, C-II}{A-II, B-III, C-I}{A-II, B-I, C-III}
\end{question}

\begin{question}{H}
Given below are two statements: one is labelled as Assertion A and the other is labelled as Reason R.\\
Assertion A: An inverted page table reduces the memory needed to store page tables.\\
Reason R: An inverted page table makes it easy to share memory between processes.
\optslong{Both A and R are true and R is the correct explanation of A}{Both A and R are true but R is NOT the correct explanation of A}{A is true but R is false}{A is false but R is true}
\end{question}

\begin{question}{M}
A system with 1 GB of physical memory and a 4 KB page size uses an inverted page table. The number of entries in the table is
\optsshort{2\textsuperscript{30}}{2\textsuperscript{20}}{2\textsuperscript{12}}{2\textsuperscript{18}}
\end{question}

\begin{question}{H}
In a two-level paging scheme for a 32-bit logical address with 4 KB pages and 4-byte page-table entries, chosen so that each inner page table fits in one page, the address is divided as
\optslong{12 bits outer page, 8 bits inner page, 12 bits offset}{16 bits outer page, 4 bits inner page, 12 bits offset}{8 bits outer page, 12 bits inner page, 12 bits offset}{10 bits outer page, 10 bits inner page, 12 bits offset}
\end{question}

\begin{question}{M}
A system has 46-bit virtual addresses, 8 KB pages and 4-byte page-table entries. How many levels of paging are needed if every page table must fit in a single page?
\optsshort{2}{3}{5}{4}
\end{question}

\begin{question}{M}
The optimal page size p that minimises page-table overhead plus internal fragmentation, for an average process size s and a page-table-entry size e, is p = sqrt(2se). For s = 1 MB and e = 8 bytes, p is
\optsshort{4 KB}{1 KB}{8 KB}{2 KB}
\end{question}

\begin{question}{M}
On average, the internal fragmentation per process in a paging system is about
\optsshort{one frame per segment}{one page}{zero}{half a page}
\end{question}

\begin{question}{M}
The valid-invalid bit in a page-table entry indicates
\optslong{whether the page is read-only}{whether the page is shared}{whether the page has been modified}{whether the page is in the process's logical address space (and, with demand paging, in memory)}
\end{question}

\begin{question}{M}
An address space identifier (ASID) stored in each TLB entry
\optslong{identifies the process the entry belongs to, so the TLB need not be flushed on every context switch}{stores the frame number}{marks the entry as dirty}{stores the page offset}
\end{question}

\begin{question}{M}
Given below are two statements.\\
Statement I: Shared pages allow several processes to share the same reentrant code, such as a library.\\
Statement II: To share code, the code must not modify itself during execution.
\optslong{Both Statement I and Statement II are true}{Both Statement I and Statement II are false}{Statement I is true but Statement II is false}{Statement I is false but Statement II is true}
\end{question}

\begin{question}{M}
In the Intel IA-32 architecture, a logical address is translated by the segmentation unit into a
\optslong{virtual page number}{segment selector}{linear address, which the paging unit translates into a physical address}{physical address directly}
\end{question}

\begin{question}{H}
Using the buddy system with a 256 KB memory segment, a request for 21 KB is satisfied with a block of size
\optsshort{21 KB}{64 KB}{24 KB}{32 KB}
\end{question}

\begin{question}{M}
In the buddy system, a 32 KB block is allocated for a 21 KB request. The internal fragmentation is
\optsshort{11 KB}{0 KB}{5 KB}{21 KB}
\end{question}

\begin{question}{M}
The slab allocator in the Linux kernel
\optslong{allocates memory in power-of-two blocks only}{uses caches of pre-allocated objects of each kernel data-structure type, avoiding fragmentation}{is used only for user processes}{is a page-replacement algorithm}
\end{question}

\begin{question}{M}
Swapping a process out of memory to a backing store and later back in is used to
\optslong{speed up I/O}{reduce internal fragmentation}{increase the size of the TLB}{allow the total memory of all processes to exceed physical memory}
\end{question}

\begin{question}{M}
In a system with base and limit registers, an address generated by a user process is legal if
\optslong{address $\geq$ base and address < base + limit}{address < base}{address equals base}{address > limit}
\end{question}

\begin{question}{M}
Match List I with List II.\\
List I (Term): A. Internal fragmentation  B. External fragmentation  C. Compaction  D. Paging\\
List II (Description): I. Memory is free in total but not contiguous  II. Shuffling memory contents to place free memory in one block  III. Allocated memory slightly larger than requested; the difference is unused  IV. Non-contiguous allocation of fixed-size frames
\optsshort{A-I, B-III, C-II, D-IV}{A-III, B-II, C-I, D-IV}{A-III, B-I, C-II, D-IV}{A-III, B-I, C-IV, D-II}
\end{question}

\begin{question}{M}
The page-table base register (PTBR)
\optslong{stores the frame number of the current page}{counts page faults}{stores the size of the TLB}{points to the page table of the currently running process}
\end{question}

\begin{question}{M}
Given below are two statements.\\
Statement I: A larger page size gives a smaller page table.\\
Statement II: A larger page size reduces internal fragmentation.
\optslong{Both Statement I and Statement II are true}{Both Statement I and Statement II are false}{Statement I is true but Statement II is false}{Statement I is false but Statement II is true}
\end{question}

\begin{question}{M}
A paging system uses a logical address space of 2\textsuperscript{16} bytes with a page size of 512 bytes. The number of pages and the number of offset bits are
\optsshort{64 and 10}{512 and 7}{256 and 8}{128 and 9}
\end{question}

\begin{question}{M}
Which part of the logical address is used to index the page table?
\optsshort{Frame number}{Offset}{Segment limit}{Page number}
\end{question}

\begin{question}{M}
In segmentation, a logical address consists of
\optslong{a frame number and a page number}{a segment number and an offset within the segment}{a page number and an offset}{a base and a limit}
\end{question}

\begin{question}{M}
Given below are two statements: one is labelled as Assertion A and the other is labelled as Reason R.\\
Assertion A: Segmentation supports the programmer's view of memory.\\
Reason R: A program is viewed as a collection of logical units such as the main program, functions, stack and arrays, each of a different size.
\optslong{Both A and R are true and R is the correct explanation of A}{Both A and R are true but R is NOT the correct explanation of A}{A is true but R is false}{A is false but R is true}
\end{question}

\begin{question}{M}
Which of the following problems is NOT solved by paging?
\optslong{The need for contiguous physical memory for a process}{Fitting variable-size chunks into the backing store}{External fragmentation}{Internal fragmentation}
\end{question}

\begin{question}{M}
In a paged system, the page-table length register (PTLR) is used to
\optslong{count the number of frames}{store the size of the TLB}{store the length of each page}{check that a page number is within the valid range for the process}
\end{question}

\begin{question}{M}
A system with 48-bit virtual addresses and 4 KB pages would need how many entries in a single-level page table?
\optsshort{2\textsuperscript{24}}{2\textsuperscript{48}}{2\textsuperscript{36}}{2\textsuperscript{12}}
\end{question}

\begin{question}{M}
In a paging system, physical memory is divided into fixed-size blocks called
\optsshort{segments}{frames}{partitions}{pages}
\end{question}

\begin{question}{L}
Which memory allocation technique places a new process in the first free hole after the point where the previous search ended?
\optsshort{Best fit}{Worst fit}{Next fit}{First fit}
\end{question}

\begin{question}{M}
In a system using paging, each memory reference to a page table entry that is not cached requires
\optslong{one extra memory access per level of the page table}{a TLB flush}{no memory access}{a disk access}
\end{question}

\begin{answers}
\ans{1}{3}{Logical: 8 $\times$ 1024 = 2\textsuperscript{13} words; physical: 32 $\times$ 1024 = 2\textsuperscript{15} words.}
\ans{2}{1}{2\textsuperscript{32} / 2\textsuperscript{12} = 2\textsuperscript{20} entries $\times$ 4 bytes = 4 MB.}
\ans{3}{3}{64 MB / 4 KB = 2\textsuperscript{26} / 2\textsuperscript{12} = 2\textsuperscript{14} frames.}
\ans{4}{4}{10000 / 1024 = 9 remainder 784.}
\ans{5}{4}{The low 12 bits (0xA7F) are the offset; the remaining bits (0x3) give the page number.}
\ans{6}{3}{EAT = 0.8 $\times$ (20 + 100) + 0.2 $\times$ (20 + 100 + 100) = 96 + 44 = 140 ns.}
\ans{7}{1}{EAT = 0.98 $\times$ 120 + 0.02 $\times$ 220 = 117.6 + 4.4 = 122 ns.}
\ans{8}{2}{130 = h $\times$ 120 + (1 $-$ h) $\times$ 220 gives 100h = 90, so h = 0.9.}
\ans{9}{1}{Each reference needs one access for the page-table entry and one for the data.}
\ans{10}{4}{TLB reach = number of entries $\times$ page size = 64 $\times$ 4 KB.}
\ans{11}{4}{Paging is the reverse: no external fragmentation, but internal fragmentation in the last page.}
\ans{12}{1}{Execution-time binding lets a process move in memory during execution.}
\ans{13}{2}{Physical address = relocation (base) value + logical address.}
\ans{14}{3}{Best fit uses 300, 500, 200 and 600 KB; first fit and worst fit leave the 426 KB process waiting.}
\ans{15}{3}{212 KB goes into 500 (leaving 288), 417 KB into 600, and 112 KB into the first hole large enough, the 288 KB remainder.}
\ans{16}{1}{Simulations show first fit and best fit are better than worst fit in time and storage use.}
\ans{17}{1}{0.5N out of 1.5N blocks is one-third.}
\ans{18}{1}{Programs and data must be movable, so addresses must be bound at run time.}
\ans{19}{1}{The offset 430 is less than the limit 600, so the address is 219 + 430 = 649.}
\ans{20}{2}{Segment 2 has a limit of 100, so offset 500 is illegal; (1, 10) $\rightarrow$ 2310 and (3, 400) $\rightarrow$ 1727 are valid.}
\ans{21}{3}{Standard page-table structures.}
\ans{22}{3}{With one entry per frame, a frame can map only one virtual page, which makes sharing difficult.}
\ans{23}{4}{2\textsuperscript{30} / 2\textsuperscript{12} = 2\textsuperscript{18} entries, one per frame.}
\ans{24}{4}{A 4 KB page holds 1024 four-byte entries, so each level uses 10 bits.}
\ans{25}{2}{The offset is 13 bits; a page holds 2048 = 2\textsuperscript{11} entries; (46 $-$ 13) / 11 = 3 levels.}
\ans{26}{1}{sqrt(2 $\times$ 2\textsuperscript{20} $\times$ 8) = sqrt(2\textsuperscript{24}) = 2\textsuperscript{12} = 4096 bytes.}
\ans{27}{4}{Only the last page of a process is partly used, on average half-full.}
\ans{28}{4}{Accessing an invalid page causes a trap to the operating system.}
\ans{29}{1}{Without ASIDs, the TLB must be flushed when the address space changes.}
\ans{30}{1}{Each process has its own data pages while mapping the same code frames.}
\ans{31}{3}{IA-32 supports segmentation with paging.}
\ans{32}{4}{Blocks are powers of two: 256 $\rightarrow$ 128 $\rightarrow$ 64 $\rightarrow$ 32; the smallest power of two $\geq$ 21 KB is 32 KB.}
\ans{33}{1}{32 $-$ 21 = 11 KB is wasted inside the block.}
\ans{34}{2}{Each cache consists of slabs of contiguous pages holding objects of one type.}
\ans{35}{4}{Modern systems swap pages rather than entire processes.}
\ans{36}{1}{Otherwise the hardware traps to the operating system.}
\ans{37}{3}{Standard memory-management terms.}
\ans{38}{4}{Changing page tables on a context switch requires changing only this register.}
\ans{39}{3}{A larger page wastes more space, on average half a page per process.}
\ans{40}{4}{2\textsuperscript{16} / 2\textsuperscript{9} = 2\textsuperscript{7} = 128 pages; offset = log2 512 = 9 bits.}
\ans{41}{4}{The page number selects the entry; the offset is copied unchanged into the physical address.}
\ans{42}{2}{The segment number indexes the segment table, which supplies the base and limit.}
\ans{43}{1}{Segmentation maps these variable-size logical units directly.}
\ans{44}{4}{Paging still wastes part of the last page of a process.}
\ans{45}{4}{Page numbers beyond the PTLR value cause a trap.}
\ans{46}{3}{2\textsuperscript{48} / 2\textsuperscript{12} = 2\textsuperscript{36} entries, which is why multilevel or hashed tables are used.}
\ans{47}{2}{Logical memory is divided into pages of the same size.}
\ans{48}{3}{Next fit is a variation of first fit that starts each search where the last one ended.}
\ans{49}{1}{A two-level table needs two extra accesses on a TLB miss.}
\end{answers}

\section{Virtual Memory (demand paging, page replacement, frame allocation, thrashing, working sets)}

\begin{question}{H}
For the reference string 7, 0, 1, 2, 0, 3, 0, 4, 2, 3, 0, 3, 2, 1, 2, 0, 1, 7, 0, 1 with 3 frames (initially empty), the number of page faults under FIFO replacement is
\optsshort{16}{9}{15}{12}
\end{question}

\begin{question}{H}
For the reference string 7, 0, 1, 2, 0, 3, 0, 4, 2, 3, 0, 3, 2, 1, 2, 0, 1, 7, 0, 1 with 3 frames (initially empty), the number of page faults under LRU replacement is
\optsshort{15}{9}{12}{10}
\end{question}

\begin{question}{H}
For the reference string 7, 0, 1, 2, 0, 3, 0, 4, 2, 3, 0, 3, 2, 1, 2, 0, 1, 7, 0, 1 with 3 frames (initially empty), the number of page faults under optimal replacement is
\optsshort{9}{8}{6}{12}
\end{question}

\begin{question}{M}
For the reference string 7, 0, 1, 2, 0, 3, 0, 4, 2, 3, 0, 3, 2, 1, 2, 0, 1, 7, 0, 1 with 4 frames (initially empty), the numbers of page faults under FIFO and LRU are respectively
\optsshort{10 and 8}{10 and 10}{12 and 8}{9 and 8}
\end{question}

\begin{question}{H}
For the reference string 1, 2, 3, 4, 1, 2, 5, 1, 2, 3, 4, 5 under FIFO replacement, the numbers of page faults with 3 frames and with 4 frames are respectively
\optsshort{9 and 8}{10 and 9}{9 and 10}{10 and 8}
\end{question}

\begin{question}{H}
Belady's anomaly, in which increasing the number of frames can increase the number of page faults, can occur with
\optsshort{optimal replacement}{FIFO}{any stack algorithm}{LRU}
\end{question}

\begin{question}{H}
Given below are two statements: one is labelled as Assertion A and the other is labelled as Reason R.\\
Assertion A: FIFO page replacement never suffers from Belady's anomaly.\\
Reason R: LRU is a stack algorithm: the set of pages in memory with n frames is always a subset of the set with n + 1 frames.
\optslong{Both A and R are true and R is the correct explanation of A}{Both A and R are true but R is NOT the correct explanation of A}{A is true but R is false}{A is false but R is true}
\end{question}

\begin{question}{H}
For the reference string 1, 2, 3, 4, 2, 1, 5, 6, 2, 1, 2, 3, 7, 6, 3, 2, 1, 2, 3, 6 with 4 frames (initially empty), the numbers of page faults under FIFO, LRU and optimal replacement are respectively
\optsshort{10, 10 and 8}{16, 15 and 11}{14, 8 and 10}{14, 10 and 8}
\end{question}

\begin{question}{M}
For the reference string 1, 2, 3, 4, 2, 1, 5, 6, 2, 1, 2, 3, 7, 6, 3, 2, 1, 2, 3, 6 with 3 frames (initially empty), the number of page faults under LRU is
\optsshort{16}{14}{11}{15}
\end{question}

\begin{question}{M}
For the reference string 4, 7, 6, 1, 7, 6, 1, 2, 7, 2 with 3 frames (initially empty), the number of page faults under optimal replacement is
\optsshort{7}{5}{4}{6}
\end{question}

\begin{question}{M}
For the reference string 0, 1, 2, 3, 0, 1, 4, 0, 1, 2, 3, 4 under FIFO replacement, the numbers of page faults with 3 and 4 frames are respectively
\optsshort{9 and 10}{8 and 10}{10 and 9}{9 and 9}
\end{question}

\begin{question}{H}
The page replacement algorithm that has the lowest page-fault rate but cannot be implemented in practice is
\optsshort{Optimal (OPT/MIN)}{FIFO}{LRU}{Second chance}
\end{question}

\begin{question}{H}
Given below are two statements.\\
Statement I: The second-chance (clock) algorithm is a FIFO algorithm that skips a page whose reference bit is 1, clearing the bit.\\
Statement II: If all reference bits are set, second-chance degenerates into FIFO.
\optslong{Both Statement I and Statement II are true}{Both Statement I and Statement II are false}{Statement I is true but Statement II is false}{Statement I is false but Statement II is true}
\end{question}

\begin{question}{M}
In the enhanced second-chance algorithm that uses the (reference bit, modify bit) pair, the best page to replace is in the class
\optsshort{(1, 1)}{(0, 0)}{(1, 0)}{(0, 1)}
\end{question}

\begin{question}{H}
In a demand-paging system, the memory access time is 200 ns and the average page-fault service time is 8 ms. If the page-fault rate is 1 in 1000, the effective access time is approximately
\optsshort{8 ms}{8.2 $\mu$s}{200 ns}{1 $\mu$s}
\end{question}

\begin{question}{H}
With a memory access time of 200 ns and a page-fault service time of 8 ms, to keep the slowdown due to paging below 10\%, the page-fault rate must be less than about
\optsshort{1 in 4,000,000}{1 in 400,000}{1 in 1000}{1 in 40,000}
\end{question}

\begin{question}{M}
A page fault takes 8 ms if a free frame is available or the replaced page is clean, and 20 ms if the replaced page is modified; 70\% of replaced pages are modified. The memory access time is 100 ns. The maximum page-fault rate for an effective access time of 200 ns is approximately
\optsshort{6.1 $\times$ 10\textsuperscript{-6}}{5 $\times$ 10\textsuperscript{-6}}{1.25 $\times$ 10\textsuperscript{-5}}{1 $\times$ 10\textsuperscript{-4}}
\end{question}

\begin{question}{H}
Arrange the steps in handling a page fault in the correct order.\\
A. Schedule a disk read of the page into a free frame  B. Trap to the operating system  C. Restart the instruction that caused the fault  D. Check that the reference was a valid memory access  E. Update the page table
\optsshort{B, A, D, E, C}{B, D, E, A, C}{D, B, A, E, C}{B, D, A, E, C}
\end{question}

\begin{question}{H}
Thrashing occurs when
\optslong{the page size is too large}{the CPU is idle because of I/O}{a process spends more time paging than executing}{the TLB hit ratio is 100\%}
\end{question}

\begin{question}{H}
Given below are two statements.\\
Statement I: Using a local replacement algorithm completely eliminates thrashing.\\
Statement II: If the sum of the working-set sizes of all processes exceeds the number of available frames, thrashing will occur.
\optslong{Both Statement I and Statement II are true}{Both Statement I and Statement II are false}{Statement I is true but Statement II is false}{Statement I is false but Statement II is true}
\end{question}

\begin{question}{H}
Using the working-set model with window $\Delta$ = 10, the working set at time t for the most recent references 2, 6, 1, 5, 7, 7, 7, 7, 5, 1 is
\optsshort{\{1, 5, 7\}}{\{7\}}{\{2, 6, 1, 5, 7, 7\}}{\{1, 2, 5, 6, 7\}}
\end{question}

\begin{question}{M}
In the page-fault frequency (PFF) scheme, if the page-fault rate of a process rises above an upper bound, the OS
\optslong{allocates another frame to the process}{takes a frame away from the process}{reduces the page size}{terminates the process}
\end{question}

\begin{question}{H}
Using proportional allocation, 62 free frames are divided between two processes of 10 pages and 127 pages. The numbers of frames allocated are approximately
\optsshort{4 and 57}{10 and 52}{31 and 31}{5 and 57}
\end{question}

\begin{question}{M}
Given below are two statements.\\
Statement I: The minimum number of frames per process is defined by the computer architecture.\\
Statement II: The maximum number of frames per process is defined by the amount of available physical memory.
\optslong{Both Statement I and Statement II are true}{Both Statement I and Statement II are false}{Statement I is true but Statement II is false}{Statement I is false but Statement II is true}
\end{question}

\begin{question}{M}
In global page replacement
\optslong{a process can select a replacement frame from the set of all frames, even one belonging to another process}{a process selects a replacement frame only from its own set of frames}{no page is ever replaced}{only kernel pages are replaced}
\end{question}

\begin{question}{H}
A program accesses every element of the array int data[128][128], stored in row-major order, with pages of 128 words (one row per page) and fewer frames than rows. If the elements are accessed column by column, the number of page faults is
\optsshort{16,384}{1}{128}{256}
\end{question}

\begin{question}{M}
Copy-on-write is used with fork() so that
\optslong{the child gets a copy of all pages immediately}{the parent and child initially share pages, and a page is copied only when either process writes to it}{the child cannot modify any page}{pages are written to disk on creation}
\end{question}

\begin{question}{M}
In demand paging, a page is brought into memory
\optslong{when the process terminates}{only when it is needed (referenced)}{before the process starts}{at fixed intervals}
\end{question}

\begin{question}{M}
Pure demand paging means
\optslong{the page size is fixed at 4 KB}{pages are never replaced}{a process starts with no pages in memory, and even the first instruction causes a page fault}{all pages are loaded at the start}
\end{question}

\begin{question}{M}
The dirty (modify) bit associated with a page is used to
\optslong{avoid writing a page back to disk if it has not been modified}{indicate the page is shared}{indicate that the page has been referenced recently}{mark invalid pages}
\end{question}

\begin{question}{H}
Match List I with List II.\\
List I (Page replacement algorithm): A. FIFO  B. Optimal  C. LRU  D. LFU\\
List II (Page replaced): I. The page used least recently  II. The page used the fewest number of times  III. The oldest page in memory  IV. The page that will not be used for the longest time
\optsshort{A-III, B-I, C-IV, D-II}{A-III, B-IV, C-I, D-II}{A-III, B-IV, C-II, D-I}{A-IV, B-III, C-I, D-II}
\end{question}

\begin{question}{M}
LRU can be implemented exactly using
\optslong{the dirty bit only}{a single reference bit}{a FIFO queue}{counters (time-of-use fields) or a stack of page numbers}
\end{question}

\begin{question}{M}
Locality of reference refers to the tendency of a program to
\optslong{access a relatively small set of pages over any short period}{access pages in strictly sequential order}{use only one page during its execution}{access memory addresses randomly}
\end{question}

\begin{question}{M}
Memory-mapped files allow
\optslong{the CPU to read the disk directly}{processes to avoid page faults}{files to be stored in the TLB}{file I/O to be treated as routine memory accesses by mapping disk blocks to pages in memory}
\end{question}

\begin{question}{M}
Which of the following reduces the number of page faults at process start-up by bringing in several pages at once?
\optsshort{Prepaging}{Compaction}{Copy-on-write}{Page buffering}
\end{question}

\begin{question}{M}
Locking a page in memory (an I/O interlock) is needed when
\optslong{the page is not in the working set}{the page holds read-only code}{the page is dirty}{an I/O device is transferring data to or from a buffer in that page}
\end{question}

\begin{question}{M}
Given below are two statements.\\
Statement I: vfork() does not use copy-on-write; the parent is suspended and the child uses the parent's address space.\\
Statement II: vfork() is intended for use when the child calls exec() immediately after creation.
\optslong{Both Statement I and Statement II are true}{Both Statement I and Statement II are false}{Statement I is true but Statement II is false}{Statement I is false but Statement II is true}
\end{question}

\begin{question}{M}
When a page fault occurs and there is no free frame, the operating system
\optslong{uses a page-replacement algorithm to select a victim frame}{increases the size of memory}{stops all other processes}{terminates the process}
\end{question}

\begin{question}{M}
If the degree of multiprogramming is increased beyond the point where thrashing begins, CPU utilisation
\optsshort{keeps increasing}{stays constant}{drops sharply}{becomes 100\%}
\end{question}

\begin{question}{H}
Given below are two statements: one is labelled as Assertion A and the other is labelled as Reason R.\\
Assertion A: Virtual memory allows the logical address space of a process to be larger than physical memory.\\
Reason R: The TLB caches recently used page-table entries.
\optslong{Both A and R are true and R is the correct explanation of A}{Both A and R are true but R is NOT the correct explanation of A}{A is true but R is false}{A is false but R is true}
\end{question}

\begin{question}{M}
A process references 5 distinct pages and has 3 frames. Under any replacement algorithm, the minimum number of page faults is
\optsshort{15}{5}{3}{8}
\end{question}

\begin{question}{M}
Which of the following algorithms approximates LRU using a reference bit and a circular list of pages?
\optsshort{LFU}{MFU}{Clock (second chance)}{FIFO}
\end{question}

\begin{question}{M}
Given below are two statements.\\
Statement I: Increasing the page size generally reduces the number of page faults for programs with good locality.\\
Statement II: A smaller page size allows better resolution, bringing in only the memory that is actually needed.
\optslong{Both Statement I and Statement II are true}{Both Statement I and Statement II are false}{Statement I is true but Statement II is false}{Statement I is false but Statement II is true}
\end{question}

\begin{question}{M}
Which of the following is a counting-based page replacement algorithm?
\optsshort{OPT}{LFU}{LRU}{FIFO}
\end{question}

\begin{question}{L}
In Windows and Linux, the area of a disk used to hold pages that have been swapped out of memory is called
\optsshort{swap space}{the page cache}{the master file table}{the boot block}
\end{question}

\begin{question}{M}
Which of the following causes a page fault?
\optslong{Accessing a page whose valid-invalid bit is invalid}{Executing a privileged instruction in kernel mode}{A TLB hit}{Accessing a page that is in memory}
\end{question}

\begin{question}{M}
The working set of a process is defined as
\optslong{the pages that are dirty}{the set of pages referenced in the most recent $\Delta$ references}{the pages in the TLB}{all pages of the process}
\end{question}

\begin{answers}
\ans{1}{3}{FIFO gives 15 faults for this standard reference string.}
\ans{2}{3}{LRU gives 12 faults.}
\ans{3}{1}{The optimal algorithm replaces the page used furthest in the future and gives 9 faults.}
\ans{4}{1}{With 4 frames FIFO gives 10 faults and LRU gives 8.}
\ans{5}{3}{More frames give more faults: this is Belady's anomaly.}
\ans{6}{2}{LRU and OPT are stack algorithms and never exhibit the anomaly.}
\ans{7}{4}{FIFO is the classic algorithm that does show Belady's anomaly; LRU does not, because it is a stack algorithm.}
\ans{8}{4}{Simulation gives FIFO 14, LRU 10 and OPT 8.}
\ans{9}{4}{Simulation gives LRU 15 (FIFO 16, OPT 11).}
\ans{10}{2}{4, 7, 6 fault; 1 replaces 4; 2 replaces 1 (not used again); total 5.}
\ans{11}{1}{This string, too, shows Belady's anomaly.}
\ans{12}{1}{It needs knowledge of the future reference string; it is used as a benchmark.}
\ans{13}{1}{The hand clears every bit in one sweep and then replaces the page it started from.}
\ans{14}{2}{(0, 0): neither recently used nor modified, so it need not be written back.}
\ans{15}{2}{EAT = 0.999 $\times$ 200 + 0.001 $\times$ 8,000,000 $\approx$ 8199.8 ns $\approx$ 8.2 $\mu$s, a 40-fold slowdown.}
\ans{16}{2}{220 > 200 + 7,999,800p gives p < 2.5 $\times$ 10\textsuperscript{-6}, i.e. fewer than 1 fault in 400,000 accesses.}
\ans{17}{1}{Average service time = 0.3 $\times$ 8 + 0.7 $\times$ 20 = 16.4 ms; 200 = (1 $-$ p) 100 + p $\times$ 16.4 $\times$ 10\textsuperscript{6} gives p $\approx$ 6.1 $\times$ 10\textsuperscript{-6}.}
\ans{18}{4}{Trap, check validity, read the page, update the tables, restart the instruction.}
\ans{19}{3}{As the degree of multiprogramming increases beyond a point, CPU utilisation drops sharply.}
\ans{20}{4}{Local replacement only stops one thrashing process from stealing frames from others; it does not solve the problem.}
\ans{21}{4}{The working set is the set of distinct pages in the last $\Delta$ references.}
\ans{22}{1}{If the rate falls below a lower bound, a frame is taken away.}
\ans{23}{1}{10/137 $\times$ 62 $\approx$ 4 and 127/137 $\times$ 62 $\approx$ 57.}
\ans{24}{1}{An instruction may need several pages at once, which sets the minimum.}
\ans{25}{1}{Global replacement gives better throughput but a process cannot control its own fault rate.}
\ans{26}{1}{Each access touches a different row (page), so there are 128 $\times$ 128 = 16,384 faults; row-by-row access gives only 128.}
\ans{27}{2}{It makes process creation fast, especially when exec() follows immediately.}
\ans{28}{2}{The lazy swapper (pager) never swaps a page in unless it is needed.}
\ans{29}{3}{Locality of reference keeps performance reasonable.}
\ans{30}{1}{A clean page can simply be overwritten when it is replaced.}
\ans{31}{2}{MFU, in contrast, replaces the most frequently used page.}
\ans{32}{4}{Both need hardware support; reference bits only approximate LRU.}
\ans{33}{1}{Locality is why demand paging and caching work.}
\ans{34}{4}{After mmap(), reads and writes to memory reach the file through the paging system.}
\ans{35}{1}{Prepaging pays off only if most of the pages brought in are actually used.}
\ans{36}{4}{Replacing the page during the transfer would corrupt the data.}
\ans{37}{1}{Any changes the child makes are visible to the parent.}
\ans{38}{1}{The victim is written out if it is dirty, then the new page is read in.}
\ans{39}{3}{Processes spend most of their time waiting on the paging device.}
\ans{40}{2}{Both are true, but the TLB only speeds up translation; it is not why virtual memory can exceed physical memory.}
\ans{41}{2}{Each distinct page faults at least once when first referenced.}
\ans{42}{3}{The clock hand advances, clearing reference bits, until it finds a page with bit 0.}
\ans{43}{1}{These are two sides of the page-size trade-off.}
\ans{44}{2}{LFU and MFU keep a count of references to each page.}
\ans{45}{1}{Linux uses swap partitions or swap files; Windows uses pagefile.sys.}
\ans{46}{1}{The OS then checks whether the reference is illegal or the page is simply not in memory.}
\ans{47}{2}{Too small a $\Delta$ misses the locality; too large overlaps several localities.}
\end{answers}

\section{Storage Management (disk structure, disk scheduling, RAID, disk management)}

\begin{question}{H}
A disk has 200 cylinders (0 to 199). The request queue is 98, 183, 37, 122, 14, 124, 65, 67 and the head is at cylinder 53. The total head movement under FCFS scheduling is
\optsshort{236}{640}{208}{382}
\end{question}

\begin{question}{H}
A disk has 200 cylinders (0 to 199). The request queue is 98, 183, 37, 122, 14, 124, 65, 67 and the head is at cylinder 53. The total head movement under SSTF scheduling is
\optsshort{299}{208}{236}{640}
\end{question}

\begin{question}{H}
A disk has 200 cylinders (0 to 199). The request queue is 98, 183, 37, 122, 14, 124, 65, 67 and the head is at cylinder 53, moving towards cylinder 0. The total head movement under SCAN scheduling is
\optsshort{236}{331}{382}{208}
\end{question}

\begin{question}{M}
A disk has 200 cylinders (0 to 199). The request queue is 98, 183, 37, 122, 14, 124, 65, 67 and the head is at cylinder 53, moving towards cylinder 0. The total head movement under LOOK scheduling is
\optsshort{299}{236}{322}{208}
\end{question}

\begin{question}{H}
A disk has 200 cylinders (0 to 199). The request queue is 98, 183, 37, 122, 14, 124, 65, 67 and the head is at cylinder 53, moving towards higher cylinders. The total head movement under C-SCAN, counting the return seek from 199 to 0, is
\optsshort{382}{331}{386}{322}
\end{question}

\begin{question}{M}
A disk has 200 cylinders (0 to 199). The request queue is 98, 183, 37, 122, 14, 124, 65, 67 and the head is at cylinder 53, moving towards higher cylinders. The total head movement under C-LOOK, counting the jump back, is
\optsshort{322}{382}{331}{299}
\end{question}

\begin{question}{H}
A disk has 200 cylinders (0 to 199). The request queue is 82, 170, 43, 140, 24, 16, 190 and the head is at cylinder 50. The total head movement under SSTF scheduling is
\optsshort{642}{314}{332}{208}
\end{question}

\begin{question}{M}
A disk has 200 cylinders (0 to 199). The request queue is 82, 170, 43, 140, 24, 16, 190 and the head is at cylinder 50, moving towards higher cylinders. The total head movement under SCAN scheduling is
\optsshort{391}{314}{332}{341}
\end{question}

\begin{question}{M}
A disk has 200 cylinders (0 to 199). The request queue is 82, 170, 43, 140, 24, 16, 190 and the head is at cylinder 50, moving towards higher cylinders. The total head movement under LOOK scheduling is
\optsshort{208}{332}{341}{314}
\end{question}

\begin{question}{M}
A disk has 5000 cylinders (0 to 4999). The head is at cylinder 143 and the queue is 86, 1470, 913, 1774, 948, 1509, 1022, 1750, 130. The total head movement under SSTF is
\optsshort{1745}{9769}{7081}{3319}
\end{question}

\begin{question}{M}
A disk has 5000 cylinders (0 to 4999). The head is at cylinder 143, moving towards higher cylinders, and the queue is 86, 1470, 913, 1774, 948, 1509, 1022, 1750, 130. The total head movement under LOOK is
\optsshort{3319}{3363}{1745}{9769}
\end{question}

\begin{question}{H}
Which disk scheduling algorithm may cause starvation of requests that are far from the current head position?
\optsshort{SSTF}{FCFS}{C-SCAN}{SCAN}
\end{question}

\begin{question}{H}
Match List I with List II.\\
List I (Disk scheduling algorithm): A. SCAN  B. C-SCAN  C. LOOK  D. SSTF\\
List II (Characteristic): I. Selects the request with the minimum seek time from the current head position  II. Like the elevator algorithm, but the arm goes only as far as the last request in each direction  III. Arm moves end to end, servicing requests in both directions (elevator algorithm)  IV. Services requests in one direction only and returns to the beginning without servicing on the return trip
\optsshort{A-II, B-IV, C-III, D-I}{A-III, B-II, C-IV, D-I}{A-III, B-IV, C-II, D-I}{A-IV, B-III, C-II, D-I}
\end{question}

\begin{question}{H}
Given below are two statements: one is labelled as Assertion A and the other is labelled as Reason R.\\
Assertion A: C-SCAN provides a more uniform waiting time than SCAN.\\
Reason R: C-SCAN treats the cylinders as a circular list, so requests at both ends wait about the same time.
\optslong{Both A and R are true and R is the correct explanation of A}{Both A and R are true but R is NOT the correct explanation of A}{A is true but R is false}{A is false but R is true}
\end{question}

\begin{question}{H}
A disk has 16 surfaces, 128 tracks per surface, 256 sectors per track and 512 bytes per sector. Its capacity is
\optsshort{64 MB}{128 MB}{512 MB}{256 MB}
\end{question}

\begin{question}{H}
A disk rotates at 7200 rpm. The average rotational latency is approximately
\optsshort{2.08 ms}{8.33 ms}{16.67 ms}{4.17 ms}
\end{question}

\begin{question}{M}
A disk rotating at 15,000 rpm has an average rotational latency of
\optsshort{8 ms}{1 ms}{4 ms}{2 ms}
\end{question}

\begin{question}{H}
A disk has an average seek time of 4 ms, rotates at 6000 rpm and has 100 sectors per track. The average time to read one sector is
\optsshort{9.1 ms}{14.1 ms}{10 ms}{9 ms}
\end{question}

\begin{question}{M}
A disk rotates at 6000 rpm and has 100 sectors of 512 bytes per track. The maximum data transfer rate is
\optsshort{512 KB/s}{5.12 MB/s}{6 MB/s}{51.2 MB/s}
\end{question}

\begin{question}{M}
The time for the disk arm to move the heads to the cylinder containing the desired sector is called
\optsshort{positioning time}{transfer time}{seek time}{rotational latency}
\end{question}

\begin{question}{H}
Match List I with List II.\\
List I (RAID level): A. RAID 0  B. RAID 1  C. RAID 5  D. RAID 6\\
List II (Organisation): I. Block-interleaved distributed parity  II. Mirroring  III. P + Q redundancy, tolerating two disk failures  IV. Block striping without redundancy
\optsshort{A-IV, B-I, C-II, D-III}{A-IV, B-II, C-III, D-I}{A-II, B-IV, C-I, D-III}{A-IV, B-II, C-I, D-III}
\end{question}

\begin{question}{H}
Given below are two statements.\\
Statement I: RAID 0 provides redundancy by mirroring every block.\\
Statement II: RAID 4 stores parity on a dedicated disk, which can become a bottleneck for writes.
\optslong{Both Statement I and Statement II are true}{Both Statement I and Statement II are false}{Statement I is true but Statement II is false}{Statement I is false but Statement II is true}
\end{question}

\begin{question}{H}
A RAID 5 array is built from 5 disks of 1 TB each. Its usable capacity is
\optsshort{5 TB}{4 TB}{2.5 TB}{1 TB}
\end{question}

\begin{question}{M}
A RAID 6 array is built from 6 disks of 2 TB each. Its usable capacity is
\optsshort{6 TB}{10 TB}{12 TB}{8 TB}
\end{question}

\begin{question}{M}
A RAID 1+0 (RAID 10) array uses 4 disks of 1 TB each. Its usable capacity is
\optsshort{1 TB}{3 TB}{2 TB}{4 TB}
\end{question}

\begin{question}{M}
In a mirrored pair, each disk has a mean time to failure of 100,000 hours and the mean time to repair is 10 hours. Assuming independent failures, the mean time to data loss is
\optsshort{10\textsuperscript{6} hours}{5 $\times$ 10\textsuperscript{8} hours}{10\textsuperscript{10} hours}{10\textsuperscript{5} hours}
\end{question}

\begin{question}{M}
Given below are two statements.\\
Statement I: RAID 1+0 (striped mirrors) is more robust than RAID 0+1 (mirrored stripes) when a single disk fails.\\
Statement II: In RAID 0+1, the failure of a single disk makes an entire stripe unavailable.
\optslong{Both Statement I and Statement II are true}{Both Statement I and Statement II are false}{Statement I is true but Statement II is false}{Statement I is false but Statement II is true}
\end{question}

\begin{question}{M}
Which RAID level uses bit-level or byte-level striping with a dedicated parity disk?
\optsshort{RAID 6}{RAID 5}{RAID 3}{RAID 1}
\end{question}

\begin{question}{M}
Arrange the steps in preparing a new disk for use in the correct order.\\
A. Logical formatting (creating a file system)  B. Low-level (physical) formatting  C. Partitioning
\optsshort{B, C, A}{A, B, C}{C, B, A}{B, A, C}
\end{question}

\begin{question}{M}
Replacing a bad sector logically with a spare sector that the controller maps to the same sector number is called
\optsshort{garbage collection}{sector sparing (forwarding)}{wear leveling}{sector slipping}
\end{question}

\begin{question}{M}
Given below are two statements.\\
Statement I: Solid-state drives have no moving parts, so they have no seek time or rotational latency.\\
Statement II: SSTF scheduling is essential for SSDs because seek time dominates their access time.
\optslong{Both Statement I and Statement II are true}{Both Statement I and Statement II are false}{Statement I is true but Statement II is false}{Statement I is false but Statement II is true}
\end{question}

\begin{question}{M}
Wear leveling in NAND flash storage is used to
\optslong{spread erase cycles evenly across blocks so that no block wears out early}{increase seek time}{reduce rotational latency}{compress data}
\end{question}

\begin{question}{M}
In NAND flash storage, which of the following is true?
\optslong{Reads are slower than writes}{There is no limit on erase cycles}{Data can be overwritten in place}{A block must be erased before its pages can be rewritten}
\end{question}

\begin{question}{M}
Match List I with List II.\\
List I (Storage attachment): A. Host-attached storage  B. Network-attached storage (NAS)  C. Storage-area network (SAN)\\
List II (Description): I. Private network connecting servers and storage units, e.g. Fibre Channel  II. Accessed through local I/O ports such as SATA or SCSI  III. Accessed remotely over an IP network using protocols such as NFS or CIFS
\optsshort{A-II, B-III, C-I}{A-II, B-I, C-III}{A-I, B-III, C-II}{A-III, B-II, C-I}
\end{question}

\begin{question}{M}
The area of a disk partition that contains the bootstrap code is called the
\optsshort{superblock}{swap space}{inode table}{boot block}
\end{question}

\begin{question}{M}
In the disk queue, if requests arrive only one at a time and the queue is always empty or has a single request, then all disk scheduling algorithms
\optsshort{cause starvation}{always perform like SSTF}{perform the same as FCFS}{perform worse than FCFS}
\end{question}

\begin{question}{M}
Which disk scheduling algorithm is also known as the elevator algorithm?
\optsshort{SCAN}{SSTF}{FCFS}{C-LOOK}
\end{question}

\begin{question}{M}
Given below are two statements.\\
Statement I: In C-LOOK, the arm goes only as far as the final request in each direction before jumping back.\\
Statement II: In C-LOOK, requests are serviced on the return trip.
\optslong{Both Statement I and Statement II are true}{Both Statement I and Statement II are false}{Statement I is true but Statement II is false}{Statement I is false but Statement II is true}
\end{question}

\begin{question}{M}
The Linux deadline I/O scheduler maintains separate queues for reads and writes and
\optslong{gives priority to writes}{disables merging of requests}{services only the oldest request}{uses expiration times to prevent starvation of requests}
\end{question}

\begin{question}{M}
In a disk, the set of tracks at the same arm position on all surfaces is called a
\optsshort{cylinder}{sector}{block}{cluster}
\end{question}

\begin{question}{M}
Disk bandwidth is defined as
\optslong{the number of requests per second}{the number of cylinders per second}{the rotational speed}{the total number of bytes transferred divided by the total time between the first request and the completion of the last transfer}
\end{question}

\begin{question}{L}
Which RAID level offers the best write performance among levels that provide redundancy, at the cost of 50\% storage overhead?
\optsshort{RAID 6}{RAID 1 (or 1+0)}{RAID 5}{RAID 4}
\end{question}

\begin{question}{M}
Swap space on Linux can be placed
\optslong{in a separate raw partition or in a file within a normal file system}{only in a raw partition}{only on optical media}{only in main memory}
\end{question}

\begin{question}{M}
With a request queue of 10, 22, 20, 2, 40, 6, 38 and the head at cylinder 20, which algorithm services the request at cylinder 20 first?
\optsshort{FCFS only}{SSTF only}{Both SSTF and SCAN}{Neither}
\end{question}

\begin{question}{M}
Given below are two statements: one is labelled as Assertion A and the other is labelled as Reason R.\\
Assertion A: Disk scheduling is less important for SSDs than for magnetic disks.\\
Reason R: SSD access time does not depend on the physical location of the data.
\optslong{Both A and R are true and R is the correct explanation of A}{Both A and R are true but R is NOT the correct explanation of A}{A is true but R is false}{A is false but R is true}
\end{question}

\begin{answers}
\ans{1}{2}{53$\rightarrow$98$\rightarrow$183$\rightarrow$37$\rightarrow$122$\rightarrow$14$\rightarrow$124$\rightarrow$65$\rightarrow$67 = 45 + 85 + 146 + 85 + 108 + 110 + 59 + 2 = 640 cylinders.}
\ans{2}{3}{Order 65, 67, 37, 14, 98, 122, 124, 183: 12 + 2 + 30 + 23 + 84 + 24 + 2 + 59 = 236.}
\ans{3}{1}{53 $\rightarrow$ 0 (53 cylinders), then 0 $\rightarrow$ 183 (183 cylinders): 236.}
\ans{4}{4}{53 $\rightarrow$ 14 (39 cylinders), then 14 $\rightarrow$ 183 (169 cylinders): 208.}
\ans{5}{1}{53 $\rightarrow$ 199 (146), 199 $\rightarrow$ 0 (199), 0 $\rightarrow$ 37 (37): 382.}
\ans{6}{1}{53 $\rightarrow$ 183 (130), jump 183 $\rightarrow$ 14 (169), 14 $\rightarrow$ 37 (23): 322.}
\ans{7}{4}{Order 43, 24, 16, 82, 140, 170, 190: 7 + 19 + 8 + 66 + 58 + 30 + 20 = 208.}
\ans{8}{3}{50 $\rightarrow$ 199 (149), then 199 $\rightarrow$ 16 (183): 332.}
\ans{9}{4}{50 $\rightarrow$ 190 (140), then 190 $\rightarrow$ 16 (174): 314.}
\ans{10}{1}{Order 130, 86, 913, 948, 1022, 1470, 1509, 1750, 1774: 13 + 44 + 827 + 35 + 74 + 448 + 39 + 241 + 24 = 1745.}
\ans{11}{1}{143 $\rightarrow$ 1774 (1631), then 1774 $\rightarrow$ 86 (1688): 3319.}
\ans{12}{1}{SSTF always serves the nearest request, so a stream of nearby requests can postpone distant ones indefinitely.}
\ans{13}{3}{Standard disk scheduling algorithms.}
\ans{14}{1}{In SCAN, cylinders in the middle are visited twice as often as those at the ends.}
\ans{15}{4}{2\textsuperscript{4} $\times$ 2\textsuperscript{7} $\times$ 2\textsuperscript{8} $\times$ 2\textsuperscript{9} = 2\textsuperscript{28} bytes = 256 MB.}
\ans{16}{4}{One rotation takes 60,000 / 7200 = 8.33 ms; the average latency is half a rotation.}
\ans{17}{4}{One rotation takes 4 ms; half of it is 2 ms.}
\ans{18}{1}{Seek 4 ms + average latency 5 ms + transfer 10 ms / 100 = 0.1 ms, total 9.1 ms.}
\ans{19}{2}{100 rotations per second $\times$ 100 $\times$ 512 bytes = 5.12 $\times$ 10\textsuperscript{6} bytes per second.}
\ans{20}{3}{Positioning (random-access) time = seek time + rotational latency.}
\ans{21}{4}{Standard RAID levels.}
\ans{22}{4}{RAID 0 is striping with no redundancy; mirroring is RAID 1.}
\ans{23}{2}{One disk's worth of space holds parity: (5 $-$ 1) $\times$ 1 TB = 4 TB.}
\ans{24}{4}{Two disks' worth of space holds the P and Q redundancy: (6 $-$ 2) $\times$ 2 TB = 8 TB.}
\ans{25}{3}{Mirroring halves the raw capacity.}
\ans{26}{2}{MTTF\textsuperscript{2} / (2 $\times$ MTTR) = 10\textsuperscript{10} / 20 = 5 $\times$ 10\textsuperscript{8} hours.}
\ans{27}{1}{In RAID 1+0 a failed disk affects only its mirror pair.}
\ans{28}{3}{RAID 4 uses block-level striping with a dedicated parity disk.}
\ans{29}{1}{Physical formatting creates sectors, partitioning divides the disk, and logical formatting writes file-system structures.}
\ans{30}{2}{Sector slipping instead shifts all sectors down by one.}
\ans{31}{3}{SSDs have no seek time, so simple FCFS-like scheduling is used.}
\ans{32}{1}{Each NAND block tolerates a limited number of erase cycles.}
\ans{33}{4}{The flash translation layer maps logical blocks to physical pages and performs garbage collection.}
\ans{34}{1}{iSCSI carries SCSI over IP networks.}
\ans{35}{4}{The boot block is read by the firmware to start the boot loader.}
\ans{36}{3}{Scheduling only helps when there is a queue of requests to reorder.}
\ans{37}{1}{The arm behaves like an elevator, serving requests on the way up and on the way down.}
\ans{38}{3}{Like C-SCAN, C-LOOK services requests in one direction only.}
\ans{39}{4}{Reads are usually given shorter deadlines because processes block on them.}
\ans{40}{1}{Switching surfaces within a cylinder needs no seek.}
\ans{41}{4}{Disk scheduling aims to improve both access time and bandwidth.}
\ans{42}{2}{Mirroring needs no parity computation; each write simply goes to both disks.}
\ans{43}{1}{A raw swap partition avoids file-system overhead; a swap file is more flexible.}
\ans{44}{3}{The head is already at cylinder 20, which is at distance zero; SCAN services it as it starts moving in either direction.}
\ans{45}{1}{There is no mechanical movement to minimise.}
\end{answers}

\section{File Systems and I/O Systems (access methods, directories, allocation, free-space management, I/O hardware and kernel I/O)}

\begin{question}{H}
A UNIX file system uses 1 KB blocks and 4-byte block addresses. An inode has 10 direct pointers, one single-indirect pointer and one double-indirect pointer. The maximum file size is approximately
\optsshort{64 MB}{16 GB}{266 KB}{10 KB}
\end{question}

\begin{question}{H}
A file system uses 4 KB blocks and 4-byte block addresses. An inode has 12 direct pointers and one single-, one double- and one triple-indirect pointer. The maximum file size is approximately
\optsshort{16 GB}{4 GB}{16 TB}{4 TB}
\end{question}

\begin{question}{M}
A file system uses 4 KB blocks and 4-byte block addresses. An inode has 12 direct pointers, one single-indirect pointer and one double-indirect pointer. The maximum file size is approximately
\optsshort{4 MB}{4 TB}{4 GB}{16 GB}
\end{question}

\begin{question}{M}
An inode has 10 direct pointers and 1 KB blocks. The byte at offset 20,000 of a file is found through the
\optsshort{direct pointers}{triple-indirect pointer}{double-indirect pointer}{single-indirect pointer}
\end{question}

\begin{question}{H}
A 1 TB disk with 4 KB blocks manages free space with a bit vector. The size of the bit vector is
\optsshort{4 MB}{256 MB}{1 GB}{32 MB}
\end{question}

\begin{question}{H}
Match List I with List II.\\
List I (File allocation method): A. Contiguous allocation  B. Linked allocation  C. Indexed allocation\\
List II (Characteristic): I. Uses an index block holding all the block addresses; supports direct access without external fragmentation  II. Supports direct access easily but suffers from external fragmentation and difficulty growing files  III. No external fragmentation, but efficient only for sequential access
\optsshort{A-II, B-III, C-I}{A-II, B-I, C-III}{A-I, B-III, C-II}{A-III, B-II, C-I}
\end{question}

\begin{question}{H}
Given below are two statements.\\
Statement I: The File Allocation Table (FAT) used by MS-DOS is a variation of linked allocation.\\
Statement II: In FAT, random access time is improved because the table can be cached and followed without reading every data block.
\optslong{Both Statement I and Statement II are true}{Both Statement I and Statement II are false}{Statement I is true but Statement II is false}{Statement I is false but Statement II is true}
\end{question}

\begin{question}{H}
A file of 100 blocks already exists, and its file-control block is in memory. How many disk I/O operations are needed to add a block at the beginning of the file under contiguous, linked and indexed allocation respectively (assuming there is room to grow only at the end for contiguous allocation)?
\optsshort{201, 2 and 1}{201, 1 and 1}{100, 1 and 2}{1, 101 and 1}
\end{question}

\begin{question}{H}
Match List I with List II.\\
List I (Directory structure): A. Single-level  B. Two-level  C. Tree-structured  D. Acyclic graph\\
List II (Feature): I. Allows files and subdirectories to be shared through links  II. All files in one directory, so names must be unique  III. A separate user file directory for each user under a master file directory  IV. Arbitrary depth with absolute and relative path names
\optsshort{A-II, B-III, C-IV, D-I}{A-II, B-IV, C-III, D-I}{A-III, B-II, C-IV, D-I}{A-II, B-III, C-I, D-IV}
\end{question}

\begin{question}{H}
Given below are two statements.\\
Statement I: A hard link can point to a file in a different file system.\\
Statement II: When the target of a symbolic link is deleted, the symbolic link is left dangling.
\optslong{Both Statement I and Statement II are true}{Both Statement I and Statement II are false}{Statement I is true but Statement II is false}{Statement I is false but Statement II is true}
\end{question}

\begin{question}{H}
The UNIX command chmod 754 file sets the permissions to
\optsshort{rw-r-xr--}{rwxrw-r--}{rwxr-x--x}{rwxr-xr--}
\end{question}

\begin{question}{H}
If the umask is 022, a newly created regular file (default mode 666) gets the permissions
\optsshort{rwxr-xr-x}{rw-------}{rw-r--r--}{rw-rw-rw-}
\end{question}

\begin{question}{M}
Match List I with List II.\\
List I (Access method): A. Sequential access  B. Direct (relative) access  C. Indexed access\\
List II (Description): I. Blocks can be read or written in any order using a relative block number  II. An index of pointers to blocks is searched to find a record  III. Records are processed in order, one after the other, like a tape
\optsshort{A-I, B-III, C-II}{A-III, B-II, C-I}{A-III, B-I, C-II}{A-II, B-I, C-III}
\end{question}

\begin{question}{H}
Arrange the layers of a layered file system from the top (closest to the application) to the bottom.\\
A. Basic file system  B. Logical file system  C. I/O control  D. File-organisation module
\optsshort{D, B, A, C}{B, D, C, A}{B, D, A, C}{B, A, D, C}
\end{question}

\begin{question}{M}
In UNIX, the file control block (FCB) is called the
\optsshort{inode}{dentry}{file descriptor}{superblock}
\end{question}

\begin{question}{M}
In the Linux virtual file system (VFS), the four primary object types are
\optslong{process, thread, task and signal}{inode, file, superblock and dentry}{block, sector, track and cylinder}{file, directory, link and pipe}
\end{question}

\begin{question}{M}
A bit vector 0000 0000 0000 0000 0000 0000 0010 1100 (16-bit words, 1 means free) is used for free-space management. The number of the first free block is
\optsshort{34}{50}{42}{26}
\end{question}

\begin{question}{M}
Which free-space management technique stores the address of the first free block and the number of free contiguous blocks that follow it?
\optsshort{Grouping}{Bit vector}{Linked list}{Counting}
\end{question}

\begin{question}{M}
Given below are two statements.\\
Statement I: A linked-list free-space scheme needs no extra space other than a pointer to the first free block.\\
Statement II: Traversing the free list to find many free blocks is efficient.
\optslong{Both Statement I and Statement II are true}{Both Statement I and Statement II are false}{Statement I is true but Statement II is false}{Statement I is false but Statement II is true}
\end{question}

\begin{question}{M}
Match List I with List II.\\
List I (Consistency semantics): A. UNIX semantics  B. Session semantics  C. Immutable-shared-file semantics\\
List II (Description): I. A file declared shared cannot be modified  II. Writes are visible to other users immediately  III. Writes become visible only to sessions that open the file after it is closed (as in AFS)
\optsshort{A-II, B-III, C-I}{A-I, B-III, C-II}{A-III, B-II, C-I}{A-II, B-I, C-III}
\end{question}

\begin{question}{M}
A file-system journal (log) is used to
\optslong{recover quickly after a crash by replaying or discarding logged metadata transactions}{encrypt data}{speed up random reads}{compress files}
\end{question}

\begin{question}{M}
The UNIX fsck utility is used for
\optslong{mounting remote file systems}{creating backups}{checking and repairing file-system consistency}{formatting a disk}
\end{question}

\begin{question}{M}
Mounting a file system means
\optslong{copying it to another disk}{deleting all files}{attaching it to the directory tree at a specified mount point}{formatting it}
\end{question}

\begin{question}{M}
Given below are two statements.\\
Statement I: Advisory file locks are enforced by the operating system.\\
Statement II: A shared lock lets several processes read a file concurrently, while an exclusive lock allows only one process.
\optslong{Both Statement I and Statement II are true}{Both Statement I and Statement II are false}{Statement I is true but Statement II is false}{Statement I is false but Statement II is true}
\end{question}

\begin{question}{M}
For sequential file access, which caching techniques improve performance?
\optsshort{Free-behind and read-ahead}{LRU only}{Write-through only}{Random replacement}
\end{question}

\begin{question}{M}
A unified buffer cache
\optslong{is a separate cache for each file}{caches file data and memory-mapped pages in one page cache, avoiding double caching}{is used only by the kernel}{is stored on disk}
\end{question}

\begin{question}{H}
Match List I with List II.\\
List I (I/O technique): A. Polling  B. Interrupt-driven I/O  C. Direct memory access (DMA)\\
List II (Description): I. The device controller notifies the CPU when it is ready  II. A controller transfers a whole block between the device and memory without involving the CPU for each byte  III. The CPU repeatedly reads the device status register (busy waiting)
\optsshort{A-III, B-II, C-I}{A-I, B-III, C-II}{A-II, B-I, C-III}{A-III, B-I, C-II}
\end{question}

\begin{question}{H}
Given below are two statements: one is labelled as Assertion A and the other is labelled as Reason R.\\
Assertion A: DMA is preferred for large data transfers.\\
Reason R: During DMA the CPU must copy every byte from the device controller into memory.
\optslong{Both A and R are true and R is the correct explanation of A}{Both A and R are true but R is NOT the correct explanation of A}{A is true but R is false}{A is false but R is true}
\end{question}

\begin{question}{M}
When the DMA controller takes the memory bus away from the CPU for one word at a time, this is called
\optsshort{cycle stealing}{burst mode}{bus mastering only}{polling}
\end{question}

\begin{question}{M}
Match List I with List II.\\
List I (Kernel I/O service): A. Buffering  B. Caching  C. Spooling  D. Device reservation\\
List II (Purpose): I. Gives a process exclusive access to a device  II. Holds output for a device that cannot accept interleaved data streams, such as a printer  III. Copes with speed and transfer-size mismatches between devices  IV. Keeps a copy of data in faster storage for reuse
\optsshort{A-III, B-II, C-IV, D-I}{A-III, B-IV, C-I, D-II}{A-IV, B-III, C-II, D-I}{A-III, B-IV, C-II, D-I}
\end{question}

\begin{question}{M}
Double buffering is used to
\optslong{increase seek time}{allow a producer to fill one buffer while a consumer empties the other}{avoid DMA}{reduce memory usage}
\end{question}

\begin{question}{M}
Match List I with List II.\\
List I (Device type): A. Block device  B. Character device  C. Network device\\
List II (Example/interface): I. Keyboard, accessed with get() and put() one character at a time  II. Disk, accessed with read(), write() and seek() on fixed-size blocks  III. Accessed through the socket interface
\optsshort{A-I, B-II, C-III}{A-II, B-III, C-I}{A-II, B-I, C-III}{A-III, B-I, C-II}
\end{question}

\begin{question}{M}
Given below are two statements.\\
Statement I: A non-blocking I/O call suspends the process until all the requested data has been transferred.\\
Statement II: An asynchronous I/O call blocks the process until the transfer completes.
\optslong{Both Statement I and Statement II are true}{Both Statement I and Statement II are false}{Statement I is true but Statement II is false}{Statement I is false but Statement II is true}
\end{question}

\begin{question}{M}
Vectored I/O (scatter-gather), as provided by readv() in UNIX, allows
\optslong{bypassing the kernel}{DMA without a controller}{several devices to be read at once}{one system call to perform I/O involving multiple memory buffers}
\end{question}

\begin{question}{M}
In memory-mapped I/O
\optslong{only DMA is used}{special I/O instructions such as IN and OUT are required}{device control registers are mapped into the address space of the processor and accessed with ordinary load and store instructions}{devices cannot be accessed by the kernel}
\end{question}

\begin{question}{M}
The software component that hides the differences among device controllers from the kernel I/O subsystem is the
\optsshort{scheduler}{file system}{device driver}{interrupt vector}
\end{question}

\begin{question}{M}
Which file attribute is used to identify a file uniquely within the file system, independent of its name?
\optslong{Type}{Identifier (e.g. inode number)}{Size}{Location}
\end{question}

\begin{question}{M}
In UNIX, a file's type can be recognised by its contents, for example by a magic number at the beginning of the file. Windows instead relies mainly on
\optsshort{access control lists}{file extensions}{inode numbers}{the superblock}
\end{question}

\begin{question}{M}
Which of the following operations is NOT a basic file operation?
\optsshort{Truncate}{Create}{Compile}{Read}
\end{question}

\begin{question}{M}
The per-process open-file table entry typically contains
\optslong{the disk block addresses}{the size of the disk}{the file pointer (current position) and access rights}{the boot block}
\end{question}

\begin{question}{M}
Given below are two statements.\\
Statement I: Linear-list directory implementation is simple but slow for searching.\\
Statement II: Hash-table directory implementation speeds up searching but must handle collisions.
\optslong{Both Statement I and Statement II are true}{Both Statement I and Statement II are false}{Statement I is true but Statement II is false}{Statement I is false but Statement II is true}
\end{question}

\begin{question}{M}
A FAT16 file system with 2\textsuperscript{16} clusters of 32 KB each can address a volume of at most
\optsshort{32 GB}{512 MB}{4 GB}{2 GB}
\end{question}

\begin{question}{M}
Which backup strategy copies only the files changed since the last backup of any kind?
\optsshort{Full backup}{Incremental backup}{Image backup}{Mirror backup}
\end{question}

\begin{question}{M}
Given below are two statements.\\
Statement I: In a general-graph directory structure, cycles can occur, so garbage collection may be needed to reclaim space.\\
Statement II: An acyclic-graph directory structure allows a file to have more than one path name.
\optslong{Both Statement I and Statement II are true}{Both Statement I and Statement II are false}{Statement I is true but Statement II is false}{Statement I is false but Statement II is true}
\end{question}

\begin{question}{M}
In the indexed allocation scheme, the linked scheme, multilevel index and combined scheme are used to
\optslong{improve security}{eliminate the index block}{handle files too large for a single index block}{reduce seek time}
\end{question}

\begin{question}{L}
Which file system is used by default on most modern Linux distributions?
\optsshort{ext4}{FAT32}{NTFS}{HFS+}
\end{question}

\begin{question}{M}
In UNIX, the hard link count of a file is stored in its
\optsshort{directory entry}{data block}{inode}{superblock}
\end{question}

\begin{question}{M}
The life cycle of a blocking read request ends with
\optslong{the process being moved from the wait queue to the ready queue after the device driver and interrupt handler complete the transfer}{the file system deleting the file}{the device controller calling the scheduler directly}{the CPU polling the device}
\end{question}

\begin{answers}
\ans{1}{1}{(10 + 256 + 256\textsuperscript{2}) $\times$ 1 KB = 65,802 KB $\approx$ 64.26 MB.}
\ans{2}{4}{Each indirect block holds 1024 pointers: (12 + 1024 + 1024\textsuperscript{2} + 1024\textsuperscript{3}) $\times$ 4 KB $\approx$ 4 TB, dominated by the triple-indirect part.}
\ans{3}{3}{(12 + 1024 + 1024\textsuperscript{2}) $\times$ 4 KB $\approx$ 4.004 GB.}
\ans{4}{4}{20,000 / 1024 gives logical block 19; blocks 0--9 are direct, so block 19 is the 10th entry of the single-indirect block.}
\ans{5}{4}{2\textsuperscript{40} / 2\textsuperscript{12} = 2\textsuperscript{28} blocks; 2\textsuperscript{28} bits = 2\textsuperscript{25} bytes = 32 MB.}
\ans{6}{1}{The three classical disk allocation methods.}
\ans{7}{1}{The FAT holds one entry per block containing the number of the next block.}
\ans{8}{2}{Contiguous: read and rewrite 100 blocks shifted by one, plus write the new block (201). Linked and indexed: write the new block only; the FCB in memory is updated.}
\ans{9}{1}{Standard directory structures.}
\ans{10}{4}{A hard link refers to an inode number, which is meaningful only within one file system.}
\ans{11}{4}{7 = rwx (owner), 5 = r-x (group), 4 = r-- (others).}
\ans{12}{3}{666 with the 022 bits cleared gives 644 = rw-r--r--; directories (777) get 755.}
\ans{13}{3}{Standard file access methods.}
\ans{14}{3}{Application programs $\rightarrow$ logical file system $\rightarrow$ file-organisation module $\rightarrow$ basic file system $\rightarrow$ I/O control $\rightarrow$ devices.}
\ans{15}{1}{The inode stores ownership, permissions, size, timestamps and block pointers.}
\ans{16}{2}{VFS lets the kernel treat different file systems uniformly.}
\ans{17}{4}{The first 16-bit word is all 0s; in the second word the first 1 is at offset 10, so the block number is 16 $\times$ 1 + 10 = 26.}
\ans{18}{4}{Counting works well when space is allocated and freed in contiguous runs.}
\ans{19}{3}{Traversing the list requires reading each block, which is slow, though it is rarely needed.}
\ans{20}{1}{Standard file-sharing consistency semantics.}
\ans{21}{1}{ext3, ext4 and NTFS are journaling file systems.}
\ans{22}{3}{It compares directory structures with the data blocks and fixes inconsistencies.}
\ans{23}{3}{The mounted file system's root then appears at the mount point.}
\ans{24}{4}{Mandatory locks are enforced by the OS; advisory locks rely on cooperating processes.}
\ans{25}{1}{A page is freed as soon as the next one is requested, and the next pages are read in advance.}
\ans{26}{2}{Without it, data read through read() and through mmap() could be cached twice.}
\ans{27}{4}{Three ways to perform I/O.}
\ans{28}{3}{The DMA controller does the copying; the CPU is interrupted only when the whole block is done.}
\ans{29}{1}{The CPU is briefly prevented from accessing main memory.}
\ans{30}{4}{Services of the kernel I/O subsystem.}
\ans{31}{2}{It decouples the producer and the consumer of the data.}
\ans{32}{3}{Common categories in the application I/O interface.}
\ans{33}{2}{Both calls return immediately: non-blocking I/O with whatever data is available, asynchronous I/O with notification later.}
\ans{34}{4}{It avoids several system calls and helps atomicity.}
\ans{35}{3}{Isolated (port-mapped) I/O uses separate I/O instructions.}
\ans{36}{3}{Each driver presents a uniform interface for its class of devices.}
\ans{37}{2}{Names are for people; identifiers are for the system.}
\ans{38}{2}{The extension (such as .exe or .docx) determines the associated program.}
\ans{39}{3}{Compiling is done by an application program, not by the file system.}
\ans{40}{3}{The system-wide table holds the location on disk, the open count and other per-file data.}
\ans{41}{1}{Hash tables with chaining are commonly used.}
\ans{42}{4}{2\textsuperscript{16} $\times$ 32 KB = 2\textsuperscript{31} bytes = 2 GB.}
\ans{43}{2}{A differential backup copies all files changed since the last full backup.}
\ans{44}{1}{Links allow sharing; cycles complicate traversal and deletion.}
\ans{45}{3}{The UNIX inode is an example of the combined scheme.}
\ans{46}{1}{ext4 uses extents and journaling; NTFS is the Windows file system.}
\ans{47}{3}{The file is deleted only when the link count drops to zero and no process has it open.}
\ans{48}{1}{The data is copied to the process's buffer and the process becomes ready.}
\end{answers}

\section{Protection, Security and Virtual Machines}

\begin{question}{H}
In the access-matrix model of protection, rows and columns represent respectively
\optsshort{domains and objects}{objects and domains}{users and passwords}{processes and files only}
\end{question}

\begin{question}{H}
Match List I with List II.\\
List I (Access-matrix implementation): A. Access control list  B. Capability list  C. Global table  D. Lock-key mechanism\\
List II (Description): I. Each domain has a list of objects with the operations allowed on them  II. Each object has unique bit patterns (locks) and each domain has keys  III. Each object has a list of domains with their access rights  IV. A set of ordered triples <domain, object, rights-set>
\optsshort{A-III, B-I, C-IV, D-II}{A-III, B-IV, C-I, D-II}{A-I, B-III, C-IV, D-II}{A-III, B-I, C-II, D-IV}
\end{question}

\begin{question}{H}
Given below are two statements.\\
Statement I: An access control list corresponds to a row of the access matrix.\\
Statement II: A capability list corresponds to a column of the access matrix.
\optslong{Both Statement I and Statement II are true}{Both Statement I and Statement II are false}{Statement I is true but Statement II is false}{Statement I is false but Statement II is true}
\end{question}

\begin{question}{H}
Given below are two statements: one is labelled as Assertion A and the other is labelled as Reason R.\\
Assertion A: Revocation of access rights is more difficult with capability lists than with access control lists.\\
Reason R: Capabilities are distributed throughout the system, so they must first be found before they can be revoked.
\optslong{Both A and R are true and R is the correct explanation of A}{Both A and R are true but R is NOT the correct explanation of A}{A is true but R is false}{A is false but R is true}
\end{question}

\begin{question}{M}
Which of the following are schemes for revoking capabilities?\\
A. Reacquisition  B. Back-pointers  C. Indirection  D. Keys  E. Aging\\
Choose the correct answer from the options given below:
\optsshort{A, B, C, D and E}{A and B only}{A, B, C and D only}{B, C and E only}
\end{question}

\begin{question}{H}
The principle of least privilege states that
\optslong{programs, users and systems should be given just enough privileges to perform their tasks}{all users should have administrator privileges}{only the kernel has privileges}{privileges should never change}
\end{question}

\begin{question}{M}
In UNIX, the setuid bit on an executable file means that
\optslong{the file is hidden}{the process runs with the user ID of the file's owner rather than of the user who runs it}{the file can be executed only by the owner}{the file cannot be deleted}
\end{question}

\begin{question}{M}
In the Intel protection-ring architecture, the operating-system kernel runs in
\optsshort{ring 2}{ring 1}{ring 0}{ring 3}
\end{question}

\begin{question}{H}
Match List I with List II.\\
List I (Program threat): A. Trojan horse  B. Trap door  C. Logic bomb  D. Buffer overflow\\
List II (Description): I. Code that misuses its environment while appearing to be useful  II. Writing beyond the bounds of an array on the stack to overwrite the return address  III. A hidden hole intentionally left by the designer that bypasses normal security  IV. Code that initiates a security incident only under certain conditions
\optsshort{A-I, B-III, C-IV, D-II}{A-I, B-IV, C-III, D-II}{A-III, B-I, C-IV, D-II}{A-I, B-III, C-II, D-IV}
\end{question}

\begin{question}{M}
Match List I with List II.\\
List I (Virus type): A. Boot virus  B. Macro virus  C. Polymorphic virus  D. Stealth virus\\
List II (Characteristic): I. Changes its signature each time it is installed to avoid detection  II. Infects the boot sector and runs every time the system is booted  III. Written in a high-level macro language, e.g. in a word-processor document  IV. Modifies parts of the system that could be used to detect it
\optsshort{A-II, B-I, C-III, D-IV}{A-II, B-III, C-I, D-IV}{A-III, B-II, C-I, D-IV}{A-II, B-III, C-IV, D-I}
\end{question}

\begin{question}{M}
The Morris worm of 1988 spread by exploiting
\optslong{the boot sector of floppy disks}{weaknesses in the UNIX finger, sendmail and rsh programs}{a macro in a word processor}{a flaw in the Windows registry}
\end{question}

\begin{question}{M}
Given below are two statements.\\
Statement I: Adding a random salt to each password before hashing defeats precomputed dictionary attacks on the password file.\\
Statement II: Salting makes two users with the same password have the same stored hash.
\optslong{Both Statement I and Statement II are true}{Both Statement I and Statement II are false}{Statement I is true but Statement II is false}{Statement I is false but Statement II is true}
\end{question}

\begin{question}{M}
A one-time password scheme protects against
\optslong{denial of service}{reuse of an eavesdropped password}{trap doors}{buffer overflows}
\end{question}

\begin{question}{M}
Two-factor authentication requires
\optslong{two different user names}{logging in twice}{two of the following: something you know, something you have, something you are}{two passwords}
\end{question}

\begin{question}{M}
Match List I with List II.\\
List I (Intrusion detection approach): A. Signature-based detection  B. Anomaly detection\\
List II (Description): I. Detects deviations from normal behaviour and can find previously unknown attacks, with more false alarms  II. Looks for patterns of known attacks and cannot detect new attacks
\optslong{A-II, B-I}{A-I, B-II}{Both use pattern matching only}{Both can detect only known attacks}
\end{question}

\begin{question}{M}
In the US Department of Defense TCSEC (Orange Book) classification, the division that provides minimal protection is
\optsshort{D}{A}{B}{C}
\end{question}

\begin{question}{M}
Tripwire is an example of a
\optslong{virus}{file-integrity checking tool that detects changes to files using stored signatures}{password cracker}{firewall}
\end{question}

\begin{question}{M}
Given below are two statements.\\
Statement I: In discretionary access control (DAC), the owner of an object decides who may access it.\\
Statement II: In mandatory access control (MAC), access is decided by system-wide policy that users cannot override.
\optslong{Both Statement I and Statement II are true}{Both Statement I and Statement II are false}{Statement I is true but Statement II is false}{Statement I is false but Statement II is true}
\end{question}

\begin{question}{M}
Role-based access control (RBAC) assigns privileges to
\optslong{roles, which are then assigned to users}{hardware devices}{network ports only}{individual files}
\end{question}

\begin{question}{M}
Which of the following is a network or system threat rather than a program threat?
\optsshort{Logic bomb}{Trap door}{Port scanning}{Trojan horse}
\end{question}

\begin{question}{H}
Match List I with List II.\\
List I (Hypervisor type): A. Type 0  B. Type 1  C. Type 2\\
List II (Example): I. VMware ESX, Citrix XenServer  II. VMware Workstation, Oracle VirtualBox  III. IBM LPARs, Oracle LDOMs (firmware-based)
\optsshort{A-III, B-II, C-I}{A-I, B-III, C-II}{A-III, B-I, C-II}{A-II, B-I, C-III}
\end{question}

\begin{question}{H}
Given below are two statements.\\
Statement I: In paravirtualization, the guest operating system is modified to work with the VMM.\\
Statement II: In paravirtualization, the guest operating system is unaware that it is running in a virtual machine.
\optslong{Both Statement I and Statement II are true}{Both Statement I and Statement II are false}{Statement I is true but Statement II is false}{Statement I is false but Statement II is true}
\end{question}

\begin{question}{H}
In the trap-and-emulate method of virtualization, when the guest kernel (running in virtual kernel mode, which is physical user mode) executes a privileged instruction
\optslong{the guest is terminated}{it causes a trap to the VMM, which emulates the instruction for the guest}{it executes directly on the hardware}{the instruction is ignored}
\end{question}

\begin{question}{H}
Given below are two statements: one is labelled as Assertion A and the other is labelled as Reason R.\\
Assertion A: Classic x86 processors could not be virtualized with pure trap-and-emulate.\\
Reason R: Intel VT-x and AMD-V add hardware support for virtualization.
\optslong{Both A and R are true and R is the correct explanation of A}{Both A and R are true but R is NOT the correct explanation of A}{A is true but R is false}{A is false but R is true}
\end{question}

\begin{question}{M}
Binary translation in a VMM such as VMware
\optslong{converts binary files to text}{translates virtual addresses to physical addresses}{rewrites sensitive guest kernel instructions into safe instruction sequences at run time}{compiles the guest OS from source}
\end{question}

\begin{question}{M}
Nested page tables (Intel EPT, AMD RVI/NPT) are a hardware feature that
\optslong{replace the hypervisor}{add more TLB entries}{encrypt guest memory}{speeds up memory virtualization by letting the hardware translate guest-physical to host-physical addresses}
\end{question}

\begin{question}{M}
Moving a running guest from one host to another with no interruption of service is called
\optsshort{cloning}{ballooning}{snapshotting}{live migration}
\end{question}

\begin{question}{M}
In memory ballooning, the VMM
\optslong{deletes guest processes}{installs a driver in the guest that allocates memory to force the guest to page out, returning the freed memory to the VMM}{adds memory chips to the host}{shares identical pages}
\end{question}

\begin{question}{M}
Match List I with List II.\\
List I (Virtualization category): A. Programming-environment virtualization  B. Emulation  C. Application containment\\
List II (Example): I. Linux containers, Solaris zones, BSD jails  II. QEMU running ARM code on an x86 host  III. Java Virtual Machine, .NET CLR
\optsshort{A-III, B-II, C-I}{A-III, B-I, C-II}{A-II, B-III, C-I}{A-I, B-II, C-III}
\end{question}

\begin{question}{M}
Given below are two statements.\\
Statement I: Containers share the host kernel and are therefore lighter weight than virtual machines.\\
Statement II: A container can run a guest operating system with a different kernel from the host.
\optslong{Both Statement I and Statement II are true}{Both Statement I and Statement II are false}{Statement I is true but Statement II is false}{Statement I is false but Statement II is true}
\end{question}

\begin{question}{M}
The Popek and Goldberg requirements for a virtual machine monitor are
\optslong{portability, compatibility and scalability}{speed, cost and size}{fidelity, performance and safety (resource control)}{isolation, encryption and compression}
\end{question}

\begin{question}{M}
The first commercially available virtual machine system was
\optsshort{IBM VM/370}{Xen}{Hyper-V}{VMware Workstation}
\end{question}

\begin{question}{M}
Which of the following is a benefit of virtual machines?\\
A. Isolation of guests from each other  B. Running several different operating systems on one machine  C. Easier testing and development of operating systems  D. Server consolidation\\
Choose the correct answer from the options given below:
\optsshort{B, C and D only}{A, B, C and D}{A, C and D only}{A and B only}
\end{question}

\begin{question}{M}
A hypervisor that is built into a general-purpose operating system, such as KVM in Linux or Hyper-V in Windows Server, is classified as
\optsshort{Type 2}{Type 0}{Type 3}{Type 1}
\end{question}

\begin{question}{M}
Overcommitment of memory in a VMM means that
\optslong{the host has more memory than all guests}{each guest gets exactly its share}{memory is never reclaimed}{guests are allocated more memory in total than physically exists}
\end{question}

\begin{question}{M}
Given below are two statements.\\
Statement I: Encrypting data protects confidentiality but not availability.\\
Statement II: A denial-of-service attack threatens availability.
\optslong{Both Statement I and Statement II are true}{Both Statement I and Statement II are false}{Statement I is true but Statement II is false}{Statement I is false but Statement II is true}
\end{question}

\begin{question}{M}
Code injection attacks such as stack-smashing can be mitigated by
\optslong{increasing the disk size}{disabling the cache}{using larger page sizes}{making the stack non-executable (NX bit) and using address-space layout randomisation}
\end{question}

\begin{question}{M}
In the protection model, a domain switch in UNIX occurs when
\optslong{a file is opened}{a page fault occurs}{a process is swapped out}{a process executes a setuid program, temporarily changing its user ID}
\end{question}

\begin{question}{M}
In the access-matrix model, a right marked with an asterisk (for example read*) indicates the
\optslong{copy right, allowing the right to be copied to other domains for the same object}{control right}{owner right}{switch right}
\end{question}

\begin{question}{L}
Which security mechanism examines packets and blocks traffic according to rules, separating trusted and untrusted networks?
\optsshort{Compiler}{Linker}{Firewall}{Assembler}
\end{question}

\begin{answers}
\ans{1}{1}{Entry access(i, j) is the set of operations a process in domain Di may perform on object Oj.}
\ans{2}{1}{An ACL stores a column of the access matrix; a capability list stores a row.}
\ans{3}{2}{Rows are domains and columns are objects: an ACL is a column (per object) and a capability list is a row (per domain).}
\ans{4}{1}{With an ACL, rights are revoked simply by editing the object's list.}
\ans{5}{3}{Aging is a scheduling technique, not a revocation scheme.}
\ans{6}{1}{It limits the damage that a faulty or compromised component can do.}
\ans{7}{2}{For example, passwd runs with root privileges to update the password file.}
\ans{8}{3}{User applications run in ring 3, the least privileged ring.}
\ans{9}{1}{Standard program threats in Silberschatz.}
\ans{10}{2}{Common virus categories.}
\ans{11}{2}{It was one of the first widely spread Internet worms.}
\ans{12}{3}{With different salts, identical passwords produce different hashes.}
\ans{13}{2}{Each password is valid for only one session.}
\ans{14}{3}{For example a password together with a hardware token or a fingerprint.}
\ans{15}{1}{The two main approaches to intrusion detection.}
\ans{16}{1}{The levels from lowest to highest are D, C1, C2, B1, B2, B3 and A1.}
\ans{17}{2}{It compares current file hashes with a stored database.}
\ans{18}{1}{SELinux implements mandatory access control.}
\ans{19}{1}{Solaris 10 and later implement RBAC.}
\ans{20}{3}{Port scanning, worms and denial-of-service attacks target systems and networks.}
\ans{21}{3}{Silberschatz's classification of virtual machine managers.}
\ans{22}{3}{A paravirtualized guest knows about the hypervisor and calls it directly.}
\ans{23}{2}{Non-privileged instructions run natively at full speed.}
\ans{24}{2}{Both are true, but the reason classic x86 failed is that some sensitive instructions (such as POPF) do not trap.}
\ans{25}{3}{Translated code is cached for reuse.}
\ans{26}{4}{They avoid the overhead of software-maintained shadow page tables.}
\ans{27}{4}{Memory pages are copied while the guest runs, then the final state is transferred.}
\ans{28}{2}{This lets the VMM reclaim memory from guests when it is overcommitted.}
\ans{29}{1}{Variations of virtualization beyond full hypervisors.}
\ans{30}{3}{Containers must use the host's kernel; a VM can run any OS.}
\ans{31}{3}{A program in the VM should behave as on real hardware, run mostly at native speed, and the VMM must control all resources.}
\ans{32}{1}{IBM VM/370 (1972) let each user run a separate CMS operating system in a VM.}
\ans{33}{2}{All four are standard benefits.}
\ans{34}{4}{Silberschatz treats them as Type 1 hypervisors because they run in kernel mode and use hardware assistance.}
\ans{35}{4}{It is supported by ballooning, page sharing and swapping.}
\ans{36}{1}{Security goals include confidentiality, integrity and availability.}
\ans{37}{4}{These defences make injected code harder to execute or locate.}
\ans{38}{4}{In UNIX a domain is associated with the user ID.}
\ans{39}{1}{Owner and control rights are separate special rights.}
\ans{40}{3}{A firewall limits the network access between security domains.}
\end{answers}

\section{Linux, Windows and Distributed Systems}

\begin{question}{H}
The Linux kernel is best described as
\optslong{a microkernel}{an exokernel}{a layered kernel with no modules}{a monolithic kernel with loadable kernel modules}
\end{question}

\begin{question}{H}
Match List I with List II.\\
List I (Linux command): A. lsmod  B. insmod  C. rmmod  D. modprobe\\
List II (Function): I. Removes a module from the kernel  II. Loads a module together with its dependencies  III. Lists the currently loaded modules  IV. Inserts a single module into the kernel
\optsshort{A-III, B-IV, C-I, D-II}{A-IV, B-III, C-I, D-II}{A-III, B-IV, C-II, D-I}{A-III, B-I, C-IV, D-II}
\end{question}

\begin{question}{H}
The default scheduler for normal (non-real-time) tasks in Linux since kernel 2.6.23 is the
\optslong{Completely Fair Scheduler (CFS)}{multilevel feedback queue scheduler}{O(1) scheduler}{round-robin scheduler}
\end{question}

\begin{question}{H}
The Completely Fair Scheduler keeps runnable tasks in a
\optslong{FIFO queue}{hash table indexed by priority}{circular list}{red-black tree ordered by virtual runtime (vruntime)}
\end{question}

\begin{question}{H}
In Linux, nice values range from
\optslong{$-$20 to +19, a lower value meaning a higher priority}{0 to 99}{1 to 140, a higher value meaning a higher priority}{$-$19 to +20, a higher value meaning a higher priority}
\end{question}

\begin{question}{M}
Given below are two statements.\\
Statement I: Linux does not distinguish between processes and threads; both are called tasks.\\
Statement II: The clone() system call creates a task, and its flags decide which resources (memory, files, signal handlers) are shared.
\optslong{Both Statement I and Statement II are true}{Both Statement I and Statement II are false}{Statement I is true but Statement II is false}{Statement I is false but Statement II is true}
\end{question}

\begin{question}{H}
Linux manages physical memory using
\optslong{fixed partitions}{an inverted page table only}{only segmentation}{the buddy system for page allocation and the slab allocator for kernel objects}
\end{question}

\begin{question}{M}
In Linux, the out-of-memory (OOM) killer
\optslong{terminates a process to free memory when the system is critically low}{kills only zombie processes}{compresses memory}{adds swap space}
\end{question}

\begin{question}{M}
Match List I with List II.\\
List I (Linux directory): A. /etc  B. /dev  C. /proc  D. /var\\
List II (Contents): I. Device files  II. Variable data such as logs and spool files  III. System configuration files  IV. Virtual file system exposing kernel and process information
\optsshort{A-III, B-I, C-IV, D-II}{A-III, B-IV, C-I, D-II}{A-III, B-I, C-II, D-IV}{A-I, B-III, C-IV, D-II}
\end{question}

\begin{question}{M}
The /proc file system in Linux
\optslong{holds compiled programs}{stores user files on disk}{is a pseudo file system whose contents are generated by the kernel when read}{stores backups}
\end{question}

\begin{question}{M}
Given below are two statements.\\
Statement I: ext2 is a journaling file system.\\
Statement II: ext4 uses extents to describe contiguous ranges of blocks.
\optslong{Both Statement I and Statement II are true}{Both Statement I and Statement II are false}{Statement I is true but Statement II is false}{Statement I is false but Statement II is true}
\end{question}

\begin{question}{M}
Which of the following is NOT an interprocess communication mechanism in Linux?
\optsshort{Page tables}{Signals}{Shared memory}{Pipes}
\end{question}

\begin{question}{M}
In the output of the Linux ps command, a process in state Z is a
\optsshort{running process}{zombie process}{sleeping process}{stopped process}
\end{question}

\begin{question}{M}
Pluggable Authentication Modules (PAM) in Linux
\optslong{encrypt the file system}{manage kernel modules}{provide a common framework that lets system programs use different authentication methods}{schedule tasks}
\end{question}

\begin{question}{M}
Linux and the GNU utilities are distributed mainly under the
\optslong{MIT licence}{Apache licence}{GNU General Public License (GPL)}{BSD licence}
\end{question}

\begin{question}{M}
Which Linux system call replaces the current process image with a new program?
\optsshort{execve()}{fork()}{clone()}{wait()}
\end{question}

\begin{question}{M}
Linux divides device drivers into three classes. Which of the following is NOT one of them?
\optsshort{Block devices}{Character devices}{Paging devices}{Network devices}
\end{question}

\begin{question}{H}
Match List I with List II.\\
List I (Windows component): A. Hardware abstraction layer (HAL)  B. Kernel (dispatcher)  C. Executive  D. Environmental subsystem (e.g. Win32)\\
List II (Role): I. Provides services such as the object manager, virtual-memory manager, I/O manager and security reference monitor  II. Hides chipset differences from the rest of the OS  III. Provides an operating-system personality and API to applications in user mode  IV. Handles thread scheduling, low-level synchronization and interrupts
\optsshort{A-II, B-IV, C-III, D-I}{A-II, B-IV, C-I, D-III}{A-II, B-I, C-IV, D-III}{A-IV, B-II, C-I, D-III}
\end{question}

\begin{question}{H}
Windows uses a priority-based preemptive scheduler with
\optslong{8 priority levels}{140 priority levels}{64 priority levels}{32 priority levels, where 16--31 form the real-time class and 1--15 the variable class}
\end{question}

\begin{question}{M}
In Windows, the unit of CPU scheduling is the
\optsshort{job}{thread}{process}{fiber}
\end{question}

\begin{question}{H}
Given below are two statements.\\
Statement I: NTFS stores all file-system metadata, including the master file table (MFT), in files.\\
Statement II: NTFS uses a log file to support recovery after a crash.
\optslong{Both Statement I and Statement II are true}{Both Statement I and Statement II are false}{Statement I is true but Statement II is false}{Statement I is false but Statement II is true}
\end{question}

\begin{question}{M}
In NTFS, the structure that contains at least one entry (record) for every file on the volume is the
\optsshort{master file table (MFT)}{inode table}{file allocation table}{superblock}
\end{question}

\begin{question}{M}
NTFS directories are indexed using a
\optsshort{bit vector}{linear list}{B+ tree}{hash table}
\end{question}

\begin{question}{M}
Which of the following is a limitation of FAT32 that NTFS overcomes?
\optslong{FAT32 supports only 8.3 names}{FAT32 cannot store files larger than 4 GB and has no file-level access control lists}{FAT32 cannot be read by Windows}{FAT32 cannot be used on USB drives}
\end{question}

\begin{question}{M}
In Windows, configuration information for the system and applications is stored in a hierarchical database called the
\optsshort{active directory}{event log}{master file table}{registry}
\end{question}

\begin{question}{M}
Given below are two statements.\\
Statement I: Terminal Services allow several remote users to have separate sessions on one Windows machine through the Remote Desktop Protocol.\\
Statement II: Fast User Switching lets several users be logged on to the console of the same machine at once, switching between them without logging off.
\optslong{Both Statement I and Statement II are true}{Both Statement I and Statement II are false}{Statement I is true but Statement II is false}{Statement I is false but Statement II is true}
\end{question}

\begin{question}{M}
Windows file and printer sharing over a network is based mainly on the
\optslong{Network File System (NFS) protocol}{File Transfer Protocol}{Simple Mail Transfer Protocol}{Server Message Block (SMB) protocol}
\end{question}

\begin{question}{M}
Active Directory in Windows networks is a directory service based on
\optsshort{SMTP and POP3}{DHCP and ARP}{NFS and NIS}{LDAP and Kerberos}
\end{question}

\begin{question}{M}
Windows uses advanced local procedure call (ALPC) for
\optslong{disk scheduling}{memory allocation}{communication between processes on the same machine}{file compression}
\end{question}

\begin{question}{M}
In Windows, each user and group is identified by a
\optsshort{process ID}{security identifier (SID)}{MAC address}{user ID number}
\end{question}

\begin{question}{M}
Match List I with List II.\\
List I (Windows kernel mechanism): A. Deferred procedure call (DPC)  B. Asynchronous procedure call (APC)  C. Dispatcher object\\
List II (Purpose): I. Postpones interrupt processing until the IRQL drops, after urgent work is done  II. Runs a routine in the context of a particular thread  III. Synchronization object such as an event, mutant, semaphore or timer
\optsshort{A-II, B-I, C-III}{A-I, B-II, C-III}{A-III, B-II, C-I}{A-I, B-III, C-II}
\end{question}

\begin{question}{H}
Given below are two statements.\\
Statement I: In a network operating system, access to remote resources is completely transparent to users.\\
Statement II: In a distributed operating system, users must log in explicitly to each remote machine to use its resources.
\optslong{Both Statement I and Statement II are true}{Both Statement I and Statement II are false}{Statement I is true but Statement II is false}{Statement I is false but Statement II is true}
\end{question}

\begin{question}{H}
Match List I with List II.\\
List I (Migration in a distributed OS): A. Data migration  B. Computation migration  C. Process migration\\
List II (Description): I. A process, or part of it, is executed at a different site  II. Data is transferred to the site where it is needed  III. The computation is carried to the site of the data and only the result is returned
\optsshort{A-III, B-II, C-I}{A-I, B-III, C-II}{A-II, B-III, C-I}{A-II, B-I, C-III}
\end{question}

\begin{question}{H}
Given below are two statements.\\
Statement I: Location transparency means that the name of a file does not reveal its physical storage location.\\
Statement II: Location independence means that the name of a file need not change when its physical storage location changes.
\optslong{Both Statement I and Statement II are true}{Both Statement I and Statement II are false}{Statement I is true but Statement II is false}{Statement I is false but Statement II is true}
\end{question}

\begin{question}{H}
Given below are two statements: one is labelled as Assertion A and the other is labelled as Reason R.\\
Assertion A: A stateful file server recovers from a crash more easily than a stateless server.\\
Reason R: Each request to a stateless server is self-contained, so no client state is lost when it crashes.
\optslong{Both A and R are true and R is the correct explanation of A}{Both A and R are true but R is NOT the correct explanation of A}{A is true but R is false}{A is false but R is true}
\end{question}

\begin{question}{M}
Match List I with List II.\\
List I (Cache-update policy in a DFS): A. Write-through  B. Delayed write  C. Write-on-close\\
List II (Description): I. Changes are sent to the server when the file is closed, as in AFS  II. Changes are written to the server immediately when they are made in the cache  III. Changes are written to the cache and sent to the server later
\optsshort{A-III, B-II, C-I}{A-I, B-III, C-II}{A-II, B-III, C-I}{A-II, B-I, C-III}
\end{question}

\begin{question}{M}
The robustness of a distributed system includes
\optslong{failure detection, reconfiguration and recovery from failure}{only data encryption}{a single central server}{faster CPUs}
\end{question}

\begin{question}{M}
In the Google File System and Hadoop HDFS, file metadata is maintained by
\optslong{every client}{a single master server (NameNode), while data is stored in chunks on chunk servers (DataNodes)}{the DNS}{a peer-to-peer ring}
\end{question}

\begin{question}{M}
Match List I with List II.\\
List I (Naming scheme in a DFS): A. host:local-name  B. Remote directory mounting  C. Total integration\\
List II (Example): I. NFS  II. A single global name structure, as in AFS  III. Name is location-dependent and not transparent
\optsshort{A-I, B-III, C-II}{A-III, B-II, C-I}{A-II, B-I, C-III}{A-III, B-I, C-II}
\end{question}

\begin{question}{M}
Which of the following is an advantage of distributed systems?\\
A. Resource sharing  B. Computation speedup  C. Reliability  D. Communication\\
Choose the correct answer from the options given below:
\optsshort{A and B only}{A, B, C and D}{B and C only}{A, C and D only}
\end{question}

\begin{question}{M}
Lamport's logical clocks are used to
\optslong{detect deadlocks}{synchronise physical clocks exactly}{order events in a distributed system consistently with the happened-before relation}{encrypt messages}
\end{question}

\begin{question}{M}
In a fully distributed mutual exclusion algorithm (Ricart-Agrawala) with n processes, the number of messages needed per critical-section entry is
\optsshort{n $-$ 1}{3}{2(n $-$ 1)}{2}
\end{question}

\begin{question}{M}
In the centralised approach to distributed mutual exclusion, entering and leaving a critical section needs
\optsshort{n messages}{3 messages}{2 messages}{1 message}
\end{question}

\begin{question}{M}
The two-phase commit protocol is used in distributed systems to ensure
\optslong{load balancing}{clock synchronization}{mutual exclusion}{atomicity of a transaction across several sites}
\end{question}

\begin{question}{M}
Which election algorithm chooses the process with the highest priority number as coordinator by sending messages to all processes with higher numbers?
\optsshort{Bakery algorithm}{Ring algorithm}{Banker's algorithm}{Bully algorithm}
\end{question}

\begin{question}{M}
Given below are two statements.\\
Statement I: Circuit switching establishes a permanent physical link for the duration of the communication.\\
Statement II: In packet switching, a dedicated path is reserved for the entire duration of the communication.
\optslong{Both Statement I and Statement II are true}{Both Statement I and Statement II are false}{Statement I is true but Statement II is false}{Statement I is false but Statement II is true}
\end{question}

\begin{question}{M}
Which of the following is a heterogeneity solution used in RPC so that machines with different data representations can communicate?
\optslong{Paging}{Bit stuffing}{Spooling}{External data representation (XDR)}
\end{question}

\begin{question}{L}
Linux was first released by Linus Torvalds in
\optsshort{1995}{1991}{2001}{1985}
\end{question}

\begin{question}{M}
In Linux, the command used to start a process with a given niceness or to change the niceness of a running process is
\optsshort{kill}{nice / renice}{chmod}{top}
\end{question}

\begin{question}{M}
The Linux process with PID 1 in modern distributions is usually
\optsshort{cron}{kthreadd}{systemd}{bash}
\end{question}

\begin{answers}
\ans{1}{4}{All kernel code runs in one address space, but drivers and file systems can be loaded and unloaded at run time.}
\ans{2}{1}{Kernel-module management commands.}
\ans{3}{1}{CFS was replaced by the similar EEVDF scheduler in kernel 6.6, but CFS remains the textbook answer.}
\ans{4}{4}{The task with the smallest vruntime (leftmost node) is chosen next.}
\ans{5}{1}{Real-time priorities occupy 0--99 and normal tasks map to 100--139.}
\ans{6}{1}{A thread is a task created with CLONE\_\allowbreak{}VM, CLONE\_\allowbreak{}FS, CLONE\_\allowbreak{}FILES and CLONE\_\allowbreak{}SIGHAND.}
\ans{7}{4}{Pages are grouped into zones such as ZONE\_\allowbreak{}DMA, ZONE\_\allowbreak{}NORMAL and ZONE\_\allowbreak{}HIGHMEM.}
\ans{8}{1}{It chooses a victim using a badness score.}
\ans{9}{1}{Standard Filesystem Hierarchy Standard directories.}
\ans{10}{3}{For example /proc/cpuinfo and /proc/<pid>/status.}
\ans{11}{4}{Journaling was added in ext3; ext2 has no journal.}
\ans{12}{1}{Linux IPC includes signals, pipes, FIFOs, shared memory, message queues, semaphores and sockets.}
\ans{13}{2}{R is running, S sleeping, T stopped, D uninterruptible sleep.}
\ans{14}{3}{Administrators can change authentication policy without rewriting programs.}
\ans{15}{3}{The GPL requires derived works that are distributed to make their source code available.}
\ans{16}{1}{exec functions such as execvp are library wrappers around execve().}
\ans{17}{3}{Block, character and network devices are the three classes.}
\ans{18}{2}{Layers of the Windows architecture.}
\ans{19}{4}{Priority 0 is reserved for the memory-management (zero-page) thread.}
\ans{20}{2}{Fibers are scheduled by the application in user mode.}
\ans{21}{1}{NTFS is a journaling file system with security descriptors, compression and encryption.}
\ans{22}{1}{Small files can be stored entirely inside their MFT records.}
\ans{23}{3}{B+ trees keep look-up cost logarithmic and avoid reorganisation.}
\ans{24}{2}{NTFS also supports journaling, encryption and much larger volumes.}
\ans{25}{4}{The registry is organised into hives such as HKEY\_\allowbreak{}LOCAL\_\allowbreak{}MACHINE.}
\ans{26}{1}{Both rely on Windows sessions.}
\ans{27}{4}{SMB (and its dialect CIFS) runs over TCP/IP.}
\ans{28}{4}{It stores information about users, computers and policies in a domain.}
\ans{29}{3}{ALPC is used, for example, between clients and environmental subsystems.}
\ans{30}{2}{Access tokens and access control lists are built from SIDs.}
\ans{31}{2}{Windows kernel mechanisms.}
\ans{32}{2}{The descriptions are swapped: network OS users access remote resources explicitly, while a distributed OS makes access transparent.}
\ans{33}{3}{Three forms of migration in distributed systems.}
\ans{34}{1}{Location independence is the stronger property.}
\ans{35}{4}{It is the stateless server (such as NFSv3) that recovers easily, for exactly the reason given in R.}
\ans{36}{3}{Cache-update policies trade reliability against performance.}
\ans{37}{1}{Heartbeat (handshaking) messages are used to detect failures.}
\ans{38}{2}{Files are divided into large chunks (64 MB or 128 MB) that are replicated.}
\ans{39}{4}{Three DFS naming approaches.}
\ans{40}{2}{These are the four main reasons for building distributed systems.}
\ans{41}{3}{If a happened before b, then the timestamp of a is less than that of b.}
\ans{42}{3}{n $-$ 1 request messages and n $-$ 1 reply messages.}
\ans{43}{2}{A request, a reply (grant) and a release message.}
\ans{44}{4}{All sites commit, or all abort.}
\ans{45}{4}{The ring algorithm circulates the election message around a logical ring.}
\ans{46}{3}{In packet switching messages are split into packets that may take different routes; no path is reserved.}
\ans{47}{4}{Data is converted to a machine-independent form before transmission.}
\ans{48}{2}{Version 0.01 was released in 1991.}
\ans{49}{2}{nice starts a process with a given niceness; renice changes it for a running process.}
\ans{50}{3}{systemd replaced the traditional SysV init in most distributions.}
\end{answers}
